\chapter{Manifolds of non-negative curvature}
\label{chap:manifolds-nonneg-curvature}
\label{nonnegcurv}

In studying singularity development in $3$-dimensional Ricci flows
one forms blow-up limits. By this we mean the following. One
considers a sequence of points $x_k$ in the flow converging to the
singularity. It will be the case that $R(x_k)$ tends to $\infty$ as
$k$ tends to $\infty$. We form a sequence of based Riemannian
manifolds labeled by $k$, where the $k^{th}$ Riemannian manifold is
obtained by taking the time-slice of $x_k$, rescaling its metric by
$R(x_k)$, and then taking $x_k$ as the base point. This creates a
sequence with the property that for each member of the sequence the
scalar curvature at the base point is one. Because of a pinching
result of Hamilton's (see Chapter~\ref{secmaxprin}), if there is a
geometric limit of this sequence, or of any subsequence of it, then
that limit is non-negatively curved. Hence, it is important to
understand the basic properties of Riemannian manifolds of
non-negative curvature in order to study singularity development. In
this chapter we review the properties that we shall need. We suppose
that $M$ is non-compact and of positive (resp., non-negative)
curvature. The key to understanding these manifolds is the Busemann
function associated to a minimizing geodesic ray.


\section{Busemann functions}

\begin{definition}[Minimizing geodesic ray]
\label{def:minimizing-geodesic-ray}
\lean{MorganTianLib.IsGeodesicRay}
\leanok
\uses{def:geodesic}
\source{morgan-tian}
A geodesic ray $\lambda\colon [0,\infty)\to M$ is said to be {\em
minimizing} if the restriction of $\lambda$ to every compact
subinterval of $[0,\infty)$ is a length-minimizing geodesic arc,
i.e., a geodesic arc whose length is equal to the distance between
its endpoints. Since a unit-speed arc realizes the distance between
its endpoints exactly when its length equals that distance, a
unit-speed ray $\lambda$ is a minimizing geodesic ray if and only if
$$d(\lambda(s),\lambda(t))=|s-t|\qquad\text{for all }s,t\ge 0;$$
we take this metric characterization as the formal definition (it
makes sense in any metric space).
\end{definition}

\begin{definition}[Minimizing geodesics on an interval]
\label{def:minimizing-geodesic-window}
\lean{MorganTianLib.IsMinGeodesicOn}
\leanok
\uses{def:minimizing-geodesic-ray}
\source{morgan-tian}
Let $(M,d)$ be a metric space and let $I\subset\Ar$. A map
$\gamma\colon I\to M$ is a {\em unit-speed minimizing geodesic on
$I$} if
$$d(\gamma(s),\gamma(t))=|s-t|\qquad\text{for all } s,t\in I.$$
For $I=[a,b]$ such a map is a {\em minimizing geodesic segment}; for
$I=[0,\infty)$ it is exactly a minimizing geodesic ray in the sense
of Definition~\ref{def:minimizing-geodesic-ray}; and for $I=\Ar$ it
is a {\em minimizing geodesic line}. Directly from the definition, a
unit-speed minimizing geodesic is $1$-Lipschitz (hence continuous) on
$I$, its restriction to any subset of $I$ is again a unit-speed
minimizing geodesic, and so is any arclength translate
$t\mapsto\gamma(t+c)$ on the translated set $I-c$.
\end{definition}

\begin{definition}[Spaces with minimizing segments]
\label{def:has-min-segments}
\lean{MorganTianLib.HasMinSegments}
\leanok
\uses{def:minimizing-geodesic-window}
A metric space $(M,d)$ {\em has minimizing segments} if for every
pair of points $x,y\in M$ there is a unit-speed minimizing geodesic
$\sigma\colon[0,d(x,y)]\to M$ with $\sigma(0)=x$ and
$\sigma(d(x,y))=y$. Recall that a metric space is {\em proper} if its
closed metric balls are compact. Every connected, complete Riemannian
manifold has minimizing segments
(Lemma~\ref{lem:minimizing-segment-exists}) and is proper (as
established in the proof of Lemma~\ref{lem:ends-exist} from
Theorem~\ref{thm:hopf-rinow}).
\end{definition}

A limit of minimizing geodesic arcs is again minimizing. The
following lemmas make this precise and record its standard
consequences; they will be used repeatedly throughout this
chapter. We first prove a purely metric limit lemma, which suffices
for all the existence arguments of this section; the Riemannian
refinement (convergence of velocities in $TM$) is
Lemma~\ref{lem:limit-of-minimizing-geodesics} below.

\begin{lemma}[Ultrafilter limits of minimizing geodesics]
\label{lem:ultrafilter-limit-minimizing}
\lean{MorganTianLib.exists_isMinGeodesicOn_hyperfilter_limit}
\leanok
\uses{def:minimizing-geodesic-window}
\source{morgan-tian}
Fix a non-principal ultrafilter $\cU$ on the natural numbers. Let $M$
be a proper metric space, let $I\subset\Ar$ contain $0$, and for
each $k\in\mathbb{N}$ let $\gamma_k$ be a unit-speed minimizing
geodesic on a set $I_k\ni 0$, such that
\begin{enumerate}
\item for every $t\in I$, the set of indices $k$ with $t\in I_k$ is
cofinite; and
\item the anchor points $\gamma_k(0)$ all lie in a single compact
set $K\subset M$.
\end{enumerate}
Then the ultrafilter limits
$$\gamma(t)=\lim_{k\rightarrow\cU}\gamma_k(t)$$
exist for every $t\in I$ and define a unit-speed minimizing geodesic
$\gamma$ on $I$ with $\gamma(0)\in K$.
\end{lemma}

\begin{proof}
\leanok
\uses{def:minimizing-geodesic-window}
Recall two basic properties of a non-principal ultrafilter $\cU$ on
$\mathbb{N}$: every cofinite set belongs to $\cU$ (an ultrafilter
contains one set of every complementary pair, and a non-principal
one contains no finite set), and any sequence in a compact set has a
limit along $\cU$ (the pushforward of $\cU$ is an ultrafilter on the
compact set, and every ultrafilter on a compact space converges).

Fix $t\in I$, and choose $x_0\in K$ and $R<\infty$ with
$K\subset\overline{B}(x_0,R)$ ($K$ is compact, hence bounded). For
the cofinitely many $k$ with $t\in I_k$ we have
$d(\gamma_k(t),\gamma_k(0))=|t-0|=|t|$, so
$$d(\gamma_k(t),x_0)\le
d(\gamma_k(t),\gamma_k(0))+d(\gamma_k(0),x_0)\le|t|+R.$$
Thus $\gamma_k(t)$ lies, for $\cU$-almost every $k$, in the compact
ball $\overline{B}(x_0,R+|t|)$, and the limit
$\gamma(t)=\lim_{k\rightarrow\cU}\gamma_k(t)$ exists.

For $s,t\in I$, the distance function $d\colon M\times M\to\Ar$ is
continuous and limits along ultrafilters commute with continuous
functions, so
$$d(\gamma(s),\gamma(t))
=\lim_{k\rightarrow\cU}d(\gamma_k(s),\gamma_k(t))=|s-t|,$$
since $d(\gamma_k(s),\gamma_k(t))=|s-t|$ for the cofinitely many $k$
with $s,t\in I_k$, and the $\cU$-limit of a sequence that is
$\cU$-almost everywhere equal to a constant is that constant. Hence
$\gamma$ is a unit-speed minimizing geodesic on $I$. Finally,
$\gamma(0)$ is the $\cU$-limit of the points $\gamma_k(0)\in K$, and
$K$ is closed, so $\gamma(0)\in K$.
\end{proof}

\begin{lemma}[Limits of minimizing geodesics: metric version]
\label{lem:metric-limit-of-minimizing}
\lean{MorganTianLib.exists_isMinGeodesicOn_tendstoLocallyUniformlyOn}
\leanok
\uses{def:minimizing-geodesic-window}
\source{morgan-tian}
Let $M$ be a proper metric space, let $I\subset\Ar$ be a closed
interval containing $0$, and for each $k\in\mathbb{N}$ let $\gamma_k$
be a unit-speed minimizing geodesic on an interval $I_k\ni 0$, such
that every point of $I$ belongs to $I_k$ for all sufficiently large
$k$, and such that the anchor points $\gamma_k(0)$ all lie in a
single compact set $K\subset M$. Then there are a unit-speed
minimizing geodesic $\gamma\colon I\to M$ with $\gamma(0)\in K$ and a
subsequence of $(\gamma_k)$ converging to $\gamma$ uniformly on
every compact subset of $I$.
\end{lemma}

\begin{proof}
\leanok
\uses{lem:ultrafilter-limit-minimizing}
Fix a non-principal ultrafilter $\cU$ on $\mathbb{N}$ and let
$\gamma$ be the unit-speed minimizing geodesic on $I$ with
$\gamma(0)\in K$ and $\gamma(t)=\lim_{k\rightarrow\cU}\gamma_k(t)$
($t\in I$) produced by
Lemma~\ref{lem:ultrafilter-limit-minimizing}. It remains to extract
a subsequence converging to $\gamma$ uniformly on compact subsets of
$I$.

{\em A subsequence via finite grids.} For $n\in\mathbb{N}$ let
$$G_n=\left\{\tfrac{j}{n+1}\ :\ j\in\ZZ,\ |j|\le n(n+1)\right\}\cap
I,$$
a finite subset of $I\cap[-n,n]$. Consider, for fixed $n$, two
conditions on an index $k$:
\begin{itemize}
\item[(a)] $I\cap[-n,n]\subset I_k$;
\item[(b)] $d(\gamma_k(q),\gamma(q))<\tfrac1{n+1}$ for every $q\in
G_n$.
\end{itemize}
Condition (a) holds for all sufficiently large $k$: the set
$I\cap[-n,n]$ is compact ($I$ is closed), so if it is nonempty its
infimum $a$ and supremum $b$ belong to it, hence to $I$; for all
large $k$ both $a$ and $b$ lie in $I_k$; and since $I_k$ is an
interval, $I\cap[-n,n]\subset[a,b]\subset I_k$. Condition (b) holds
for a set of indices belonging to $\cU$: it is a finite intersection
of the sets $\{k\,:\,d(\gamma_k(q),\gamma(q))<\tfrac1{n+1}\}$, each
of which belongs to $\cU$ by the $\cU$-convergence
$\gamma_k(q)\rightarrow\gamma(q)$ at the grid point $q\in G_n\subset
I$. Since cofinite sets belong to $\cU$ and every member of $\cU$ is
infinite, for each $n$ the set of indices satisfying both (a) and
(b) is infinite. Choose then indices $k_0<k_1<k_2<\cdots$ such that
$k_n$ satisfies (a) and (b) at stage $n$.

{\em Uniform convergence near each point.} Fix $x\in I$ and
$\epsilon>0$, and let $n\ge\max(|x|+1,\,3/\epsilon)$. Let $y\in I$
with $|y-x|<1$; then $|y|\le|x|+1\le n$. Round $y$ toward $0$ at
scale $\tfrac1{n+1}$: set $q=\lfloor y(n+1)\rfloor/(n+1)$ if
$y\ge0$, and $q=\lceil y(n+1)\rceil/(n+1)$ if $y<0$, so that $q$
lies between $0$ and $y$, $|y-q|\le\tfrac1{n+1}$, and
$q=\tfrac{j}{n+1}$ with $|j|\le|y|(n+1)\le n(n+1)$. Since $I$ is an
interval containing $0$ and $y$, the point $q$ lies in $I$, hence
$q\in G_n$; and $y$ and $q$ lie in $I\cap[-n,n]$, hence in
$I_{k_n}$ by (a). Therefore
$$d(\gamma(y),\gamma_{k_n}(y))\le
d(\gamma(y),\gamma(q))+d(\gamma(q),\gamma_{k_n}(q))
+d(\gamma_{k_n}(q),\gamma_{k_n}(y))
<|y-q|+\tfrac1{n+1}+|q-y|\le\tfrac3{n+1}\le\epsilon,$$
using that $\gamma$ and $\gamma_{k_n}$ are unit-speed minimizing
geodesics on $I$ and on $I_{k_n}$ respectively, together with (b).
Thus $\sup\{d(\gamma(y),\gamma_{k_n}(y)):y\in I,\ |y-x|<1\}\le
\epsilon$ for all $n\ge\max(|x|+1,3/\epsilon)$: the subsequence
$(\gamma_{k_n})$ converges to $\gamma$ uniformly on a neighborhood
of each point of $I$, and hence uniformly on every compact subset
of $I$.
\end{proof}

\begin{lemma}[Rays from noncompact spaces: metric version]
\label{lem:ray-exists-metric}
\lean{MorganTianLib.exists_isGeodesicRay_of_noncompact}
\leanok
\uses{def:minimizing-geodesic-window,def:has-min-segments}
\source{morgan-tian}
Let $M$ be a noncompact proper metric space with minimizing
segments, and let $p\in M$. Then there is a minimizing geodesic ray
$\lambda\colon[0,\infty)\to M$ with $\lambda(0)=p$.
\end{lemma}

\begin{proof}
\leanok
\uses{lem:ultrafilter-limit-minimizing}
First, $M$ is unbounded: if $d(p,\cdot)$ were bounded by some
$R<\infty$, then $M=\overline{B}(p,R)$ would be compact (closed
balls of a proper metric space are compact), contrary to hypothesis.
Choose points $q_k\in M$ with $d_k:=d(p,q_k)\ge k$, and let
$\gamma_k\colon[0,d_k]\to M$ be unit-speed minimizing segments from
$p$ to $q_k$. Take $I=[0,\infty)$ and $I_k=[0,d_k]$: since $d_k\ge
k$, every $t\ge0$ lies in $I_k$ for all $k\ge t$; and
$\gamma_k(0)=p$ for every $k$.
Lemma~\ref{lem:ultrafilter-limit-minimizing}, applied with
$K=\{p\}$, yields a unit-speed minimizing geodesic
$\lambda\colon[0,\infty)\to M$ with $\lambda(0)=p$, i.e., a
minimizing geodesic ray emanating from $p$.
\end{proof}

\begin{lemma}[Distinct ends differ at some compact set]
\label{lem:ends-differ-at-compact}
\lean{MorganTianLib.SpaceOfEnds.exists_apply_ne}
\leanok
\uses{def:space-of-ends}
Let $X$ be a topological space and let $\cE_1\neq\cE_2$ be distinct
points of the space of ends of $X$
(Definition~\ref{def:space-of-ends}). Then there is a compact set
$K\subset X$ such that the connected components of $X\setminus K$
determined by $\cE_1$ and by $\cE_2$ are distinct.
\end{lemma}

\begin{proof}
\leanok
An end is an element of the inverse limit of the system $K\mapsto
\pi_0(X\setminus K)$, i.e., a compatible family of components
indexed by the compact subsets of $X$; two families that agree at
every index are equal. Contrapositively, $\cE_1\neq\cE_2$ forces
disagreement at some compact set $K$.
\end{proof}

\begin{lemma}[Ends reach infinity]
\label{lem:end-component-unbounded}
\lean{MorganTianLib.SpaceOfEnds.exists_mem_connectedComponentIn_forall_le_dist}
\leanok
\uses{def:space-of-ends,def:ends-transition}
\source{morgan-tian}
Let $M$ be a proper metric space, let $\cE$ be an end of $M$, let
$K\subset M$ be compact, and let $U$ be the connected component of
$M\setminus K$ determined by $\cE$. Then for every $R<\infty$ there
is $z\in U$ with $d(z,w)\ge R$ for all $w\in K$.
\end{lemma}

\begin{proof}
\leanok
\uses{def:ends-transition}
Suppose not: every $z\in U$ has some $w\in K$ with $d(z,w)<R$.
Then, $K$ being bounded (it is compact), $U$ is bounded as well, so
its closure $\overline{U}$ is compact, since $M$ is proper. The set
$K'=K\cup\overline{U}$ is then a compact set containing $K$. The
end $\cE$ determines a component $U'$ of $M\setminus K'$; connected
components are nonempty, so we may choose $y\in U'$. By the
compatibility of $\cE$ with the transition map
(Definition~\ref{def:ends-transition}) associated to the inclusion
$K\subset K'$, the component of $M\setminus K$ containing $y$ is
exactly the component determined by $\cE$ at $K$, namely $U$; that
is, $y\in U\subset\overline{U}\subset K'$. But $y\in M\setminus K'$
--- a contradiction.
\end{proof}

\begin{lemma}[Lines from two ends: metric version]
\label{lem:line-exists-metric}
\lean{MorganTianLib.exists_isMinGeodesicOn_univ_of_ends_ne}
\leanok
\uses{def:minimizing-geodesic-window,def:has-min-segments,def:space-of-ends}
\source{morgan-tian}
Let $M$ be a proper metric space with minimizing segments. If the
space of ends of $M$ (Definition~\ref{def:space-of-ends}) has two
distinct points, then $M$ contains a minimizing geodesic line.
\end{lemma}

\begin{proof}
\leanok
\uses{lem:ultrafilter-limit-minimizing,lem:ends-differ-at-compact,lem:end-component-unbounded}
{\em A separating compact set.} Let $\cE_1\neq\cE_2$ be distinct
ends of $M$. By Lemma~\ref{lem:ends-differ-at-compact} there is a
compact set $K\subset M$ such that the components $U_1$ and $U_2$
of $M\setminus K$ determined by $\cE_1$ and $\cE_2$ are distinct.
Distinct connected components are disjoint: two components that
share a point coincide.

{\em Crossing segments.} By
Lemma~\ref{lem:end-component-unbounded}, choose, for each
$k\in\mathbb{N}$, points
$x_k\in U_1$ and $y_k\in U_2$ with $d(x_k,w)\ge k$ and $d(y_k,w)\ge
k$ for all $w\in K$, and let $\sigma_k\colon[0,L_k]\to M$ be a
unit-speed minimizing segment from $x_k$ to $y_k$, where
$L_k=d(x_k,y_k)$. The image of $\sigma_k$ is connected (a
continuous image of an interval); if it avoided $K$ it would lie in
a single connected component of $M\setminus K$, whereas it contains
the point $x_k$ of $U_1$ and the point $y_k$ of the disjoint
component $U_2$. Hence there is $s_k\in[0,L_k]$ with
$\sigma_k(s_k)\in K$.

{\em The limit line.} Recenter: let $\tau_k(t)=\sigma_k(t+s_k)$, a
unit-speed minimizing geodesic on $[-s_k,L_k-s_k]$ with
$\tau_k(0)=\sigma_k(s_k)\in K$. Since $\sigma_k$ is minimizing,
$$s_k=d(x_k,\sigma_k(s_k))\ge k,\qquad
L_k-s_k=d(\sigma_k(s_k),y_k)\ge k,$$
by the choice of $x_k$ and $y_k$ (and $\sigma_k(s_k)\in K$). Thus
for every $t\in\Ar$ we have $t\in[-s_k,L_k-s_k]$ for all $k\ge|t|$,
and the anchor points $\tau_k(0)$ lie in the compact set $K$.
Lemma~\ref{lem:ultrafilter-limit-minimizing}, applied with $I=\Ar$
and anchor set $K$, yields a unit-speed minimizing geodesic
$\lambda\colon\Ar\to M$: a minimizing geodesic line.
\end{proof}

\begin{lemma}[Geodesics depend continuously on their initial conditions]
\label{lem:geodesic-continuous-dependence}
\lean{MorganTianLib.exists_globalGeodesic_initial,
  MorganTianLib.eqOn_of_initial_data,
  MorganTianLib.convAt_of_convAt_zero,
  MorganTianLib.tendsto_apply_of_convAt_zero,
  MorganTianLib.tendstoUniformlyOn_of_convAt_zero}
\leanok
\uses{def:geodesic,thm:hopf-rinow}
\source{morgan-tian}
Let $(M,g)$ be a complete Riemannian manifold.
\begin{enumerate}
\item For each $q\in M$ and each $v\in T_qM$ there is a unique geodesic
$\gamma_{q,v}\colon \Ar\to M$ with $\gamma_{q,v}(0)=q$ and
$\gamma_{q,v}'(0)=v$. Moreover, any geodesic $\beta\colon I\to M$
defined on an interval $I$ containing $0$ and satisfying $\beta(0)=q$
and $\beta'(0)=v$ is the restriction of $\gamma_{q,v}$ to $I$. (If $0$
is an endpoint of $I$, then $\beta'(0)$ denotes the one-sided
derivative.)
\item Suppose that $(q_k,v_k)$ is a sequence with $v_k\in T_{q_k}M$
converging to $(q,v)$ in the tangent bundle $TM$; that is, $q_k\to q$
in $M$ and, in some coordinate chart whose domain contains $q$ (and
hence contains $q_k$ for all large $k$), the coordinates of $q_k$ and
the components of $v_k$ converge to the coordinates of $q$ and the
components of $v$. (The choice of chart is immaterial, since
coordinate transitions and their differentials are continuous.) Then
for every $T<\infty$
$$\lim_{k\rightarrow\infty}\ \sup_{t\in[-T,T]}
d\left(\gamma_{q_k,v_k}(t),\gamma_{q,v}(t)\right)=0,$$
and for every $t\in\Ar$ the pairs
$\left(\gamma_{q_k,v_k}(t),\gamma_{q_k,v_k}'(t)\right)$ converge to
$\left(\gamma_{q,v}(t),\gamma_{q,v}'(t)\right)$ in $TM$.
\end{enumerate}
\end{lemma}

% Formalization status: statement FULLY FORMALIZED (axiom-clean),
% `MorganTianLib/Ch02/GeodesicContinuousDependence.lean`.  Part (1): existence is
% `exists_globalGeodesic_initial` (Hopf--Rinow c⟹d, `exists_global_geodesic`);
% uniqueness/restriction is `eqOn_of_initial_data` (intrinsic uniqueness on a
% preconnected time set).  Part (2): "convergence in $TM$" is rendered exactly as
% the statement itself defines it — convergence of chart states (position and
% chart-at-the-limit coordinate velocity), i.e. the invariant
% `Riemannian.Geodesic.ConvAt g γ γs`.  `convAt_of_convAt_zero` propagates that
% invariant from time 0 to every time (giving the pairwise $TM$-convergence),
% `tendsto_apply_of_convAt_zero` is pointwise position convergence, and
% `tendstoUniformlyOn_of_convAt_zero` is the $\sup_{[-T,T]}\to 0$ clause
% (`TendstoUniformlyOn`).  The Lean route reuses the flow-box convergence step
% `MorganTianLib.convStepProperty` of `lem:ends-exist` (its clopen propagation is
% base-point agnostic, so the initial *points* $q_k$ may also vary) together with
% a uniform Lipschitz bound from the conserved speeds; it is an alternative to the
% self-contained Grönwall argument written below, which proves the same statement.
% As with `lem:ends-exist`, the Ch.\ 1 dependencies `def:geodesic`/`thm:hopf-rinow`
% are formalized in `OpenGALib` though their Ch.\ 1 blueprint nodes are not yet
% `\leanok`.

\begin{proof}
\uses{def:geodesic,thm:hopf-rinow}
The only fact we import from the classical theory of ordinary
differential equations is the local existence of a geodesic with
prescribed initial position and velocity, as invoked in the
discussion following Definition~\ref{def:geodesic}; uniqueness and
continuous dependence are established below by a Gr\"onwall argument.

{\em Chart states.} Fix a coordinate chart $\varphi\colon W\to
\Omega_0$, where $W\subset M$ is open and $\Omega_0\subset\Ar^n$ is
open. For a geodesic $\gamma$ and a time $t$ with $\gamma(t)\in W$,
write $x(t)=\varphi(\gamma(t))$ for the coordinate representation and
call
$$Z(t)=(x(t),\dot x(t))\in \Omega:=\Omega_0\times\Ar^n$$
the {\em chart state} of $\gamma$ at time $t$. In the chart the
geodesic equation is the second-order system $\ddot
x^k=-\Gamma^k_{ij}(x)\,\dot x^i\dot x^j$ displayed after
Definition~\ref{def:geodesic}; equivalently, setting
$$F(x,y)=\left(y,\ -\Gamma^k_{ij}(x)\,y^iy^j\right),$$
a smooth map $F\colon\Omega\to\Ar^{2n}$, a smooth curve in $W$ is a
geodesic if and only if its chart state satisfies the first-order
system $\dot Z(t)=F(Z(t))$. Likewise every point $(q',v')$ of $TM$
with $q'\in W$ has a chart state in $\Omega$, and convergence in $TM$
of a sequence whose limit lies over $W$ is, by the description in
part (2) of the statement, exactly convergence of chart states.

{\em Step 1: local estimates.} Let $z^*=(x^*,y^*)\in\Omega$. Choose
$r>0$ so small that $\overline{B}(x^*,2r)\subset\Omega_0$; then the
Euclidean closed ball $\overline{B}(z^*,2r)\subset\Ar^{2n}$ is a
compact convex subset of $\Omega$. Set
$$C=1+\max_{\overline{B}(z^*,2r)}|F|,\qquad
\epsilon=\frac{r}{2C},$$
and let $L=\max_{\overline{B}(z^*,2r)}\|DF\|$, so that by the mean
value inequality on the convex set $\overline{B}(z^*,2r)$ the map $F$
is $L$-Lipschitz there. We prove three estimates.

(i) ({\em No escape.}) If $\gamma$ is a geodesic defined at least on
$[t_0-\epsilon,t_0+\epsilon]$ with $\gamma(t_0)\in W$ and chart state
$Z(t_0)\in\overline{B}(z^*,r)$, then $\gamma(t)\in W$ for all
$|t-t_0|\le\epsilon$, and $Z(t)\in\overline{B}(z^*,3r/2)$ there.

Indeed, let $A$ be the set of $\tau\in[t_0,t_0+\epsilon]$ such that
$\gamma([t_0,\tau])\subset W$ and $Z(s)\in\overline{B}(z^*,2r)$ for
all $s\in[t_0,\tau]$. Then $t_0\in A$ and $A$ is a subinterval. For
$\tau\in A$ and $s\in[t_0,\tau]$, integrating $\dot Z=F(Z)$ gives
$$|Z(s)-Z(t_0)|\le\int_{t_0}^{s}|F(Z(u))|\,du\le C\,(s-t_0)\le
C\epsilon=\frac{r}{2},$$
so that in fact $Z(s)\in\overline{B}(z^*,3r/2)$ for all
$s\in[t_0,\tau]$. The set $A$ is closed in $[t_0,t_0+\epsilon]$: if
$\tau_j\in A$ increase to $\tau$, then for $s\in[t_0,\tau)$ the
positions $x(s)$ lie in the compact set
$\overline{B}(x^*,3r/2)\subset\Omega_0$, hence the points $\gamma(s)$
lie in the compact set
$\varphi^{-1}(\overline{B}(x^*,3r/2))\subset W$; by continuity of
$\gamma$ the point $\gamma(\tau)$ lies in this compact set as well,
so $\gamma(\tau)\in W$, the chart state $Z(\tau)$ is defined, and by
continuity $Z(\tau)=\lim_{s\rightarrow\tau^-}Z(s)\in
\overline{B}(z^*,3r/2)$; thus $\tau\in A$. Consequently
$\sigma=\sup A$ belongs to $A$; and if $\sigma<t_0+\epsilon$ then,
since $Z(\sigma)$ lies in $\overline{B}(z^*,3r/2)$ --- in the
interior of $\overline{B}(z^*,2r)$ --- and $W$ is open, continuity of
$\gamma$ and of $Z$ near $\sigma$ would put points $\tau'>\sigma$
into $A$, a contradiction. Hence $A=[t_0,t_0+\epsilon]$. The same
argument run backward in time covers $[t_0-\epsilon,t_0]$, proving
(i).

(ii) ({\em Gr\"onwall estimate.}) Let $J$ be an interval containing
$t_0$ and let $\gamma_1,\gamma_2$ be geodesics defined on $J$ whose
points lie in $W$ and whose chart states $Z_1,Z_2$ remain in
$\overline{B}(z^*,2r)$ throughout $J$. Then
$$|Z_1(t)-Z_2(t)|\le e^{L|t-t_0|}\,|Z_1(t_0)-Z_2(t_0)|
\qquad\text{for all } t\in J.$$
Indeed, for $t\in J$ with $t\ge t_0$ set $\phi(t)=|Z_1(t)-Z_2(t)|$.
From $Z_i(t)=Z_i(t_0)+\int_{t_0}^tF(Z_i(s))\,ds$ and the
$L$-Lipschitz property of $F$ on $\overline{B}(z^*,2r)$,
$$\phi(t)\le\phi(t_0)+L\int_{t_0}^t\phi(s)\,ds=:\psi(t).$$
Then $\psi'=L\phi\le L\psi$, so
$\left(e^{-L(t-t_0)}\psi(t)\right)'\le0$, whence
$\phi(t)\le\psi(t)\le e^{L(t-t_0)}\phi(t_0)$. The same computation
backward in time gives the estimate for $t\le t_0$.

(iii) ({\em Chart distance controls Riemannian distance.}) Let
$\Lambda<\infty$ satisfy $g_{ij}(x)u^iu^j\le\Lambda^2|u|^2$ for all
$x\in\overline{B}(x^*,2r)$ and all $u\in\Ar^n$; such $\Lambda$ exists
because the $g_{ij}$ are continuous and $\overline{B}(x^*,2r)$ is
compact. For $x_1,x_2\in\overline{B}(x^*,2r)$ the straight coordinate
segment from $x_1$ to $x_2$ lies in the convex set
$\overline{B}(x^*,2r)$ and has Riemannian length at most
$\Lambda|x_1-x_2|$; hence
$d\left(\varphi^{-1}(x_1),\varphi^{-1}(x_2)\right)\le
\Lambda|x_1-x_2|$.

{\em Step 2: proof of (1).} By local existence there is a geodesic
through $q$ with velocity $v$ defined on an open interval around $0$,
and by Theorem~\ref{thm:hopf-rinow} it extends to a geodesic
$\gamma_{q,v}\colon\Ar\to M$; thus $\gamma_{q,v}(0)=q$ and
$\gamma_{q,v}'(0)=v$. Now let $\beta\colon I\to M$ be any geodesic
with $0\in I$, $\beta(0)=q$ and $\beta'(0)=v$, and let $A\subset I$
be the set of $t\in I$ with $\beta(t)=\gamma_{q,v}(t)$ and
$\beta'(t)=\gamma_{q,v}'(t)$. Then $0\in A$, and $A$ is closed in $I$
because it is the coincidence set of the two continuous maps
$t\mapsto(\beta(t),\beta'(t))$ and
$t\mapsto(\gamma_{q,v}(t),\gamma_{q,v}'(t))$ into the Hausdorff space
$TM$. The set $A$ is also open in $I$: given $t_0\in A$, fix a chart
$W$ around the common point $\beta(t_0)$ and let $z^*$ be the common
chart state at $t_0$, with $r$ and $L$ as in Step 1; by continuity
there is $\eta>0$ such that for $t\in J:=I\cap[t_0-\eta,t_0+\eta]$
both curves lie in $W$ with chart states in $\overline{B}(z^*,2r)$,
and then estimate (ii), whose right-hand side vanishes, forces the
chart states --- hence the curves and their velocities --- to agree
on $J$. Since $I$ is connected, $A=I$. Taking $I=\Ar$ this also shows
that $\gamma_{q,v}$ is unique, proving (1).

{\em Step 3: proof of (2).} Write $\gamma_k=\gamma_{q_k,v_k}$ and
$\gamma=\gamma_{q,v}$. The geodesic equation is invariant under
$t\mapsto-t$, because it is of second order with right-hand side
quadratic in the velocity; hence $t\mapsto\gamma(-t)$ is the geodesic
with initial condition $(q,-v)$, by the uniqueness in part (1), and
similarly for each $\gamma_k$. Since $(q_k,-v_k)\rightarrow(q,-v)$ in
$TM$, it therefore suffices to prove, for each fixed $T<\infty$, that
$$\sup_{t\in[0,T]}d(\gamma_k(t),\gamma(t))\rightarrow0
\qquad\text{and}\qquad
(\gamma_k(t),\gamma_k'(t))\rightarrow(\gamma(t),\gamma'(t))
\ \text{in } TM \text{ for each } t\in[0,T].$$

Let $S$ be the set of $\tau\in[0,T]$ such that
$\sup_{t\in[0,\tau]}d(\gamma_k(t),\gamma(t))\rightarrow0$ and
$(\gamma_k(t),\gamma_k'(t))\rightarrow(\gamma(t),\gamma'(t))$ in $TM$
for every $t\in[0,\tau]$. The hypothesis of (2) says exactly that
$0\in S$, and $S$ is an interval containing $0$, since the defining
conditions for $\tau$ contain those for every $\tau'\le\tau$. Let
$\tau^*=\sup S$.

Fix a chart $\varphi\colon W\to\Omega_0$ with $\gamma(\tau^*)\in W$,
let $z^*$ be the chart state of $\gamma$ at $\tau^*$, and let
$r,C,L,\epsilon,\Lambda$ be as in Step 1 for this chart and this
$z^*$. Since the chart state of $\gamma$ is a continuous function of
$t$ near $\tau^*$, there is $\delta\in(0,\epsilon/2]$ such that for
all $t\in[0,T]$ with $|t-\tau^*|\le\delta$ the point $\gamma(t)$ lies
in $W$ with chart state in $\overline{B}(z^*,r/2)$. Choose
$\tau_0\in S$ with $\tau^*-\tau_0\le\delta$: if $\tau^*\in S$ take
$\tau_0=\tau^*$, and otherwise use that $S$ contains points
arbitrarily close below its supremum.

Since $\tau_0\in S$, the states $(\gamma_k(\tau_0),\gamma_k'(\tau_0))$
converge in $TM$ to $(\gamma(\tau_0),\gamma'(\tau_0))$, whose chart
state $Z(\tau_0)$ lies in $\overline{B}(z^*,r/2)$; by the description
of convergence in $TM$, for all large $k$ the point
$\gamma_k(\tau_0)$ lies in $W$ and the chart states $Z_k(\tau_0)$
converge to $Z(\tau_0)$; in particular
$Z_k(\tau_0)\in\overline{B}(z^*,r)$ for all large $k$. By (i), for
all such $k$ the geodesics $\gamma$ and $\gamma_k$ remain in $W$ with
chart states in $\overline{B}(z^*,3r/2)\subset\overline{B}(z^*,2r)$
for $|t-\tau_0|\le\epsilon$, and then (ii) gives
$$\sup_{|t-\tau_0|\le\epsilon}|Z_k(t)-Z(t)|
\le e^{L\epsilon}\,|Z_k(\tau_0)-Z(\tau_0)|\longrightarrow0.$$
By (iii) it follows that
$\sup\{d(\gamma_k(t),\gamma(t))\,:\,|t-\tau_0|\le\epsilon\}
\le\Lambda\,\sup_{|t-\tau_0|\le\epsilon}|Z_k(t)-Z(t)|
\rightarrow0$, and for each fixed $t$ with $|t-\tau_0|\le\epsilon$
the convergence $Z_k(t)\rightarrow Z(t)$ of chart states is exactly
the convergence $(\gamma_k(t),\gamma_k'(t))\rightarrow
(\gamma(t),\gamma'(t))$ in $TM$. Combining this with $\tau_0\in S$,
we conclude that $\tau_1:=\min(T,\tau_0+\epsilon)\in S$. Since
$\tau_0+\epsilon\ge\tau^*-\epsilon/2+\epsilon>\tau^*$, the case
$\tau^*<T$ would give an element $\tau_1>\tau^*$ of $S$,
contradicting $\tau^*=\sup S$. Hence $\tau^*=T$ and $\tau_1=T\in S$,
which is assertion (2).
\end{proof}

\begin{lemma}[Limits of minimizing geodesics]
\label{lem:limit-of-minimizing-geodesics}
\lean{MorganTianLib.exists_isMinGeodesicOn_convAt_of_tendsto}
\leanok
\uses{def:geodesic}
\source{morgan-tian}
Let $M$ be a complete Riemannian manifold, let $I\subset\Ar$ be an
interval containing $0$ and containing more than one point, and let
$I_k\subset\Ar$ be closed intervals containing $0$ with the property
that every compact subset of $I$ is contained in $I_k$ for all
sufficiently large $k$. For each $k$ let $\gamma_k\colon I_k\to M$ be
a unit-speed minimizing geodesic, i.e., a unit-speed geodesic with
$$d(\gamma_k(s),\gamma_k(t))=|s-t|\qquad\text{for all } s,t\in I_k,$$
and suppose that the initial points $\gamma_k(0)$ converge to a point
$p\in M$. Then, after passing to a subsequence, the $\gamma_k$
converge to a unit-speed minimizing geodesic $\gamma\colon I\to M$
with $\gamma(0)=p$, in the following sense: for every compact subset
$J\subset I$ we have $J\subset I_k$ for all large $k$ and
$$\lim_{k\rightarrow\infty}\ \sup_{t\in J}
d(\gamma_k(t),\gamma(t))=0;$$
moreover the initial velocities converge:
$(\gamma_k(0),\gamma_k'(0))\rightarrow(\gamma(0),\gamma'(0))$ in
$TM$.
\end{lemma}

% Formalization status: FULLY FORMALIZED (axiom-clean),
% `MorganTianLib.exists_isMinGeodesicOn_convAt_of_tendsto` in
% `MorganTianLib/Ch02/LimitMinimizingGeodesics.lean`.  The Lean statement uses the
% ambient formulation with *global* geodesics `γs k : ℝ → M` (`IsGeodesic`,
% continuous, unit conserved speed `speedSq g (γs k) 0 = 1`, and
% `IsMinGeodesicOn (γs k) (In k)`): the blueprint's preliminary extension of the
% window-geodesics `γ_k : I_k → M` to global ones (its appeal to
% Lemma~\ref{lem:geodesic-continuous-dependence}(1)) is folded into these
% hypotheses.  "Convergence in `TM`" of the initial data
% `(γ_k(0), γ_k'(0)) → (γ(0), γ'(0))` is rendered exactly as the invariant
% `Riemannian.Geodesic.ConvAt g γ (γs ∘ φ) 0` (positions and chart-at-the-limit
% velocities converge), and the uniform-on-compacts clause as
% `TendstoUniformlyOn (γs ∘ φ) γ atTop (Icc (-T) T)` for every `T`, which covers
% every compact `J ⊆ I`.  The only new ingredient over
% `lem:geodesic-continuous-dependence` is Step 1: the chart velocities are bounded
% through the conserved unit speed by the local coordinate-norm bound
% `exists_sq_norm_deriv_le_speedSq`, so Bolzano–Weierstrass on a closed ball of the
% finite-dimensional model (`IsCompact.tendsto_subseq`) extracts a convergent
% velocity subsequence.
\begin{proof}
\leanok
\uses{lem:geodesic-continuous-dependence,thm:hopf-rinow,thm:levi-civita-connection,def:geodesic}
Since $I$ contains $0$ and more than one point, it contains a compact
interval $J_0$ of positive length with $0\in J_0$; discarding
finitely many terms of the sequence, we may assume $J_0\subset I_k$
for every $k$, so that each $I_k$ has positive length and each
$\gamma_k$ has a well-defined (possibly one-sided) unit velocity
$v_k:=\gamma_k'(0)$ at $0$. Write $q_k:=\gamma_k(0)$, so $q_k\to p$.
By Lemma~\ref{lem:geodesic-continuous-dependence}(1), $\gamma_k$ is
the restriction to $I_k$ of the complete geodesic
$\bar\gamma_k:=\gamma_{q_k,v_k}\colon\Ar\to M$ (whose global
existence rests on Theorem~\ref{thm:hopf-rinow}).

{\em Step 1: a convergent subsequence of initial conditions.} Fix a
coordinate chart $\varphi\colon W\to\Omega_0$ around $p$ and a closed
coordinate ball $\overline{B}\subset\Omega_0$ centered at
$\varphi(p)$; for all large $k$ we have $q_k\in
\varphi^{-1}(\overline B)$. The function
$(x,u)\mapsto g_{ij}(x)u^iu^j$ is continuous and positive on the
compact set $\overline{B}\times S^{n-1}$, where $S^{n-1}$ is the
Euclidean unit sphere of $\Ar^n$; let $c>0$ be its minimum. By
homogeneity, $g_{ij}(x)u^iu^j\ge c|u|^2$ for all
$x\in\overline{B}$ and $u\in\Ar^n$. Since each $v_k$ is a unit
vector, $g_{ij}(\varphi(q_k))v_k^iv_k^j=1$, so the Euclidean norms of
the component vectors of the $v_k$ are bounded by $c^{-1/2}$. The
coordinates of the $q_k$ converge to those of $p$; hence, by the
Bolzano--Weierstrass theorem, after passing to a subsequence the
components of the $v_k$ converge to some $\bar v\in\Ar^n$, defining a
vector $v\in T_pM$ with $(q_k,v_k)\rightarrow(p,v)$ in $TM$. By
continuity of the $g_{ij}$,
$$\langle v,v\rangle
=\lim_{k\rightarrow\infty}g_{ij}(\varphi(q_k))v_k^iv_k^j=1,$$
so $v$ is a unit vector.

{\em Step 2: the limit geodesic.} Let
$\bar\gamma:=\gamma_{p,v}\colon\Ar\to M$ be the geodesic of
Lemma~\ref{lem:geodesic-continuous-dependence}(1) and set
$\gamma:=\bar\gamma|_I$. Then $\gamma(0)=p$ and $\gamma'(0)=v$, so
along the chosen subsequence
$(\gamma_k(0),\gamma_k'(0))=(q_k,v_k)\rightarrow(p,v)
=(\gamma(0),\gamma'(0))$ in $TM$. Since the Levi-Civita connection is
compatible with the metric (Theorem~\ref{thm:levi-civita-connection}),
$$\frac{d}{dt}\langle\bar\gamma',\bar\gamma'\rangle
=2\langle\nabla_{\bar\gamma'}\bar\gamma',\bar\gamma'\rangle=0,$$
so $|\gamma'(t)|=|v|=1$ for all $t\in I$: the curve $\gamma$ is a
unit-speed geodesic.

{\em Step 3: uniform convergence on compact subsets of $I$.} Let
$J\subset I$ be compact. By hypothesis $J\subset I_k$ for all large
$k$, and $J\subset[-T,T]$ for some $T<\infty$. Applying
Lemma~\ref{lem:geodesic-continuous-dependence}(2) to
$(q_k,v_k)\rightarrow(p,v)$ gives
$\sup_{t\in[-T,T]}d(\bar\gamma_k(t),\bar\gamma(t))\rightarrow0$;
since $\gamma_k=\bar\gamma_k|_{I_k}$ and $J\subset I_k$ for large
$k$, this yields
$\sup_{t\in J}d(\gamma_k(t),\gamma(t))\rightarrow0$.

{\em Step 4: the limit is minimizing.} Let $s,t\in I$ with $s\le t$.
The interval $[s,t]$ is a compact subset of $I$, so $s,t\in I_k$ for
all large $k$ and, by Step 3, $\gamma_k(s)\rightarrow\gamma(s)$ and
$\gamma_k(t)\rightarrow\gamma(t)$. The distance function on $M\times
M$ is continuous, since by the triangle inequality
$|d(a,b)-d(a',b')|\le d(a,a')+d(b,b')$. Hence
$$d(\gamma(s),\gamma(t))
=\lim_{k\rightarrow\infty}d(\gamma_k(s),\gamma_k(t))=t-s.$$
Thus the restriction of the unit-speed geodesic $\gamma$ to any
compact subinterval of $I$ is an arc whose length equals the distance
between its endpoints; that is, $\gamma\colon I\to M$ is a unit-speed
minimizing geodesic, and $\gamma(0)=p$.
\end{proof}

\begin{lemma}[Minimizing segments exist in complete manifolds]
\label{lem:minimizing-segment-exists}
\lean{MorganTianLib.hasMinSegments_of_complete}
\leanok
\uses{def:geodesic}
Let $M$ be a connected, complete Riemannian manifold and let $x,q\in M$ be
distinct points. Then there is a unit-speed minimizing geodesic
$\sigma\colon [0,d(x,q)]\to M$ with $\sigma(0)=x$ and
$\sigma(d(x,q))=q$, i.e., a geodesic arc from $x$ to $q$ whose
length is $d(x,q)$.
\end{lemma}

% Formalization note: the Lean formalization
% `MorganTianLib.hasMinSegments_of_complete` proves the metric-space predicate
% `HasMinSegments M` (any two points joined by a unit-speed minimizing segment,
% covering also the degenerate x = q case by the constant curve). It takes the
% equivalent Hopf-Rinow route rather than the cut-locus argument below: do
% Carmo's proportional-to-arclength minimizing geodesic
% (Riemannian.Geodesic.exists_minimizing_geodesic, do Carmo Ch. 7 Thm 2.8 f,
% dist (gamma s)(gamma t) = |s - t| * d(x,q) on [0,1]) reparametrized to unit
% speed by s |-> gamma(s / d(x,q)) on [0, d(x,q)].
\begin{proof}
\leanok
\uses{def:cut-locus,lem:cut-time-star-shaped,def:exponential-map,thm:hopf-rinow,lem:limit-of-minimizing-geodesics}
Set $d=d(x,q)>0$. Let $U_x\subset T_xM$ be the open set of
Definition~\ref{def:cut-locus}, so that the cut locus ${\mathcal
C}_x=M\setminus \exp_x(U_x)$ is a closed subset of $M$ of measure
zero.

First suppose $q\notin{\mathcal C}_x$. Then $q=\exp_x(w)$ for some
$w\in U_x$, and by Lemma~\ref{lem:cut-time-star-shaped}(2) the
geodesic $\gamma_w$ is a minimal geodesic from $x$ to $q$ with
$|w|=d(x,\exp_x(w))=d$. Its unit-speed reparameterization
$\sigma(s)=\exp_x(sw/|w|)$, $0\le s\le d$, is the required
minimizing segment.

Now suppose that $q\in{\mathcal C}_x$. Since ${\mathcal C}_x$ is
closed of measure zero, its complement is dense: every non-empty open subset
of $M$ has positive measure (in a coordinate chart it contains the
image of a Euclidean ball, and the Riemannian volume has a positive
smooth density with respect to Lebesgue measure in the chart), so
no metric ball is contained in ${\mathcal C}_x$. Choose therefore a
sequence $q_k\rightarrow q$ with $q_k\notin{\mathcal C}_x$; since
$q\not=x$ we may also assume $q_k\not=x$ for all $k$. By the
previous paragraph there are unit-speed minimizing geodesics
$\tau_k\colon [0,d_k]\to M$ from $x$ to $q_k$, where
$d_k=d(x,q_k)$. By the triangle inequality $|d_k-d|\le d(q_k,q)$,
so $d_k\rightarrow d$.

We apply Lemma~\ref{lem:limit-of-minimizing-geodesics} with
$I=[0,d)$ and $I_k=[0,d_k]$: every compact subset of $[0,d)$ is
contained in $[0,d-\epsilon]$ for some $\epsilon>0$, and
$d_k>d-\epsilon$ for all sufficiently large $k$; also
$\tau_k(0)=x$ for every $k$. Hence, after passing to a subsequence,
the $\tau_k$ converge, uniformly on compact subsets of $[0,d)$, to
a unit-speed minimizing geodesic $\sigma_0\colon [0,d)\to M$ with
$\sigma_0(0)=x$. For each fixed $s\in[0,d)$ and each large $k$, the
restriction $\tau_k|_{[s,d_k]}$ is an arc of length $d_k-s$ from
$\tau_k(s)$ to $q_k$, so $d(\tau_k(s),q_k)\le d_k-s$ and hence
$$d(\sigma_0(s),q)\le
d(\sigma_0(s),\tau_k(s))+d(\tau_k(s),q_k)+d(q_k,q)\le
d(\sigma_0(s),\tau_k(s))+(d_k-s)+d(q_k,q).$$
Letting $k\rightarrow\infty$ gives $d(\sigma_0(s),q)\le d-s$.

Since $M$ is complete, by Theorem~\ref{thm:hopf-rinow} the geodesic
$\sigma_0$ extends to a geodesic $\sigma$ defined for all time; in
particular $\sigma$ is defined and continuous on $[0,d]$ and agrees
with $\sigma_0$ on $[0,d)$. As $s\rightarrow d^-$ we have
$d(\sigma(s),q)=d(\sigma_0(s),q)\le d-s\rightarrow 0$, so by
continuity $\sigma(d)=\lim_{s\rightarrow d^-}\sigma(s)=q$. Thus
$\sigma|_{[0,d]}$ is a unit-speed geodesic arc from $x$ to $q$ of
length $d=d(x,q)$, i.e., a minimizing geodesic segment from $x$ to
$q$.
\end{proof}

\begin{lemma}[Non-compact manifolds have ends]
\label{lem:ends-exist}
\label{ends}
\lean{MorganTianLib.exists_isGeodesicRay_of_complete_noncompact', MorganTianLib.exists_isMinGeodesicOn_univ_of_complete_ends_ne', MorganTianLib.properSpace_of_complete', MorganTianLib.geodesicEndpointContinuity, MorganTianLib.convStepProperty}
\leanok
\uses{def:geodesic,def:space-of-ends}
\source{morgan-tian}
Suppose that $M$ is a complete, connected, non-compact Riemannian
manifold and let $p$ be a point of $M$. Then $M$ has a minimizing
geodesic ray emanating from $p$. If $M$ has more than one end, then
it has a minimizing line.
\end{lemma}

% Formalization status: FULLY FORMALIZED (axiom-clean).  The ray/line are
% produced by the unconditional `MorganTianLib.exists_isGeodesicRay_of_complete_noncompact'`
% / `exists_isMinGeodesicOn_univ_of_complete_ends_ne'` (`Ch02/EndsConvStep.lean`)
% from `ProperSpace M` (`properSpace_of_complete'`), `HasMinSegments M`
% (`hasMinSegments_of_complete`), and the metric backbones
% `exists_isGeodesicRay_of_noncompact` / `exists_isMinGeodesicOn_univ_of_ends_ne`.
% Properness follows from Hopf--Rinow c⟹d (`isGeodesicallyComplete_of_complete`)
% plus OpenGALib's `properSpace_of_forall_geodesic`, whose endpoint-continuity
% input `GeodesicEndpointContinuity g p` is now the unconditional
% `MorganTianLib.geodesicEndpointContinuity`: the reduction
% `geodesicEndpointContinuity_of_convStep` (clopen/connectedness on the flow-box
% convergence invariant + uniform-Lipschitz upgrade) applied to the discharged
% flow-box step `MorganTianLib.convStepProperty` (the universally quantified form of
% do Carmo's sorried `exists_conv_step`, re-derived in `EndsConvStep.lean` from
% the readback `exists_uniform_flow_readback`, the flow's Lipschitz dependence on
% initial data, chart-to-chart velocity transfer, and an affine rescaling).

\begin{proof}
\uses{def:geodesic,thm:hopf-rinow,def:exponential-map,lem:minimizing-segment-exists,def:has-min-segments,lem:ray-exists-metric,lem:line-exists-metric,def:space-of-ends}
We verify that the metric space underlying $M$ satisfies the
hypotheses of Lemmas~\ref{lem:ray-exists-metric}
and~\ref{lem:line-exists-metric}: it is noncompact by hypothesis, it
has minimizing segments, and it is proper.

$M$ has minimizing segments in the sense of
Definition~\ref{def:has-min-segments}: any two distinct points
$x,y\in M$ are joined by a unit-speed minimizing geodesic segment by
Lemma~\ref{lem:minimizing-segment-exists}, and for $x=y$ the
constant map works.

$M$ is proper, i.e., closed metric balls in $M$ are compact. Indeed,
fix $x\in M$ and $R<\infty$. If $|w|\le R$, then
$s\mapsto\exp_x(sw)$, $s\in[0,1]$, is a geodesic of constant speed
$|w|$ from $x$ to $\exp_x(w)$, so $d(x,\exp_x(w))\le|w|\le R$.
Conversely, if $d(x,y)\le R$, let $\sigma$ be a unit-speed
minimizing segment from $x$ to $y$; then $s\mapsto\sigma(s\,d(x,y))$
is a geodesic parameterized by $[0,1]$ --- if $\sigma$ is a geodesic
and $c$ a constant, then $\tilde\sigma(s)=\sigma(cs)$ satisfies
$\nabla_{\tilde\sigma'}\tilde\sigma'
=c^2\,\nabla_{\sigma'}\sigma'=0$ --- with initial velocity
$w=d(x,y)\,\sigma'(0)$ of norm $d(x,y)\le R$, and $y=\exp_x(w)$.
Hence
$$\overline{B}(x,R)=\exp_x\left(\{w\in T_xM\ :\ |w|\le R\}\right),$$
the image of a compact set under the map $\exp_x$, which is smooth
and defined on all of $T_xM$ by completeness
(Definition~\ref{def:exponential-map},
Theorem~\ref{thm:hopf-rinow}); so $\overline{B}(x,R)$ is compact.

Given this, Lemma~\ref{lem:ray-exists-metric} yields a minimizing
geodesic ray emanating from $p$. If moreover $M$ has more than one
end --- that is, if the space of ends of
Definition~\ref{def:space-of-ends} has two distinct points --- then
Lemma~\ref{lem:line-exists-metric} yields a minimizing geodesic
line in $M$.
\end{proof}



\begin{definition}[Busemann approximants]
\label{def:busemann-approx}
\lean{MorganTianLib.busemannAux}
\leanok
\uses{def:minimizing-geodesic-ray}
\source{morgan-tian}
Suppose that $\lambda\colon [0,\infty)\to M$ is a minimizing
geodesic ray. For each $t\ge 0$ we consider the function
$$B_{\lambda,t}(x)=d(\lambda(t),x)-t,\qquad x\in M.$$
\end{definition}

\begin{lemma}[Busemann approximants are $1$-Lipschitz]
\label{lem:busemann-approx-lipschitz}
\lean{MorganTianLib.lipschitzWith_busemannAux}
\leanok
\uses{def:busemann-approx}
Let $\lambda\colon [0,\infty)\to M$ be a minimizing geodesic ray.
For each $t\ge 0$ and all $x,y\in M$ we have
$$|B_{\lambda,t}(x)-B_{\lambda,t}(y)|\le d(x,y).$$
\end{lemma}

\begin{proof}
\leanok
By the triangle inequality
$|d(\lambda(t),x)-d(\lambda(t),y)|\le d(x,y)$, and the constant
$-t$ cancels in the difference.
\end{proof}

\begin{lemma}[Busemann approximants are non-increasing in $t$]
\label{lem:busemann-approx-antitone}
\lean{MorganTianLib.busemannAux_antitoneOn}
\leanok
\uses{def:busemann-approx}
Let $\lambda\colon [0,\infty)\to M$ be a minimizing geodesic ray.
For each $x\in M$ the function $t\mapsto B_{\lambda,t}(x)$ is
non-increasing on $[0,\infty)$.
\end{lemma}

\begin{proof}
\leanok
For $0\le s\le t$ the triangle inequality and the fact that
$\lambda$ is minimizing give
$$d(\lambda(t),x)\le d(\lambda(t),\lambda(s))+d(\lambda(s),x)
=(t-s)+d(\lambda(s),x),$$
so $B_{\lambda,t}(x)=d(\lambda(t),x)-t\le
d(\lambda(s),x)-s=B_{\lambda,s}(x)$.
\end{proof}

\begin{lemma}[Busemann approximants are bounded below]
\label{lem:busemann-approx-lower-bound}
\lean{MorganTianLib.neg_dist_le_busemannAux}
\leanok
\uses{def:busemann-approx}
Let $\lambda\colon [0,\infty)\to M$ be a minimizing geodesic ray.
For each $x\in M$ and each $t\ge 0$ we have
$$B_{\lambda,t}(x)\ge -d(x,\lambda(0)).$$
In particular, $B_{\lambda,t}(\lambda(0))=0$ for all $t\ge 0$.
\end{lemma}

\begin{proof}
\leanok
Since $\lambda$ is minimizing,
$t=d(\lambda(0),\lambda(t))\le d(\lambda(0),x)+d(x,\lambda(t))$,
whence $d(\lambda(t),x)-t\ge -d(x,\lambda(0))$. For $x=\lambda(0)$
we have $B_{\lambda,t}(\lambda(0))=d(\lambda(t),\lambda(0))-t=t-t=0$.
\end{proof}

\begin{definition}[Busemann function]
\label{def:busemann-function}
\lean{MorganTianLib.busemann}
\leanok
\uses{def:busemann-approx,lem:busemann-approx-antitone,lem:busemann-approx-lower-bound,lem:ends-exist}
\source{morgan-tian}
Let $\lambda\colon [0,\infty)\to M$ be a minimizing geodesic ray. By
Lemmas~\ref{lem:busemann-approx-antitone}
and~\ref{lem:busemann-approx-lower-bound}, for each $x\in M$ the
function $t\mapsto B_{\lambda,t}(x)$ is non-increasing and bounded
below, so
$$B_\lambda(x)=\lim_{t\rightarrow \infty} B_{\lambda,t}(x)
=\inf_{t\ge 0}B_{\lambda,t}(x)$$
exists. The function $B_\lambda$ is the {\sl Busemann function} for
$\lambda$.
\end{definition}

\begin{lemma}[Busemann approximants converge to the Busemann function]
\label{lem:busemann-limit}
\lean{MorganTianLib.tendsto_busemannAux_atTop}
\leanok
\uses{def:busemann-function}
Let $\lambda\colon [0,\infty)\to M$ be a minimizing geodesic ray.
For each $x\in M$ we have $B_{\lambda,t}(x)\rightarrow B_\lambda(x)$
as $t\rightarrow\infty$.
\end{lemma}

\begin{proof}
\leanok
\uses{lem:busemann-approx-antitone,lem:busemann-approx-lower-bound}
A function of $t$ that is non-increasing and bounded below converges
to its infimum as $t\rightarrow\infty$.
\end{proof}

\begin{lemma}[The Busemann function bounds its approximants from below]
\label{lem:busemann-le-approx}
\lean{MorganTianLib.busemann_le_busemannAux}
\leanok
\uses{def:busemann-function}
Let $\lambda\colon [0,\infty)\to M$ be a minimizing geodesic ray.
For each $y\in M$ and each $t\ge 0$ we have
$$B_\lambda(y)\le B_{\lambda,t}(y).$$
\end{lemma}

\begin{proof}
\leanok
$B_\lambda(y)$ is the infimum of the values $B_{\lambda,t}(y)$ over
$t\ge 0$.
\end{proof}

\begin{lemma}[The Busemann function is $1$-Lipschitz]
\label{lem:busemann-lipschitz}
\lean{MorganTianLib.lipschitzWith_busemann}
\leanok
\uses{def:busemann-function}
Let $\lambda\colon [0,\infty)\to M$ be a minimizing geodesic ray.
For all $x,y\in M$ we have
$$|B_\lambda(x)-B_\lambda(y)|\le d(x,y).$$
In particular $B_\lambda$ is continuous.
\end{lemma}

\begin{proof}
\leanok
\uses{lem:busemann-approx-lipschitz}
By Lemma~\ref{lem:busemann-approx-lipschitz} we have
$B_{\lambda,t}(x)\le B_{\lambda,t}(y)+d(x,y)$ for every $t\ge 0$;
taking infima over $t\ge 0$ gives $B_\lambda(x)\le
B_\lambda(y)+d(x,y)$, and symmetrically in $x$ and $y$.
\end{proof}

\begin{lemma}[Lower bound for the Busemann function]
\label{lem:busemann-lower-bound}
\lean{MorganTianLib.neg_dist_le_busemann}
\leanok
\uses{def:busemann-function}
Let $\lambda\colon [0,\infty)\to M$ be a minimizing geodesic ray.
For all $x\in M$ we have $B_\lambda(x)\ge -d(x,\lambda(0))$.
\end{lemma}

\begin{proof}
\leanok
\uses{lem:busemann-approx-lower-bound}
Each $B_{\lambda,t}(x)$, $t\ge 0$, is bounded below by
$-d(x,\lambda(0))$ by Lemma~\ref{lem:busemann-approx-lower-bound},
hence so is their infimum.
\end{proof}

\begin{lemma}[The Busemann function along its ray]
\label{lem:busemann-along-ray}
\lean{MorganTianLib.busemann_apply_ray}
\leanok
\uses{def:busemann-function}
Let $\lambda\colon [0,\infty)\to M$ be a minimizing geodesic ray.
For all $s\ge 0$ we have $B_\lambda(\lambda(s))=-s$.
\end{lemma}

\begin{proof}
\leanok
For $t\ge 0$ we have
$$B_{\lambda,t}(\lambda(s))=d(\lambda(t),\lambda(s))-t=|s-t|-t,$$
which equals $s-2t\ge -s$ for $0\le t\le s$ and equals $-s$ for
$t\ge s$. Hence the infimum over $t\ge 0$ is $-s$.
\end{proof}

Since $B_\lambda$ is Lipschitz, $\nabla B_\lambda$ is
well-defined as an $L^2$-vector field.

\begin{definition}[Asymptotic rays]
\label{def:asymptotic-ray}
\lean{MorganTianLib.IsAsymptoticRay}
\leanok
\uses{def:minimizing-geodesic-ray,def:minimizing-geodesic-window}
\source{morgan-tian}
Let $M$ be a metric space, let $\lambda\colon[0,\infty)\to M$ be a
minimizing geodesic ray, and let $\mu\colon[0,\infty)\to M$ be a
unit-speed minimizing geodesic ray. We call $\mu$ a {\em ray
asymptotic to $\lambda$} if there are a sequence
$t_n\rightarrow\infty$ and unit-speed minimizing geodesic segments
$\sigma_n\colon[0,d(\mu(0),\lambda(t_n))]\to M$ from $\mu(0)$ to
$\lambda(t_n)$ such that $\sigma_n$ converges to $\mu$ uniformly on
compact subsets of $[0,\infty)$. (The condition makes sense: since
$\lambda$ is minimizing, $d(\lambda(0),\lambda(t))=t$, so by the
triangle inequality
$d(\mu(0),\lambda(t_n))\ge t_n-d(\mu(0),\lambda(0))
\rightarrow\infty$, and every compact subset of $[0,\infty)$ is
contained in the domain of $\sigma_n$ for all sufficiently large
$n$.)
\end{definition}

\begin{lemma}[Existence of asymptotic rays]
\label{lem:asymptotic-ray-exists}
\lean{MorganTianLib.exists_isAsymptoticRay}
\leanok
\uses{def:asymptotic-ray,def:has-min-segments}
\source{morgan-tian}
Let $M$ be a proper metric space with minimizing segments (for
example, a connected, complete Riemannian manifold), let
$\lambda\colon[0,\infty)\to M$ be a minimizing geodesic ray, and
let $x\in M$. Then there is at least one ray asymptotic to
$\lambda$ emanating from $x$.
\end{lemma}

\begin{proof}
\leanok
\uses{lem:metric-limit-of-minimizing}
Set $t_n=n$ and $d_n=d(x,\lambda(t_n))$. Since $\lambda$ is
minimizing, $d(\lambda(0),\lambda(t_n))=t_n$, so the triangle
inequality gives $d_n\ge t_n-d(x,\lambda(0))$; in particular every
$s\ge0$ satisfies $s\le d_n$ for all sufficiently large $n$. Let
$\sigma_n\colon[0,d_n]\to M$ be unit-speed minimizing segments from
$x$ to $\lambda(t_n)$, which exist since $M$ has minimizing
segments. We apply Lemma~\ref{lem:metric-limit-of-minimizing} with
$I=[0,\infty)$, $I_n=[0,d_n]$ and $K=\{x\}$: every point of
$[0,\infty)$ lies in $I_n$ for all large $n$, and $\sigma_n(0)=x$
for every $n$. Hence there are a unit-speed minimizing geodesic ray
$\mu\colon[0,\infty)\to M$ with $\mu(0)=x$ and a subsequence of the
$\sigma_n$ converging to $\mu$ uniformly on compact subsets of
$[0,\infty)$. Passing to this subsequence and relabeling the $t_n$
accordingly exhibits $\mu$ as a ray asymptotic to $\lambda$
emanating from $x$.
\end{proof}

\begin{lemma}[Busemann values along asymptotic rays]
\label{lem:asymptotic-ray}
\lean{MorganTianLib.IsAsymptoticRay.busemann_apply}
\leanok
\uses{def:busemann-function,def:asymptotic-ray}
Let $M$ be a metric space, let $\lambda\colon[0,\infty)\to M$ be a
minimizing geodesic ray, and let $\mu$ be a ray asymptotic to
$\lambda$, with $x=\mu(0)$. Then
$$B_\lambda(\mu(s))=B_\lambda(x)-s\ \ \text{for all } s\ge 0.$$
\end{lemma}

\begin{proof}
\leanok
\uses{lem:busemann-approx-antitone,lem:busemann-limit,lem:busemann-lipschitz,lem:busemann-le-approx,lem:busemann-along-ray}
Throughout we use the basic properties of the family
$B_{\lambda,t}(y)=d(\lambda(t),y)-t$ established above: $B_\lambda$
is $1$-Lipschitz (Lemma~\ref{lem:busemann-lipschitz}):
\begin{equation}\label{eqnBlip}
|B_\lambda(y)-B_\lambda(z)|\le d(y,z)\ \ \text{for all } y,z\in M,
\end{equation}
for each fixed $y$ the function $t\mapsto B_{\lambda,t}(y)$ is
non-increasing on $[0,\infty)$
(Lemma~\ref{lem:busemann-approx-antitone}) with limit
$B_\lambda(y)$ (Lemma~\ref{lem:busemann-limit}), and
(Lemma~\ref{lem:busemann-le-approx})
\begin{equation}\label{eqnBmono}
B_\lambda(y)\le B_{\lambda,t}(y)\ \ \text{for all } y\in M \text{
and all } t\ge 0.
\end{equation}
Fix a witnessing sequence for $\mu$: $t_n\rightarrow\infty$ and
unit-speed minimizing segments $\sigma_n\colon[0,d_n]\to M$ from
$x$ to $\lambda(t_n)$, $d_n=d(x,\lambda(t_n))$, converging to $\mu$
uniformly on compact subsets of $[0,\infty)$; as recorded in
Definition~\ref{def:asymptotic-ray}, $d_n\ge
t_n-d(x,\lambda(0))$, so for each $s\ge0$ we have $s\le d_n$ for
all large $n$.

{\em Step 1: for all $s\ge 0$ and all $n$ with $0\le t_n$ and
$s\le d_n$,}
$$B_{\lambda,t_n}(\sigma_n(s))=B_{\lambda,t_n}(x)-s.$$
Since $\sigma_n$ is a unit-speed minimizing geodesic on $[0,d_n]$
with $\sigma_n(d_n)=\lambda(t_n)$, we have directly
$d(\sigma_n(s),\lambda(t_n))=d(\sigma_n(s),\sigma_n(d_n))=d_n-s$,
and $d(\lambda(t_n),x)=d(\lambda(t_n),\sigma_n(0))=d_n$. Therefore
$$B_{\lambda,t_n}(\sigma_n(s))=d(\lambda(t_n),\sigma_n(s))-t_n
=(d_n-t_n)-s=B_{\lambda,t_n}(x)-s.$$

{\em Step 2: $B_\lambda(\mu(s))\le B_\lambda(x)-s$ for all $s\ge
0$.} Fix $s\ge 0$ and let $\epsilon>0$. Since $t\mapsto
B_{\lambda,t}(x)$ is non-increasing on $[0,\infty)$ with limit
$B_\lambda(x)$, there is $T\ge0$ such that
$$B_{\lambda,T}(x)\le B_\lambda(x)+\epsilon.$$
Since $t_n\rightarrow\infty$ (hence also $d_n\rightarrow\infty$)
and $\sigma_n(s)\rightarrow\mu(s)$ (uniform convergence on
compacts implies pointwise convergence), we may choose $n$ with
$s\le d_n$, $T\le t_n$ and $d(\sigma_n(s),\mu(s))\le\epsilon$.
Then
$$B_\lambda(\mu(s))\le B_\lambda(\sigma_n(s))+\epsilon\le
B_{\lambda,t_n}(\sigma_n(s))+\epsilon
=B_{\lambda,t_n}(x)-s+\epsilon\le
B_{\lambda,T}(x)-s+\epsilon\le B_\lambda(x)-s+2\epsilon,$$
where the first inequality is the Lipschitz
estimate~(\ref{eqnBlip}), the second is the monotonicity
estimate~(\ref{eqnBmono}) at the point $\sigma_n(s)$ with $t=t_n$,
the equality is Step 1, the third inequality is the monotonicity of
$t\mapsto B_{\lambda,t}(x)$ together with $0\le T\le t_n$, and the
last is the choice of $T$. Since $\epsilon>0$ was arbitrary,
$B_\lambda(\mu(s))\le B_\lambda(x)-s$.

{\em Step 3: the reverse inequality.} Since $\mu$ is a unit-speed
minimizing geodesic ray with $\mu(0)=x$, we have $d(x,\mu(s))=s$.
The Lipschitz estimate~(\ref{eqnBlip}) then gives
$$B_\lambda(x)-B_\lambda(\mu(s))\le d(x,\mu(s))=s,$$
i.e., $B_\lambda(\mu(s))\ge B_\lambda(x)-s$. Combining with Step 2
completes the proof.
\end{proof}

\begin{corollary}[Busemann upper bound from an asymptotic ray]
\label{cor:asymptotic-ray-busemann-bound}
\lean{MorganTianLib.IsAsymptoticRay.busemann_le}
\leanok
\uses{def:busemann-function,def:asymptotic-ray,lem:asymptotic-ray}
Let $M$ be a metric space, let $\lambda$ be a minimizing geodesic
ray in $M$, let $\mu$ be a ray asymptotic to $\lambda$, and write
$x=\mu(0)$. Then for every $y\in M$ and every $t\ge 0$,
$$B_\lambda(y)\le B_\lambda(x)-t+d(y,\mu(t)).$$
\end{corollary}

\begin{proof}
\leanok
\uses{lem:busemann-lipschitz,lem:asymptotic-ray}
By the Lipschitz bound of Lemma~\ref{lem:busemann-lipschitz} and
Lemma~\ref{lem:asymptotic-ray},
$$B_\lambda(y)\le B_\lambda(\mu(t))+d(y,\mu(t))
=B_\lambda(x)-t+d(y,\mu(t)).$$
\end{proof}

\begin{lemma}[Busemann gradient along an asymptotic ray]
\label{lem:busemann-gradient-norm-one}
\uses{def:busemann-function,def:asymptotic-ray,lem:asymptotic-ray-exists,lem:asymptotic-ray}
Let $M$ be a connected, complete Riemannian manifold, let $\lambda\colon
[0,\infty)\to M$ be a minimizing geodesic ray, let $x\in M$, and
let $\mu$ be a ray asymptotic to $\lambda$ emanating from $x$
(Definition~\ref{def:asymptotic-ray},
Lemma~\ref{lem:asymptotic-ray-exists}). If the Busemann function
$B_\lambda$ is differentiable at $x$, then
$$\nabla B_\lambda(x)=-\mu'(0);$$
in particular $|\nabla B_\lambda(x)|=1$.
\end{lemma}

\begin{proof}
\uses{def:geodesic,thm:levi-civita-connection}
That $B_\lambda$ is differentiable at $x$ means that in a smooth
coordinate chart around $x$ the function $B_\lambda$ admits a
first-order Taylor expansion at $x$; the resulting linear
functional $dB_\lambda|_x\colon T_xM\to\Ar$ does not depend on the
chart (composing with the smooth transition map preserves
first-order approximability and transforms the linear term by the
differential of the transition map), and the gradient $\nabla
B_\lambda(x)\in T_xM$ is the vector representing it:
$$\langle\nabla B_\lambda(x),v\rangle=dB_\lambda|_x(v)\ \
\text{for all } v\in T_xM.$$
Recall also from Definition~\ref{def:busemann-function} that
$|B_{\lambda,t}(y)-B_{\lambda,t}(z)|\le d(y,z)$ for all $t\ge0$;
letting $t\rightarrow\infty$ shows that $B_\lambda$ is
$1$-Lipschitz:
\begin{equation}\label{eqnBlipgrad}
|B_\lambda(y)-B_\lambda(z)|\le d(y,z)\ \ \text{for all } y,z\in M.
\end{equation}

{\em Step 1: chain rule at the point of differentiability.} Let
$c\colon [0,\delta)\to M$ be a curve with $c(0)=x$ which is
differentiable at $0$. We claim that $s\mapsto B_\lambda(c(s))$ has
a one-sided derivative at $s=0$ equal to $\langle\nabla
B_\lambda(x),c'(0)\rangle$. Work in a coordinate chart around $x$,
identifying a neighborhood of $x$ with an open subset of $\Ar^n$
and $x$ with the origin. Differentiability of $B_\lambda$ at $x$
means
$$B_\lambda(y)=B_\lambda(x)+dB_\lambda|_x(y)+R(y),\ \ \text{where }
R(y)/|y|\rightarrow 0 \text{ as } y\rightarrow 0,$$
and differentiability of $c$ at $0$ means $c(s)=s\,c'(0)+r(s)$ with
$r(s)/s\rightarrow 0$ as $s\rightarrow 0^+$. Substituting and using
the linearity of $dB_\lambda|_x$,
$$B_\lambda(c(s))=B_\lambda(x)+s\,dB_\lambda|_x(c'(0))
+dB_\lambda|_x(r(s))+R(c(s)).$$
Here $dB_\lambda|_x(r(s))=o(s)$ because $dB_\lambda|_x$ is linear,
hence bounded, and $r(s)=o(s)$; and $R(c(s))=o(s)$ because
$|c(s)|\le \left(|c'(0)|+1\right)s$ for small $s$ while
$R(y)=o(|y|)$. Dividing by $s$ and letting $s\rightarrow 0^+$
proves the claim:
$$\lim_{s\rightarrow 0^+}\frac{B_\lambda(c(s))-B_\lambda(x)}{s}
=dB_\lambda|_x(c'(0))=\langle\nabla B_\lambda(x),c'(0)\rangle.$$

{\em Step 2: derivative along $\mu$.} By
Lemma~\ref{lem:asymptotic-ray}(2) we have
$B_\lambda(\mu(s))=B_\lambda(x)-s$ for all $s\ge 0$. The ray $\mu$
is a geodesic, hence a smooth curve, with $\mu(0)=x$ and
$|\mu'(0)|=1$ (it has unit speed). Applying Step 1 with $c=\mu$
and differentiating $s\mapsto B_\lambda(x)-s$ at $s=0$ gives
$$\langle\nabla B_\lambda(x),\mu'(0)\rangle=-1.$$

{\em Step 3: $|\nabla B_\lambda(x)|\le 1$.} Let $v\in T_xM$ be any
unit vector and let $\gamma_v$ be the geodesic with
$\gamma_v(0)=x$ and $\gamma_v'(0)=v$, defined for $s$ in some
interval $[0,\delta)$ by the local existence theorem for the
geodesic equation (Definition~\ref{def:geodesic}). Since the
Levi-Civita connection is compatible with the metric
(Theorem~\ref{thm:levi-civita-connection}),
$$\frac{d}{ds}\langle\gamma_v',\gamma_v'\rangle
=2\langle\nabla_{\gamma_v'}\gamma_v',\gamma_v'\rangle=0,$$
so $\gamma_v$ has constant speed $|\gamma_v'|\equiv|v|=1$; hence
for $0\le s<\delta$ the arc $\gamma_v|_{[0,s]}$ has length $s$ and
therefore $d(x,\gamma_v(s))\le s$. By the Lipschitz
estimate~(\ref{eqnBlipgrad}),
$$\left|B_\lambda(\gamma_v(s))-B_\lambda(x)\right|
\le d(\gamma_v(s),x)\le s.$$
Dividing by $s$, letting $s\rightarrow 0^+$ and invoking Step 1
with $c=\gamma_v$ yields $|\langle\nabla B_\lambda(x),v\rangle|\le
1$. Since
$$|\nabla B_\lambda(x)|=\sup_{v\in T_xM,\,|v|=1}
\langle\nabla B_\lambda(x),v\rangle,$$
it follows that $|\nabla B_\lambda(x)|\le 1$.

{\em Step 4: conclusion.} Using Steps 2 and 3 and $|\mu'(0)|=1$ we
compute
$$\left|\nabla B_\lambda(x)+\mu'(0)\right|^2
=|\nabla B_\lambda(x)|^2
+2\langle\nabla B_\lambda(x),\mu'(0)\rangle+|\mu'(0)|^2
\le 1-2+1=0.$$
(This is the equality case of the Cauchy-Schwarz inequality: a
pairing of $-1$ between vectors of norms $\le 1$ and $=1$ forces
them to be opposite unit vectors.) Hence $\nabla
B_\lambda(x)=-\mu'(0)$, and in particular $|\nabla
B_\lambda(x)|=|\mu'(0)|=1$.
\end{proof}

\begin{proposition}[Busemann function has non-positive Laplacian]
\label{prop:busemann-laplacian}
\label{Blambda}
\uses{def:busemann-function,def:ricci-curvature}
\source{morgan-tian}
Suppose that $M$ is complete and of non-negative Ricci curvature.
Then, for any minimizing geodesic ray $\lambda$, the Busemann
function $B_\lambda$ satisfies $\Delta B_\lambda\le 0$ in the weak
sense.
\end{proposition}


\begin{proof}
\uses{def:busemann-function,ex:distance-function-laplacian,rem:laplacian-weak-sense,thm:hopf-rinow,def:geodesic,lem:laplacian-local-formula}
The argument is due to Schoen and Yau; it is Proposition~1.1 on
pp.~7--8 of \cite{SchoenYau}. Let $n$ be the dimension of $M$ and
set $p=\lambda(0)$. Recall from
Definition~\ref{def:busemann-function} that
$B_{\lambda,t}(x)=d(\lambda(t),x)-t$, that each $B_{\lambda,t}$
satisfies $|B_{\lambda,t}(x)-B_{\lambda,t}(y)|\le d(x,y)$ and
$B_{\lambda,t}(x)\ge-d(x,p)$, that $B_{\lambda,t}(p)=0$, i.e.,
$d(p,\lambda(t))=t$, and that for each fixed $x$ the function
$t\mapsto B_{\lambda,t}(x)$ is non-increasing and converges to
$B_\lambda(x)$ as $t\to\infty$. In particular $B_\lambda$, being a
pointwise limit of $1$-Lipschitz functions, is itself
$1$-Lipschitz. By Remark~\ref{rem:laplacian-weak-sense}, the
assertion $\triangle B_\lambda\le0$ in the weak sense means
precisely that
$$\int_MB_\lambda\,\triangle\varphi\,d\vol\le0$$
for every non-negative test function $\varphi$, i.e., every
non-negative compactly supported $C^\infty$-function; the
integration-by-parts identity of that remark, which is valid for
any Lipschitz function, rewrites the left-hand side as
$-\int_M\langle\nabla B_\lambda,\nabla\varphi\rangle\,d\vol$, so
this is also equivalent to the formulation of the weak inequality
in terms of the $L^2$-vector field $\nabla B_\lambda$. Fix such a
$\varphi$ and let $C\subset M$ be its (compact) support.

We first dispose of the degenerate case $n=1$, in which
Example~\ref{ex:distance-function-laplacian} is unavailable: its
hypothesis $n\ge2$ fails, and its conclusion is genuinely false in
dimension one, since for any $x_0\in\Ar$ and any test function
$\chi$ on $\Ar$ two
integrations by parts, performed separately on $(-\infty,x_0]$ and
on $[x_0,\infty)$, give
$\int_{\Ar}|x-x_0|\,\chi''(x)\,dx=2\chi(x_0)$, which is positive
when $\chi\ge0$ and $\chi(x_0)>0$.
By Theorem~\ref{thm:hopf-rinow}, $\lambda$
extends to a unit-speed geodesic $\gamma\colon\Ar\to M$. We claim
that $\gamma$ is a bijection onto $M$. It is surjective: since
$\dim M=1$ and $\gamma'$ never vanishes, $\gamma$ is a local
diffeomorphism, so its image $E=\gamma(\Ar)$ is open; $E$ is also
closed, for if $y$ lies in the closure of $E$, we may choose a unit
vector $w\in T_yM$ and let $\eta$ be the unit-speed geodesic with
$\eta(0)=y$, $\eta'(0)=w$, defined on all of $\Ar$
(Definition~\ref{def:geodesic} and Theorem~\ref{thm:hopf-rinow});
then $\eta$ restricted to a small interval $(-\delta,\delta)$ is
again a local diffeomorphism, so $V=\eta((-\delta,\delta))$ is an
open neighborhood of $y$, and there are $s_0$ and
$\sigma\in(-\delta,\delta)$ with $\gamma(s_0)=\eta(\sigma)$; as
both curves are unit-speed and $\dim M=1$, we get
$\gamma'(s_0)=\pm\eta'(\sigma)$, so by the uniqueness of geodesics
with prescribed initial point and velocity
(Definition~\ref{def:geodesic}) $\gamma(s_0+a)=\eta(\sigma\pm a)$
for all $a\in\Ar$, and taking $a=\mp\sigma$ gives
$y=\eta(0)=\gamma(s_0\mp\sigma)\in E$. Since $M$ is connected and
$E$ is non-empty, open and closed, $E=M$. The map $\gamma$ is also
injective: suppose
$\gamma(s_1)=\gamma(s_2)$ with $s_1<s_2$. Since $\dim M=1$ and
$|\gamma'|\equiv1$, we have $\gamma'(s_2)=\pm\gamma'(s_1)$. If
$\gamma'(s_2)=\gamma'(s_1)$, then $s\mapsto\gamma(s+(s_2-s_1))$ and
$\gamma$ are geodesics with the same value and velocity at $s=s_1$,
hence equal; so $\gamma$ is periodic and
$M=\gamma(\Ar)=\gamma([s_1,s_2])$ is compact, hence of finite
diameter, contradicting $d(p,\lambda(t))=t$ for all $t\ge0$. If $\gamma'(s_2)=-\gamma'(s_1)$, then
$u\mapsto\gamma(s_2-u)$ and $u\mapsto\gamma(s_1+u)$ are geodesics
with the same initial point and initial velocity, hence
$\gamma(s_2-u)=\gamma(s_1+u)$ for all $u$; differentiating at
$u=(s_2-s_1)/2$ gives $-\gamma'(m)=\gamma'(m)$ for
$m=(s_1+s_2)/2$, contradicting $|\gamma'(m)|=1$. Thus $\gamma$ is a
bijective local diffeomorphism, hence a
diffeomorphism, and $\langle\gamma',\gamma'\rangle\equiv1$ says
that it pulls the metric of $M$ back to the standard metric
$ds^2$ on $\Ar$. Consequently $\gamma$ matches lengths of curves
and therefore distances, $d(\gamma(s),\gamma(s'))=|s-s'|$, it
carries Lebesgue measure to the Riemannian measure, and in the
resulting global coordinate $s$ the Laplacian is $d^2/ds^2$
(Lemma~\ref{lem:laplacian-local-formula} with $g\equiv1$). Now
$B_{\lambda,t}(\gamma(s))=|t-s|-t=-s$ whenever $t\ge s$, so
$B_\lambda(\gamma(s))=-s$, and with
$\psi=\varphi\circ\gamma\in C^\infty_c(\Ar)$ we get
$$\int_MB_\lambda\,\triangle\varphi\,d\vol
=\int_{\Ar}(-s)\,\psi''(s)\,ds
=\bigl[-s\,\psi'(s)\bigr]_{-\infty}^{+\infty}
+\int_{\Ar}\psi'(s)\,ds=0,$$
both terms vanishing because $\psi$ has compact support. This
proves the proposition when $n=1$; from now on we assume $n\ge2$.

Fix $t>0$ and set $f_t(x)=d(\lambda(t),x)$, so that
$B_{\lambda,t}=f_t-t$. Since $M$ is complete of dimension $n\ge2$
and $\Ric\ge0$, Example~\ref{ex:distance-function-laplacian}
applied with base point $\lambda(t)$ gives $\triangle
f_t\le(n-1)/f_t$ in the weak sense, which by
Remark~\ref{rem:laplacian-weak-sense} means
$$\int_Mf_t\,\triangle\varphi\,d\vol
\le\int_M\frac{n-1}{f_t}\,\varphi\,d\vol.$$
This inequality transfers to the shift $B_{\lambda,t}=f_t-t$
because constants are weakly harmonic: applying the
integration-by-parts identity of
Remark~\ref{rem:laplacian-weak-sense} to the constant (in
particular Lipschitz) function $1$ yields
$\int_M\triangle\varphi\,d\vol
=-\int_M\langle\nabla1,\nabla\varphi\rangle\,d\vol=0$, and hence
$$\int_MB_{\lambda,t}\,\triangle\varphi\,d\vol
=\int_Mf_t\,\triangle\varphi\,d\vol
-t\int_M\triangle\varphi\,d\vol
=\int_Mf_t\,\triangle\varphi\,d\vol
\le\int_M\frac{n-1}{f_t}\,\varphi\,d\vol.$$

Next we show that the right-hand side tends to $0$ as
$t\to\infty$. Let $D=\max_{x\in C}d(p,x)$, which is finite since
$C$ is compact. For $x\in C$ the triangle inequality and
$d(p,\lambda(t))=t$ give
$$f_t(x)=d(\lambda(t),x)\ge d(\lambda(t),p)-d(p,x)\ge t-D.$$
Hence, for every $t>D$,
$$0\le\int_M\frac{n-1}{f_t}\,\varphi\,d\vol
=\int_C\frac{n-1}{f_t}\,\varphi\,d\vol
\le\frac{n-1}{t-D}\int_M\varphi\,d\vol,$$
and since $\varphi$, and with it $\int_M\varphi\,d\vol$, is fixed,
the last expression tends to $0$ as $t\to\infty$.

Finally we pass to the limit on the left-hand side. On $C$ we have,
for every $t\ge0$, the uniform bounds
$$-D\le-d(x,p)\le B_\lambda(x)\le B_{\lambda,t}(x)\le
B_{\lambda,0}(x)=d(p,x)\le D,$$
where the first and last inequalities hold by the definition of $D$, the
second follows by letting $t\to\infty$ in
$B_{\lambda,t}(x)\ge-d(x,p)$, and the middle two are the
monotonicity in $t$. Thus $0\le B_{\lambda,t}-B_\lambda\le2D$ on
$C$, and $B_{\lambda,t}-B_\lambda\to0$ pointwise as $t\to\infty$.
Since $C$ is compact it has finite volume, so the constant $2D$ is
integrable on $C$, and the dominated convergence theorem gives
$$\lim_{t\rightarrow\infty}
\int_C\bigl(B_{\lambda,t}-B_\lambda\bigr)\,d\vol=0.$$
(Alternatively, Dini's theorem applies: a monotone pointwise limit
of continuous functions with continuous limit function is uniform
on the compact set $C$; the dominated-convergence route has the
minor advantage of not using the continuity of $B_\lambda$.)
Consequently
$$\Bigl|\int_MB_{\lambda,t}\,\triangle\varphi\,d\vol
-\int_MB_\lambda\,\triangle\varphi\,d\vol\Bigr|
\le\Bigl(\max_M|\triangle\varphi|\Bigr)
\int_C\bigl(B_{\lambda,t}-B_\lambda\bigr)\,d\vol
\longrightarrow0\qquad\text{as }t\rightarrow\infty,$$
the maximum being finite because $\triangle\varphi$ is continuous
with compact support. Combining this convergence with the
inequality
$\int_MB_{\lambda,t}\,\triangle\varphi\,d\vol
\le\frac{n-1}{t-D}\int_M\varphi\,d\vol$,
valid for all $t>D$ by the two preceding paragraphs, we obtain
$$\int_MB_\lambda\,\triangle\varphi\,d\vol
=\lim_{t\rightarrow\infty}
\int_MB_{\lambda,t}\,\triangle\varphi\,d\vol
\le\lim_{t\rightarrow\infty}\frac{n-1}{t-D}\int_M\varphi\,d\vol=0.$$
Since $\varphi$ was an arbitrary non-negative test function, this
is precisely the statement that $\triangle B_\lambda\le0$ in the
weak sense.
\end{proof}

\section{Comparison results in non-negative curvature}

Let us review some elementary comparison results for manifolds of
non-negative curvature. These form the basis for Toponogov
theory, \cite{Toponogov}. For any pair
of points $x,y$ in a complete Riemannian manifold $s_{xy}$ denotes a
minimizing geodesic from $x$ to $y$. We set $|s_{xy}|=d(x,y)$ and
call it the {\sl length} of the side. A {\sl triangle} in a
Riemannian manifold consists of three vertices $a,b,c$ and three
sides $s_{ab}$,$s_{ac}$,$s_{bc}$. We denote by $\angle_a$ the angle
of the triangle at $a$, i.e., the angle at $a$ between the geodesic
rays $s_{ab}$ and $s_{ac}$.

\begin{theorem}[Length comparison theorem]
\label{thm:length-comparison}
\label{lengthcompar}
\uses{def:sectional-curvature}
{\bf (Length comparison)}
Let $(M,g)$ be a manifold of non-negative curvature. Suppose that
$\triangle(a,b,c)$ is a triangle in $M$ and let
$\triangle(a',b',c')$ be a Euclidean triangle.
\begin{enumerate}
\item Suppose that the corresponding
sides of $\triangle(a,b,c)$ and $\triangle(a',b',c')$ have the same
lengths. Then the angle at each vertex of the Euclidean triangle is
no larger than the corresponding angle of $\triangle(a,b,c)$.
Furthermore, for any $\alpha$ and $\beta$ less than $|s_{ab}|$ and
$|s_{ac}|$ respectively, let $x$, resp. $x'$, be the point on
$s_{ab}$, resp. $s_{a'b'}$, at distance $\alpha$ from $a$, resp.
$a'$, and let $y$, resp. $y'$, be the point on $s_{ac}$, resp.
$s_{a'c'}$, at distance $\beta$ from $a$, resp. $a'$. Then
$d(x,y)\ge d(x',y')$.
\item Suppose that $|s_{ab}|=|s_{a'b'}|$, that
$|s_{ac}|=|s_{a'c'}|$ and that $\angle_a=\angle_{a'}$. Then
$|s_{b'c'}|\ge |s_{bc}|$.
\end{enumerate}
\end{theorem}

See the figure in [Morgan--Tian]. For a proof of this result see
Theorem 4.2 on page 161 of \cite{Sakai}, or Theorem 2.2 on page 42
of \cite{CheegerEbin}.


One corollary is a monotonicity result. Suppose that
$\triangle(a,b,c)$ is a triangle in a complete manifold of
non-negative curvature. Define a function $EA(u,v)$ defined for
$0\le u\le |s_{ab}|$ and $0\le v\le |s_{ac}|$ as follows. For $u$
and $v$ in the indicated ranges, let $x(u)$ be the point on $s_{ab}$
at distance $u$ from $a$ and let $y(v)$ be the point of $s_{ac}$ at
distance $v$ from $a$. Let $EA(u,v)$ be the angle at $a'$ of the
Euclidean triangle with side lengths $|s_{a'b'}|=u$, $|s_{a'c'}|=v$
and $|s_{b'c'}|=d(x(u),y(v))$.

\begin{corollary}[Angle monotonicity in non-negative curvature]
\label{cor:angle-monotonicity}
\label{anglemono}
\uses{thm:length-comparison}
Under the assumptions of the previous theorem, $EA(u,v)$ is a
monotone non-increasing function of each variable $u$ and $v$ when
the other variable is held fixed.
\end{corollary}



Suppose that $\alpha, \beta,\gamma$ are three geodesics
emanating from a point $p$ in a Riemannian manifold. Let $\angle_p
(\alpha,\beta)$, $\angle_p(\beta,\gamma)$ and
$\angle_p(\alpha,\gamma)$ be the angles of these geodesics at $p$ as
measured by the Riemannian metric. Then of course
$$\angle_p(\alpha,\beta)+\angle_p(\beta,\gamma)+\angle_p(\alpha,\gamma)\le
2\pi$$ since this inequality holds for the angles between straight
lines in Euclidean $n$-space. There is a second corollary of
Theorem~\ref{lengthcompar} which gives an analogous result for the
associated Euclidean angles.



\begin{corollary}[Angle comparison in non-negative curvature]
\label{cor:angle-comparison}
\label{anglecompar}
\uses{thm:length-comparison}
Let $(M,g)$ be a complete Riemannian manifold of non-negative
curvature. Let $p,a,b,c$ be four points in $M$ and let
$\alpha,\beta,\gamma$ be minimizing geodesic arcs from the point
$p$ to $a,b,c$ respectively. Let $T(a,p,b)$, $T(b,p,c)$ and
$T(c,p,a)$ be the triangles in $M$ made out of these minimizing
geodesics and minimizing geodesics between $a,b,c$. Let
$T(a',p',b')$, $T(b',p',c')$ and $T(c',p',a')$ be planar triangles
with the same side lengths. Then
$$\angle_{p'}T(a',p',b')+\angle_{p'}T(b',p',c')+\angle_{p'}T(c',p',a')\le 2\pi.$$
\end{corollary}

\begin{proof}
\uses{thm:length-comparison}
Consider the sum of these angles as the geodesic arcs in $M$ are
shortened without changing their direction. By the first property of
Theorem~\ref{thm:length-comparison} the sum of the angles of these triangles
is a monotone decreasing function of the lengths. Of course, the
limit as the lengths all go to zero is the corresponding Euclidean
angle. The result is now clear.
\end{proof}



\section{The soul theorem}

A subset $X$ of a Riemannian manifold $(M,g)$ is said to be {\em
totally convex} if every geodesic segment with endpoints in $X$ is
contained in $X$. Thus, a point $p$ in $M$ is totally convex if and
only if there is no broken geodesic arc in $M$ broken exactly at
$x$.


\begin{theorem}[Soul Theorem (Cheeger-Gromoll)]
\label{thm:soul-theorem}
\label{soul}
\uses{def:sectional-curvature}
(Cheeger-Gromoll, see
\cite{CheegerGromollbul} and \cite{CheegerGromoll}) Suppose that
$(M,g)$ is a connected, complete, non-compact Riemannian manifold of
non-negative sectional curvature. Then $M$ contains a soul $S\subset
M$. By definition a {\em soul} is a compact, totally
geodesic, totally convex submanifold (automatically of positive
codimension). Furthermore, $M$ is diffeomorphic to the total space
of the normal bundle of the $S$ in $M$. If $(M,g)$ has positive
curvature, then any soul for it is a point, and consequently $M$ is
diffeomorphic to $\Ar^n$.
\end{theorem}


\begin{remark}[Soul theorem for positive curvature]
\label{rem:soul-theorem-positive-curvature}
\uses{thm:soul-theorem}
We only use the soul theorem for manifolds with positive curvature
and the fact that any soul of such a manifold is a point. A proof of
this result first appears in \cite{GromollMeyer}.
\end{remark}



The rest of this section is devoted to a sketch of the proof of this
result. Our discussion follows closely that in \cite{Petersen}
starting on p.~349. We shall need more information about complete,
non-compact manifolds of non-negative curvature, so we review a
little of their theory as we sketch the proof of the soul theorem.

\begin{lemma}[Angles of geodesics converge at infinity]
\label{lem:angle-convergence-infinity}
\uses{def:sectional-curvature}
Let $(M,g)$ be a complete, non-compact Riemannian manifold of
non-negative sectional curvature and let $p\in M$. For every
$\epsilon>0$ there is a compact subset $K=K(p,\epsilon)\subset M$
such that for all points $q\notin K$, if $\gamma$ and $\mu$ are
minimizing geodesics from $p$ to $q$, then the angle that $\gamma$
and $\mu$ make at $q$ is less than $\epsilon$.
\end{lemma}

See the figure in [Morgan--Tian].

\begin{proof}
\uses{thm:length-comparison}
The proof is by contradiction. Fix $0<\epsilon<1$ sufficiently small
so that $\cos(\epsilon/2)<1-\epsilon^2/12$. Suppose that there
is a sequence of points $q_n$ tending to infinity such that for each
$n$ there are minimizing geodesics $\gamma_n$ and $\mu_n$ from $p$
to $q_n$ making angle at least $\epsilon$ at $q_n$. For each $n$ let
$d_n=d(p,q_n)$. By passing to a subsequence we can suppose that for
all $n$ and $m$ the cosine of the angle at $p$ between $\gamma_n$
and $\gamma_m$ at least $1-\epsilon^2/24$, and the cosine of the
angle at $p$ between $\mu_n$ and $\mu_m$ is at least
$1-\epsilon^2/24$. We can also assume that for all $n\ge 1$ we have
$d_{n+1}\ge (100/\epsilon^2) d_n$. Let $\delta_n=d(q_n,q_{n+1})$.
Applying the first Toponogov property (Theorem~\ref{thm:length-comparison}) at $p$, we see that
$\delta_n^2\le d_n^2+d_{n+1}^2-2d_nd_{n+1}(1-\epsilon^2/24)$.
Applying the same property at $q_n$ we have
$$d_{n+1}^2\le d_n^2+\delta_n^2-2d_n\delta_n\cos(\theta),$$ where
$\theta\le \pi$ is the angle at $q_n$ between $\gamma_n$ and a
minimal geodesic joining $q_n$ to $q_{n+1}$. Thus,
$$\cos(\theta)\le
\frac{d_n-d_{n+1}(1-\epsilon^2/24)}{\delta_n}.$$ By the triangle
inequality (and the fact that $\epsilon<1$) we have $\delta_n\ge
(99/\epsilon)d_n$ and $\delta_n\ge d_{n+1}(1-(\epsilon^2/100))$.
Thus,
$$\cos(\theta)\le
\epsilon^2/99-(1-\epsilon^2/24)/(1-(\epsilon^2/100))<-(1-\epsilon^2/12).$$
This implies that $\cos(\pi-\theta)>(1-\epsilon^2/12)$, which
implies that $\pi-\theta<\epsilon/2$. That is to say, the angle at
$q_n$ between $\gamma_n$ and a shortest geodesic from $q_n$ to
$q_{n+1}$ is between $\pi-\epsilon/2$ and $\pi$. By symmetry, the
same is true for the angle between $\mu_n$ and the same shortest
geodesic from $q_n$ to $q_{n+1}$. Thus, the angle between $\gamma_n$
and $\mu_n$ at $q_n$ is less than $\epsilon$, contradicting our
assumption.
\end{proof}




\begin{corollary}[Distance level sets are homeomorphic]
\label{cor:distance-level-sets-homeomorphic}
\label{homeo}
\uses{lem:angle-convergence-infinity}
Let $(M,g)$ be a complete, non-compact manifold of non-negative
sectional curvature. Let $p\in M$ and define a function $f\colon
M\to \Ar$ by $f(q)=d(p,q)$. Then there is $R<\infty$ such that for
$R\le s<s'$ we have:
\begin{enumerate}
\item $f^{-1}([s,s'])$ is homeomorphic to $f^{-1}(s)\times[s,s']$ and
in particular the level sets $f^{-1}(s)$ and $f^{-1}(s')$ are
homeomorphic; \item $f^{-1}([s,\infty)$ is homeomorphic to
$f^{-1}(s)\times[s,\infty)$.
\end{enumerate}
\end{corollary}

\begin{proof}
Given $(M,g)$ and $p\in M$ as in the statement of the corollary, choose a
constant $R<\infty$ such that any two minimal geodesics from $p$ to a point $q$
with $d(p,q)\ge R/2$ make an angle at most $\pi/6$ at $q$. Now following
\cite{Petersen} p. 335, it is possible to find a smooth unit vector field $X$
on $U=M-\overline{B(p,R/2)}$ with the property that $f(\cdot)=d(p,\cdot)$ is
increasing along any integral curve for $X$ at a rate bounded below by
$\cos(\pi/3)$. In particular, for any $s\ge R$ each integral curve of $X$
crosses the level set $f^{-1}(s)$ in a single point. Using this vector field
we see that for any $s,s'>R$, the pre-image $f^{-1}([s,s'])$ is homeomorphic to
$f^{-1}(s)\times[s,s']$ and that the end $f^{-1}\left([s,\infty)\right)$ is
homeomorphic to $f^{-1}(s)\times [s,\infty)$.
\end{proof}


In a complete, non-compact $n$-manifold of positive curvature any soul is a
point. While the proof of this result uses the same ideas as discussed above,
we shall not give a proof. Rather we refer the reader to Theorem 84 of
\cite{Petersen} on p. 349. A soul has the property that if two minimal
geodesics emanate from $p$ and end at the same point $q\not= p$, then the angle
that they make at $q$ is less than $\pi/2$. Also, of course, the exponential
mapping is a diffeomorphism sufficiently close to the soul. Applying the above
lemma and a standard compactness argument, we see that in fact there is
$\epsilon>0$ such that all such pairs of minimal geodesics from $p$ ending at
the same point make angle less than $\pi/2-\epsilon$ at that point. Hence, in
this case there is a vector field $X$ on all of $M$ vanishing only at the soul,
and agreeing with the gradient of the distance function near the soul, so that
the distance function from $p$ is strictly increasing to infinity along each
flow line of $X$ (except the fixed point). Using $X$ one establishes that $M$
is diffeomorphic to $\Ar^n$. It also follows that all the level surfaces
$f^{-1}(s)$ for $s>0$ are homeomorphic to $S^{n-1}$ and for $0<s<s'$ the
preimage $f^{-1}([s,s'])$ is homeomorphic to $S^{n-1}\times[s,s']$.

There is an analogue of this result for the distance function from
any point, not just a soul.

\begin{corollary}[Distance annuli are homotopy spheres]
\label{cor:distance-annuli-homotopy-sphere}
\label{htpytype}
\uses{cor:distance-level-sets-homeomorphic,thm:soul-theorem}
Let $(M,g)$ be a complete, non-compact Riemannian $n$-manifold of positive
curvature. Then for any point $p\in M$ there is a constant $R=R(p)$ such that
for any $s<s'$ with $R\le s$ both $f^{-1}(s,s')$ and $f^{-1}(s,\infty)$ are
homotopy equivalent to $S^{n-1}$.
\end{corollary}


\begin{proof}
\uses{cor:distance-level-sets-homeomorphic}
Given $(M,g)$ and $p$ fix $R<\infty$ sufficiently large so that
Corollary~\ref{cor:distance-level-sets-homeomorphic} holds. Since $M$ is diffeomorphic to $\Ar^n$
it has only one end and hence the level sets $f^{-1}(s)$ for $s\ge
R$ are connected. Given any compact subset $K\subset M$ there is a
larger compact set $B$ (a ball) such that $M\setminus B$ has trivial
fundamental group and trivial homology groups $H_i$ for $i<n-1$.
Hence for any subset $Z\subset M\setminus B$, the inclusion of $Z\to
M\setminus K$ induces the trivial map on $\pi_1$ and on $H_i$ for
$i<n-1$. Clearly, for any $R\le s<b$ the inclusion
$f^{-1}(b,\infty)\to f^{-1}(s,\infty)$ is a homotopy equivalence.
Thus, it must be the case that $f^{-1}(b,\infty)$ has trivial
fundamental group and $H_i$ for $i<n-1$. Hence, the same is true for
$f^{-1}(s,\infty)$ for any $s\ge R$. Lastly, since
$f^{-1}(s,\infty)$ is connected and simply connected ( hence
orientable) and has two ends, it follows by the non-compact form of
Poincar\'e duality that $H_{n-1}(f^{-1}(s,\infty))\cong \Zee$.
Hence, by the Hurewicz theorem $f^{-1}(s,\infty)$ is homotopy
equivalent to $S^{n-1}$ for any $s\ge R$. Of course, it is also true
for $R\le s\le s'$ that $f^{-1}(s,s')$ is homotopy equivalent to
$S^{n-1}$.
\end{proof}

\section{Ends of a manifold}
\label{ENDS}

Let us review the basic notions about ends of a manifold.


\begin{definition}[Transition maps between complementary components]
\label{def:ends-transition}
\lean{MorganTianLib.endsTransition}
\leanok
\source{morgan-tian}
Let $X$ be a topological space. For compact subsets $K\subset K'$ of
$X$, the inclusion $X\setminus K'\subset X\setminus K$ induces a
{\em transition map}
$$\pi_0(X\setminus K')\longrightarrow \pi_0(X\setminus K)$$
sending a component of $X\setminus K'$ to the component of
$X\setminus K$ containing it. Here $\pi_0(Y)$ denotes the set of
connected components of $Y$, endowed with the quotient topology.
\end{definition}

\begin{lemma}[The transition maps are continuous]
\label{lem:ends-transition-continuous}
\lean{MorganTianLib.continuous_endsTransition}
\leanok
\uses{def:ends-transition}
For compact subsets $K\subset K'$ of $X$, the transition map
$\pi_0(X\setminus K')\to\pi_0(X\setminus K)$ is continuous.
\end{lemma}

\begin{proof}
\leanok
The composition $X\setminus K'\hookrightarrow X\setminus
K\rightarrow\pi_0(X\setminus K)$ is continuous and constant on the
connected components of $X\setminus K'$ (the continuous image of a
connected set is connected), so it descends to a continuous map on
the quotient $\pi_0(X\setminus K')$.
\end{proof}

\begin{lemma}[Transition maps: identity]
\label{lem:ends-transition-id}
\lean{MorganTianLib.endsTransition_refl}
\leanok
\uses{def:ends-transition}
For each compact $K\subset X$, the transition map associated to
$K\subset K$ is the identity of $\pi_0(X\setminus K)$.
\end{lemma}

\begin{proof}
\leanok
The component of $X\setminus K$ containing a given component is that
component itself.
\end{proof}

\begin{lemma}[Transition maps: composition]
\label{lem:ends-transition-comp}
\lean{MorganTianLib.endsTransition_trans}
\leanok
\uses{def:ends-transition}
For compact subsets $K\subset K'\subset K''$ of $X$, the transition
map of $K\subset K''$ is the composition of the transition map of
$K'\subset K''$ followed by that of $K\subset K'$. Thus the sets
$\pi_0(X\setminus K)$, indexed by the compact subsets $K\subset X$
directed by inclusion, form an inverse system.
\end{lemma}

\begin{proof}
\leanok
A component of $X\setminus K''$ is contained in a unique component
of $X\setminus K'$, which in turn is contained in a unique component
of $X\setminus K$; containment is transitive.
\end{proof}

\begin{definition}[Space of ends]
\label{def:space-of-ends}
\lean{MorganTianLib.SpaceOfEnds}
\leanok
\uses{def:ends-transition,lem:ends-transition-comp,lem:compact-set-in-codim-zero-submanifold}
\source{morgan-tian}
Let $M$ be a connected manifold. The {\em space of ends} of $M$ is
the inverse limit of the inverse system $K\mapsto\pi_0(M\setminus
K)$ of Definition~\ref{def:ends-transition}, indexed by the compact
subsets $K\subset M$ directed by inclusion: its points are the
families ${\mathcal E}=(e_K)_K$, with $e_K$ a component of
$M\setminus K$, that are compatible under the transition maps; it is
topologized as a subspace of the product $\prod_K\pi_0(M\setminus
K)$. An {\em end} of $M$ is a point of the space of ends.

Classically the system is indexed instead by the compact,
codimension-$0$ submanifolds $K\subset M$, for which each
$\pi_0(M\setminus K)$ is a finite set; by
Lemma~\ref{lem:compact-set-in-codim-zero-submanifold} these are
cofinal among all compact subsets, so the two inverse systems have
the same inverse limit. Note also that $M\setminus K$ is an open
subset of a manifold, hence locally connected, so its components are
open and $\pi_0(M\setminus K)$ with the quotient topology is
discrete.

An end ${\mathcal E}$ determines a complementary component of each
compact, codimension-$0$ submanifold $K\subset M$, called a {\em
neighborhood} of the end. Conversely, by definition these
neighborhoods are cofinal in the set of all neighborhoods of the
end. A sequence $\{x_n\}$ in $M$ {\em converges to the end}
${\mathcal E}$ if it is eventually in every neighborhood of the end.
\end{definition}

\begin{remark}[The end compactification]
\label{rem:ends-compactification}
\uses{def:space-of-ends}
The space of ends of a connected manifold is a compact space, and
one can define a topology on the union of $M$ and its space of ends
making this union a compact, connected Hausdorff space which is a
compactification of $M$. We will not use these facts in this
chapter.
\end{remark}

\begin{lemma}[Compact sets lie in compact codimension-zero submanifolds]
\label{lem:compact-set-in-codim-zero-submanifold}
Let $M$ be a smooth $n$-manifold. Then every compact subset $C\subset
M$ is contained in a compact, codimension-zero, smooth submanifold
with boundary $K\subset M$.
\end{lemma}

\begin{proof}
Recall the standard bump function: the function on $\Ar^n$ equal to
$e^{-1/(1-|u|^2)}$ for $|u|<1$ and to $0$ for $|u|\ge1$ is smooth ---
all derivatives of $s\mapsto e^{-1/s}$ tend to $0$ as
$s\rightarrow0^+$ --- non-negative, positive exactly on the open unit
ball, and supported in the closed unit ball.

For each $x\in C$ choose a coordinate chart around $x$ whose domain
contains the closed unit coordinate ball centered at $x$, and let
$\psi_x\colon M\to[0,\infty)$ be the function equal in these
coordinates to the standard bump and equal to $0$ outside the chart;
then $\psi_x$ is smooth, has compact support, and is positive on an
open neighborhood $V_x$ of $x$ (the open unit coordinate ball). By
compactness, $C\subset V_{x_1}\cup\cdots\cup V_{x_m}$ for finitely
many points $x_1,\dots,x_m\in C$. Put
$\psi=\psi_{x_1}+\cdots+\psi_{x_m}$: it is smooth, non-negative, has
compact support, and is positive on an open set containing $C$. Since
$C$ is compact and $\psi$ is continuous and positive on $C$, the
number $\delta:=\min_C\psi$ is positive.

Let $\theta\colon\Ar\to[0,\infty)$ be smooth and positive exactly on
the interval $(\delta/3,2\delta/3)$ (a rescaled one-variable standard
bump), and define
$$h(s)=\frac{\int_{-\infty}^{s}\theta(u)\,du}
{\int_{-\infty}^{\infty}\theta(u)\,du}.$$
Then $h\colon\Ar\to[0,1]$ is smooth and non-decreasing, with
$h\equiv0$ on $(-\infty,\delta/3]$ and $h\equiv1$ on
$[2\delta/3,\infty)$. The composition $\varphi=h\circ\psi\colon
M\to[0,1]$ is smooth; it satisfies $\varphi\equiv1$ on the open set
$\{\psi>2\delta/3\}$, which contains $C$ (as $\psi\ge\delta$ there),
and $\varphi=0$ outside the compact set $\supp\psi$.

By Sard's theorem the critical values of $\varphi$ form a set of
Lebesgue measure zero in $\Ar$, so we may choose a regular value
$c\in(0,1)$ of $\varphi$. Set
$$K:=\varphi^{-1}([c,1])=\{\varphi\ge c\}.$$
Then $K\supset C$, since $\varphi\equiv1$ on $C$ and $c<1$; and $K$
is compact, being a closed subset of the compact set $\supp\psi$
(indeed $\varphi=0$ off $\supp\psi$ and $c>0$). Finally, $K$ is a
smooth $n$-dimensional submanifold with boundary
$\varphi^{-1}(c)$: each point of the open set $\{\varphi>c\}$ is an
interior point of $K$; and at each point $z$ with $\varphi(z)=c$ we
have $d\varphi_z\neq0$ because $c$ is a regular value, so
$\varphi-c$ may be completed to a system of local coordinates near
$z$, in which $K$ is the half-space where the first coordinate is
non-negative. (If $\varphi^{-1}(c)$ is empty, the same conclusion
holds with empty boundary.) Thus $K$ is a compact, codimension-zero,
smooth submanifold with boundary containing $C$.
\end{proof}

A proper map between topological manifolds induces a map on the
space of ends, and in fact induces a map on the compactifications
sending the subspace of ends of the compactification of the domain
to the subspace of ends of the compactification of the range.


We say that a path $\gamma\colon [a,b)\to M$ is a path {\em to the
end ${\mathcal E}$} if it is a proper map and it sends the end
$\{b\}$ of $[a,b)$ to the end ${\mathcal E}$ of $M$. This condition
is equivalent to saying that given a neighborhood $U$ of ${\mathcal
E}$ there is a neighborhood of the end $\{b\}$ of $[a,b)$ that maps
to $U$.

Now suppose that $M$ has a Riemannian metric $g$. Then we can
distinguish between ends at finite and infinite distance. An end is
at {\em finite distance} if there is a rectifiable path of finite
length to the end. Otherwise, the end is at infinite distance. If an
end is at finite distance we have the notion of the distance from a
point $x\in M$ to the end. It is the infimum of the lengths of
rectifiable paths from $x$ to the end. This distance is always
positive. Also, notice that the Riemannian manifold is complete if
and only if has no end at finite distance.

\section{The splitting theorem}

In this section we give a proof of the following theorem which is
originally due to Cheeger-Gromoll \cite{CheegerGromoll2}. The weaker
version giving the same conclusion under the stronger hypothesis of
non-negative sectional curvature (which is in fact all we need in
this work) was proved earlier by Toponogov, see \cite{Toponogov}.

The proof of the splitting theorem rests on the maximum principle
for functions that are subharmonic in the weak sense, and on the
regularity theorem for weakly harmonic functions; we develop both
now. Exactly two ingredients are taken from the theory of elliptic
partial differential equations with references in place of proofs;
they are stated precisely as
Claim~\ref{claim:dirichlet-problem-small-balls} and
Claim~\ref{claim:weak-subsolution-comparison} below.

\begin{lemma}[Green identity for compactly supported test functions]
\label{lem:green-identity-compact-support}
\uses{lem:laplacian-local-formula}
Let $(M,g)$ be a Riemannian manifold, let $\Omega\subset M$ be an
open set, let $u$ be a smooth function on $\Omega$, and let $\psi$
be a compactly supported $C^\infty$-function on $\Omega$. Then
$$\int_\Omega\left(u\,\Delta\psi-\psi\,\Delta u\right)d\vol=0.$$
\end{lemma}

\begin{proof}
\uses{lem:laplacian-local-formula}
First suppose that the support of $\psi$ is contained in the domain
$V\subset\Omega$ of a single coordinate chart. By
Lemma~\ref{lem:laplacian-local-formula}, in the coordinates on $V$
we have, writing $\gamma:=\sqrt{\det(g)}$,
$$u\,\Delta\psi-\psi\,\Delta u
=\frac{1}{\gamma}\left(u\,\partial_i\!\left(g^{ij}\gamma\,
\partial_j\psi\right)-\psi\,\partial_i\!\left(g^{ij}\gamma\,
\partial_ju\right)\right)
=\frac{1}{\gamma}\,\partial_i\!\left(g^{ij}\gamma\left(u\,
\partial_j\psi-\psi\,\partial_ju\right)\right),$$
where the second equality holds because expanding the last
expression by the product rule produces the two displayed terms
together with the cross terms $g^{ij}\gamma\,\partial_iu\,
\partial_j\psi$ and $-g^{ij}\gamma\,\partial_i\psi\,\partial_ju$,
which cancel by the symmetry $g^{ij}=g^{ji}$ (relabel the summation
indices). Since $d\vol=\gamma\,dx^1\cdots dx^n$ in the chart, and
since the functions $g^{ij}\gamma\left(u\partial_j\psi-\psi
\partial_ju\right)$ are smooth with compact support contained in
the coordinate image of $V$ (they vanish off the support of
$\psi$), integrating in the coordinates and applying Fubini's
theorem and the fundamental theorem of calculus in each variable
$x^i$ gives
$$\int_\Omega\left(u\Delta\psi-\psi\Delta u\right)d\vol
=\int\partial_i\!\left(g^{ij}\gamma\left(u\partial_j\psi
-\psi\partial_ju\right)\right)dx^1\cdots dx^n=0.$$
For general $\psi$, cover the compact set $\supp(\psi)$ by
finitely many chart domains $V_1,\dots,V_N\subset\Omega$ and choose
a smooth partition of unity $\chi_1,\dots,\chi_N$ with
$\chi_k\ge0$, $\supp(\chi_k)\subset V_k$ compact, and
$\sum_k\chi_k=1$ on a neighborhood of $\supp(\psi)$. Then
$\psi=\sum_k\chi_k\psi$, each $\chi_k\psi$ is a compactly supported
$C^\infty$-function with support in the single chart $V_k$, and the
asserted identity follows by summing the single-chart case over
$k$, using the linearity of $\Delta$ and of the integral.
\end{proof}

\begin{claim}[Dirichlet problem on small geodesic balls]
\label{claim:dirichlet-problem-small-balls}
\uses{def:exponential-map,cor:exponential-local-diffeomorphism,lem:laplacian-local-formula}
Let $(M,g)$ be a Riemannian manifold of dimension $n\ge1$ and let
$p\in M$. Then there is $r_D=r_D(p)>0$ such that for every
$0<s\le r_D$ the following hold.
\begin{enumerate}
\item The metric ball $B(p,s)=\{x\in M\,|\,d(p,x)<s\}$ is the image
of the ball $B(0,s)\subset T_pM$ under the diffeomorphism furnished
by the exponential map
(Corollary~\ref{cor:exponential-local-diffeomorphism} and the
discussion following Definition~\ref{def:exponential-map}); in
particular $B(p,s)$ is connected and is contained in the domain of
a single coordinate chart (Gaussian normal coordinates centered at
$p$). Its closure is the compact set
$\overline{B}(p,s)=\{x\,|\,d(p,x)\le
s\}=\exp_p\left(\overline{B}(0,s)\right)$, and its topological
boundary is the metric sphere
$$S(p,s)=\{x\in M\,|\,d(p,x)=s\}=\exp_p\left(\{v\in
T_pM\,|\,|v|=s\}\right),$$ a non-empty smooth embedded hypersurface
diffeomorphic to $S^{n-1}$.
\item For every continuous function $\varphi\colon S(p,s)\to \Ar$
there is a function $h$, continuous on $\overline{B}(p,s)$ and
smooth on $B(p,s)$, with
$$\Delta h=0\ \text{ in } B(p,s)\qquad \text{and}\qquad
h|_{S(p,s)}=\varphi.$$
\end{enumerate}
This claim is the first of the two analytic claims of this section
that enter the blueprint as black boxes, with references in place
of a proof. (The regularity theorem for weakly harmonic functions,
Theorem~\ref{thm:weak-harmonic-regularity}, which closes this
section, is proved from these two claims.) The geometric assertions of Item (1)
constitute the normal-ball theorem; see
Corollary~\ref{cor:exponential-local-diffeomorphism} and the
discussion following Definition~\ref{def:exponential-map} together
with Chapter 3 of \cite{doCarmo} (Proposition 2.9 there: $\exp_p$
is a diffeomorphism of a small ball in $T_pM$ onto an open set; and
Proposition 3.6: inside such a normal ball the radial geodesics
minimize length, so that $d(p,\exp_p(v))=|v|$ for small $v$, whence
the displayed descriptions of the ball, of its closure and of the
boundary sphere). For Item (2), note that in Gaussian normal
coordinates the operator $\Delta$ becomes, by
Lemma~\ref{lem:laplacian-local-formula}, a second-order uniformly
elliptic operator with smooth coefficients on a Euclidean ball; the
solvability of the Dirichlet problem for such operators with
continuous boundary data, with smoothness in the interior, is
furnished by Perron's method together with interior Schauder
estimates: see Gilbarg--Trudinger, {\it Elliptic Partial
Differential Equations of Second Order}, Chapter 6 (in particular
Theorems 6.11 and 6.13 and the interior regularity results of that
chapter). A closely related fact---solvability of the Dirichlet
problem on a neighborhood with smooth boundary contained in a
coordinate chart---is invoked in the proof of the regularity
theorem for harmonic functions in \cite{Petersen}, p.~279ff.
\end{claim}

\begin{claim}[Comparison principle for continuous weak subsolutions]
\label{claim:weak-subsolution-comparison}
\uses{rem:laplacian-weak-sense,lem:laplacian-local-formula,lem:green-identity-compact-support}
Let $(M,g)$ be a Riemannian manifold and let $\Omega\subset M$ be a
non-empty {\it connected} open set whose closure $\overline{\Omega}$
is compact and whose topological boundary $\partial\Omega$ is
non-empty. Let $w\colon \overline{\Omega}\to\Ar$ be a continuous
function satisfying $\Delta w\ge 0$ in the weak sense on $\Omega$,
meaning (cf.\ Remark~\ref{rem:laplacian-weak-sense}) that for every
non-negative compactly supported $C^\infty$-function $\psi$ on
$\Omega$ we have
$$\int_\Omega w\,\Delta \psi\, d\vol\ \ge\ 0.$$
Then
$$\max_{\overline{\Omega}}w=\max_{\partial\Omega}w.$$
This claim is the second of the two analytic claims of this section
that enter the blueprint as black boxes. The connectedness
hypothesis cannot simply be dropped: if $M$ is the disjoint union
of a closed (compact, boundaryless) manifold $N$ and a Euclidean
space, and $\Omega$ is the union of $N$ with a small open ball $D$
in the second component, then the function $w$ equal to $1$ on $N$
and to $0$ on $\overline{D}$ is continuous on the compact set
$\overline{\Omega}$ and satisfies the weak inequality---for every
test function $\psi\ge0$ on $\Omega$ one has $\int_\Omega
w\,\Delta\psi\,d\vol=\int_N\Delta\psi\,d\vol=0$ by
Lemma~\ref{lem:green-identity-compact-support} applied on $N$ with
$u\equiv1$---yet $\max_{\overline{\Omega}}w=1>0
=\max_{\partial\Omega}w$. (More generally the statement holds
whenever every connected component of $\Omega$ has non-empty
topological boundary; the connected case stated here is all that is
used below.) We emphasize that the hypothesis here is the {\it
distributional} one displayed above; it is not the barrier
(support-function) sense in which subharmonicity is taken in the
treatment of the maximum principle in \cite{Petersen}, p.~279ff.,
and neither sense formally subsumes the other for merely continuous
functions, so the two developments must not be confused. In a
coordinate chart $\Delta$ is a second-order elliptic operator with
smooth coefficients (Lemma~\ref{lem:laplacian-local-formula}), and
for such operators the statement---indeed the stronger assertion
that a continuous distributional subsolution obeys the strong
maximum principle---is proved by W.~Littman, {\it A strong maximum
principle for weakly $L$-subharmonic functions}, J.~Math.\ Mech.\
{\bf 8} (1959), 761--770; for subsolutions with locally
square-integrable derivatives the corresponding weak maximum
principle is Gilbarg--Trudinger, {\it op.\ cit.}, Chapter 8
(Theorem 8.1). (Only the boundary-maximum statement above is
consumed by this blueprint; the strong maximum principle is {\it
proved} below, in Theorem~\ref{thm:maximum-principle}, from this
claim and Claim~\ref{claim:dirichlet-problem-small-balls}.)
\end{claim}

\begin{lemma}[Second-derivative test, minimum version]
\label{lem:second-derivative-test-min}
\mathlibok
Let $\phi\colon\Ar\to\Ar$ be continuous at $t_0$ with $\phi'(t_0)=0$,
and suppose that $\phi'$ is defined on a neighborhood of $t_0$ and
differentiable at $t_0$ with $\phi''(t_0)=(\phi')'(t_0)>0$. Then
$\phi$ has a local minimum at $t_0$.
\end{lemma}

\begin{lemma}[Second-derivative test at a local maximum]
\label{lem:second-derivative-test-max}
\lean{MorganTianLib.deriv_deriv_nonpos_of_isLocalMax}
\leanok
Let $\phi\colon\Ar\to\Ar$ be continuous at $t_0$ and have a local
maximum at $t_0$. If $\phi$ is differentiable on a neighborhood of
$t_0$ and $\phi'$ is differentiable at $t_0$, then
$$\phi''(t_0)=(\phi')'(t_0)\le 0.$$
\end{lemma}

\begin{proof}
\leanok
\uses{lem:second-derivative-test-min}
First, $\phi'(t_0)=0$ by Fermat's lemma: for $t>t_0$ near $t_0$ the
difference quotients $\frac{\phi(t)-\phi(t_0)}{t-t_0}$ are $\le0$,
while for $t<t_0$ they are $\ge0$, so their common limit $\phi'(t_0)$
is $0$. Suppose for contradiction that $\phi''(t_0)>0$. By
Lemma~\ref{lem:second-derivative-test-min}, $\phi$ then has a local
minimum at $t_0$. Being simultaneously a local maximum and a local
minimum, $\phi$ is constant on a neighborhood of $t_0$; hence $\phi'$
vanishes identically near $t_0$, so $\phi''(t_0)=0$, contradicting
$\phi''(t_0)>0$.
\end{proof}

\begin{lemma}[Existence of integral curves]
\label{lem:integral-curve-exists}
\mathlibok
Let $M$ be a smooth manifold without boundary, let $X$ be a smooth
vector field on $M$, and let $q\in M$. Then there is a curve
$\gamma\colon\Ar\to M$ with $\gamma(0)=q$ which is an integral curve
of $X$ near $0$: for every $t$ in some neighborhood of $0$, $\gamma$
is differentiable at $t$ with $\gamma'(t)=X(\gamma(t))$.
\end{lemma}

\begin{lemma}[Vanishing differential at a local maximum]
\label{lem:differential-vanishes-at-max}
\lean{MorganTianLib.mfderiv_eq_zero_of_isLocalMax}
\leanok
Let $M$ be a smooth manifold without boundary and let $f\colon M\to\Ar$
be smooth. If $f$ has a local maximum at $q\in M$, then $df_q=0$.
\end{lemma}

\begin{proof}
\leanok
\uses{lem:integral-curve-exists}
Let $v\in T_qM$ and choose a smooth vector field $X$ on $M$ with
$X(q)=v$ (extend $v$ by a bump function supported in a chart at $q$).
Let $\gamma$ be an integral curve of $X$ with $\gamma(0)=q$
(Lemma~\ref{lem:integral-curve-exists}). The composite
$\phi=f\circ\gamma$ is differentiable near $0$ with
$\phi'(t)=df_{\gamma(t)}\!\left(X(\gamma(t))\right)$ by the chain
rule, and it has a local maximum at $0$ because $\gamma$ is continuous
with $\gamma(0)=q$. Fermat's lemma gives
$0=\phi'(0)=df_q(v)$. As $v\in T_qM$ was arbitrary, $df_q=0$.
\end{proof}

\begin{lemma}[Hessian is negative semi-definite at a local maximum]
\label{lem:hessian-nonpositive-at-max}
\lean{MorganTianLib.hessianAt_nonpos_of_isLocalMax}
\leanok
\uses{lem:hessian-symmetric}
Let $M$ be a smooth manifold without boundary equipped with an affine
connection $\nabla$, and let $f\colon M\to\Ar$ be smooth. If $f$ has a
local maximum at $q\in M$, then the Hessian of $f$ at $q$ is negative
semi-definite: for every $v\in T_qM$,
$$\Hess(f)_q(v,v)\le0.$$
\end{lemma}

\begin{proof}
\leanok
\uses{lem:differential-vanishes-at-max,lem:integral-curve-exists,lem:second-derivative-test-max,lem:hessian-symmetric}
Extend $v$ to a smooth vector field $X$ on $M$ with $X(q)=v$; that the
value $\Hess(f)(X,X)(q)$ depends only on $v=X(q)$, so that
$\Hess(f)_q$ is a well-defined bilinear form on $T_qM$, is the
tensoriality of the Hessian (Lemma~\ref{lem:hessian-symmetric}).
Recall
$$\Hess(f)(X,X)(q)=X(X(f))(q)-\left(\nabla_XX\right)(f)(q).$$
By Lemma~\ref{lem:differential-vanishes-at-max} we have $df_q=0$, so
the second term $\left(\nabla_XX\right)(f)(q)
=df_q\!\left((\nabla_XX)(q)\right)$ vanishes. For the first term, let
$\gamma$ be an integral curve of $X$ with $\gamma(0)=q$
(Lemma~\ref{lem:integral-curve-exists}) and set $\phi=f\circ\gamma$.
By the chain rule, $\phi'(t)=(Xf)(\gamma(t))$ for $t$ near $0$, and
differentiating once more, again by the chain rule applied to the
smooth function $Xf$ along $\gamma$,
$$\phi''(0)=\big(X(Xf)\big)(\gamma(0))=X(X(f))(q).$$
The function $\phi$ is continuous near $0$ and has a local maximum at
$0$ ($\gamma$ is continuous and $\gamma(0)=q$), so
Lemma~\ref{lem:second-derivative-test-max} gives $\phi''(0)\le0$.
Hence $\Hess(f)_q(v,v)=X(X(f))(q)=\phi''(0)\le0$.
\end{proof}

\begin{lemma}[Local smooth extension]
\label{lem:smooth-local-extension}
\lean{MorganTianLib.exists_contMDiff_eventuallyEq_of_contMDiffOn}
\leanok
Let $M$ be a smooth manifold, let $U\subset M$ be open, let $q\in U$,
and let $f\colon U\to\Ar$ be smooth. Then there is a smooth function
$f'\colon M\to\Ar$ agreeing with $f$ on a neighborhood of $q$.
\end{lemma}

\begin{proof}
\leanok
Choose a smooth bump function $\chi$ on $M$ at $q$ with compact
support contained in $U$ and with $\chi\equiv1$ on a neighborhood of
$q$ (transport a Euclidean bump function through a chart at $q$,
scaled so that its support lands in the chart image of $U$). Define
$f'=\chi f$ on $U$ and $f'=0$ on the complement of the support of
$\chi$. The two prescriptions agree on the overlap, are smooth on
their respective open domains, and the two domains cover $M$; hence
$f'$ is smooth on $M$, and $f'=f$ on the neighborhood of $q$ where
$\chi\equiv1$.
\end{proof}

\begin{lemma}[Laplacian is non-positive at an interior maximum]
\label{lem:laplacian-nonpositive-at-max}
\lean{MorganTianLib.laplacianAt_nonpos_of_isLocalMaxOn}
\leanok
\uses{lem:laplacian-local-formula}
Let $(M,g)$ be a Riemannian manifold, let $U\subset M$ be open, and
let $v\colon U\to \Ar$ be a smooth function. If $q\in U$ is a local
maximum of $v$, then $dv(q)=0$ and $\Delta v(q)\le 0$.
\end{lemma}

\begin{proof}
\leanok
\uses{lem:smooth-local-extension,lem:differential-vanishes-at-max,lem:hessian-nonpositive-at-max,lem:laplacian-local-formula}
Both conclusions involve only the restriction of $v$ to an arbitrarily
small neighborhood of $q$, so by Lemma~\ref{lem:smooth-local-extension}
we may replace $v$ by a smooth function on all of $M$ that agrees with
$v$ near $q$; it still has a local maximum at $q$.

The first assertion, $dv(q)=0$, is
Lemma~\ref{lem:differential-vanishes-at-max}.

For the second, recall that the Laplacian is the metric trace of the
Hessian, $\Delta v=g^{ij}\Hess(v)(\partial_i,\partial_j)$ in local
coordinates (the definition preceding
Lemma~\ref{lem:laplacian-local-formula}). Consequently, for any
orthonormal basis $e_1,\dots,e_n$ of $(T_qM,g_q)$,
$$\Delta v(q)=\sum_{i=1}^n\Hess(v)_q(e_i,e_i):$$
both sides compute the trace of the endomorphism of $T_qM$ obtained
from the bilinear form $\Hess(v)_q$ by raising an index with $g_q$
--- the left-hand side in the coordinate basis, the right-hand side in
the orthonormal basis --- and the trace is basis-independent. By
Lemma~\ref{lem:hessian-nonpositive-at-max} each diagonal term
satisfies $\Hess(v)_q(e_i,e_i)\le0$, whence $\Delta v(q)\le0$.
\end{proof}

\begin{lemma}[Strong maximum principle for smooth subharmonic functions]
\label{lem:hopf-strong-maximum}
\uses{lem:laplacian-christoffel-formula}
\lean{MorganTianLib.hopf_strong_maximum}
\leanok
(E.~Hopf) Let $(M,g)$ be a Riemannian manifold, let $U\subset M$ be
a connected open set, and let $h\colon U\to\Ar$ be a smooth
function with $\Delta h\ge 0$ at every point of $U$. If $h$ attains
its supremum $m:=\sup_Uh$ at some point of $U$, then $h$ is
identically equal to $m$.
\end{lemma}

\begin{proof}
\uses{lem:laplacian-christoffel-formula,lem:laplacian-nonpositive-at-max}
\leanok
Let $E:=\{x\in U\,|\,h(x)=m\}$. By hypothesis $E$ is non-empty, and
$E$ is closed in $U$ since $h$ is continuous. Since $U$ is
connected, it suffices to show that $E$ is open in $U$; then $E$,
being a non-empty subset of $U$ that is both open and closed in
$U$, equals $U$.

Fix $z\in E$. Choose a coordinate chart defined on an open set
$V\subset U$ containing $z$, with coordinate image
$\widehat{V}\subset\Ar^n$; we identify functions on $V$ with
functions on $\widehat{V}$ and write $\bar z\in\widehat{V}$ for the
coordinates of $z$. Throughout this proof $|\cdot|$ denotes the
Euclidean norm on the coordinates, $D(y,r)$ and $\overline{D}(y,r)$
the open and closed Euclidean balls. Choose $\rho>0$ so small that
$\overline{D}(\bar z,3\rho)\subset\widehat{V}$, and set
$K:=\overline{D}(\bar z,2\rho)$, a compact set. By the Christoffel
form of the coordinate Laplacian
(Lemma~\ref{lem:laplacian-christoffel-formula}), we have on
$\widehat{V}$, for every smooth $v$,
$$\Delta v=g^{ij}\partial_i\partial_jv+b^j\partial_jv,
\qquad b^j:=-g^{ik}\Gamma^j_{ik},$$
with $(g^{ij})$ symmetric positive definite at each point and with
smooth coefficient functions $b^j$, both as recorded in that lemma.
The function $(x,\xi)\mapsto g^{ij}(x)\xi_i\xi_j$ is continuous and
strictly positive on the compact set $K\times S^{n-1}$, so it has a
positive minimum there; likewise $\sum_ig^{ii}$ and the Euclidean
norm $|b|=\left(\sum_j(b^j)^2\right)^{1/2}$ are continuous on $K$
and hence bounded there. Fix constants $\lambda>0$ and
$T,C<\infty$ with
$$g^{ij}(x)\xi_i\xi_j\ge\lambda|\xi|^2,\qquad
\sum_ig^{ii}(x)\le T,\qquad |b(x)|\le C \qquad\text{for all } x\in
K,\ \xi\in\Ar^n.$$

We claim that $h\equiv m$ on $D(\bar z,\rho)$; since the preimage
of $D(\bar z,\rho)$ under the coordinates is an open neighborhood
of $z$ in $U$, this proves that $E$ is open and finishes the
argument. Suppose, for contradiction, that there is a point
$x_0\in D(\bar z,\rho)$ with $h(x_0)<m$.

{\bf The touching ball.} Let $F:=\{x\in K\,|\,h(x)=m\}$. Then $F$
is compact (a closed subset of the compact set $K$) and non-empty
($\bar z\in F$). Since $h(x_0)<m$ we have $x_0\notin F$, so
$$R:=\dist(x_0,F)>0,$$
and $R\le|x_0-\bar z|<\rho$ because $\bar z\in F$. Every
$x\in D(x_0,R)$ satisfies $|x-\bar z|\le|x-x_0|+|x_0-\bar
z|<R+\rho\le2\rho$, so $D(x_0,R)\subset K$; moreover $x\notin F$
(its distance to $x_0$ is less than $R$), and $h\le m$ everywhere
on $U$ (as $m=\sup_Uh$), so
$$h<m\quad\text{on }D(x_0,R).$$
The infimum defining $R$ is attained: choosing $y_k\in F$ with
$|y_k-x_0|\to R$ and passing to a convergent subsequence in the
compact set $F$ yields $y\in F$ with $|y-x_0|=R$. The point $y$
lies in the open coordinate domain $V\subset U$, and $h\le m=h(y)$
on all of $U$; so $y$ is a local (indeed global)
maximum of $h$, and the first assertion of
Lemma~\ref{lem:laplacian-nonpositive-at-max} (applied to $h$ on the
coordinate domain) gives $dh(y)=0$; in the coordinates, all first
partial derivatives of $h$ vanish at $y$:
\begin{equation}
\partial_jh(y)=0,\qquad j=1,\dots,n.\label{eqngradvanish}
\end{equation}

{\bf The barrier.} Consider the closed annulus
$$A:=\{x\,|\,R/2\le|x-x_0|\le R\},$$
which is contained in $K$ (each $x\in A$ has $|x-\bar
z|\le R+\rho<2\rho$). For $\alpha>0$ define
$$w_\alpha(x):=e^{-\alpha|x-x_0|^2}-e^{-\alpha R^2}.$$
Then $w_\alpha$ is smooth, $w_\alpha>0$ on $\{|x-x_0|<R\}$,
$w_\alpha=0$ on $\{|x-x_0|=R\}$, and $0<w_\alpha<1$ on
$\{|x-x_0|<R\}$. Direct computation gives
$$\partial_jw_\alpha=-2\alpha(x-x_0)_je^{-\alpha|x-x_0|^2},\qquad
\partial_i\partial_jw_\alpha=\left(4\alpha^2(x-x_0)_i(x-x_0)_j
-2\alpha\delta_{ij}\right)e^{-\alpha|x-x_0|^2},$$
so that on $A$
\begin{align*}
\Delta w_\alpha&=e^{-\alpha|x-x_0|^2}\left(4\alpha^2\,
g^{ij}(x-x_0)_i(x-x_0)_j-2\alpha\sum_ig^{ii}
-2\alpha\,b^j(x-x_0)_j\right)\\
&\ge e^{-\alpha|x-x_0|^2}\left(4\alpha^2\lambda|x-x_0|^2
-2\alpha T-2\alpha C|x-x_0|\right)\\
&\ge e^{-\alpha|x-x_0|^2}\,\alpha\left(\alpha\lambda
R^2-2T-2CR\right),
\end{align*}
using $R/2\le|x-x_0|\le R$ in the last step. Fix
$$\alpha:=\frac{2T+2CR}{\lambda R^2}+1,$$
so that $\Delta w_\alpha>0$ at every point of $A$. Note also that
for every unit vector $\nu$ the radial derivative on the outer
sphere is
\begin{equation}
\frac{d}{dt}\Big|_{t=R}w_\alpha(x_0+t\nu)=-2\alpha
Re^{-\alpha R^2}<0.\label{eqnradial}
\end{equation}

{\bf The contradiction.} The inner sphere $\{|x-x_0|=R/2\}$ is a
compact subset of $D(x_0,R)$, on which $h<m$; hence
$$\delta:=m-\max_{|x-x_0|=R/2}h>0.$$
Set $\eps:=\delta/2$ and
$$v:=h+\eps\, w_\alpha\qquad\text{on a neighborhood of } A.$$
We record three properties.
\begin{enumerate}
\item On the inner sphere $\{|x-x_0|=R/2\}$ we have $v\le
(m-\delta)+\eps\cdot 1=m-\delta/2<m$.
\item $v(y)=h(y)+\eps w_\alpha(y)=m+0=m$, since $|y-x_0|=R$.
\item At every point of $A$ we have $\Delta v=\Delta
h+\eps\Delta w_\alpha>0$, since $\Delta h\ge0$ on $U$ and
$\Delta w_\alpha>0$ on $A$.
\end{enumerate}
The continuous function $v$ attains its maximum over the compact
set $A$ at some point $q$, and by (2) that maximum is $\ge m$. If
$q$ were an interior point of $A$, then $q$ would be a local
maximum of the smooth function $v$ on the open subset of $M$
corresponding to the interior of $A$, so
Lemma~\ref{lem:laplacian-nonpositive-at-max} would give $\Delta
v(q)\le0$, contradicting (3). Hence $q$ lies on the boundary of
$A$, which is the union of the two spheres. By (1) the value of
$v$ on the inner sphere is $<m\le v(q)$, so $|q-x_0|=R$. Then
$w_\alpha(q)=0$, so $h(q)=v(q)\ge m$; since $h\le m$ this forces
$h(q)=m$ and
$$\max_Av=v(q)=m.$$
In particular $q$, like $y$, is a point of the coordinate domain at
which $h$ attains its global maximum $m$, so
Equation~(\ref{eqngradvanish}) holds at $q$ as well:
$\partial_jh(q)=0$ for all $j$.

Let $\nu:=(q-x_0)/R$, a Euclidean unit vector, and consider
$$\phi(t):=v(x_0+t\nu),\qquad t\in[R/2,R].$$
For each such $t$ the point $x_0+t\nu$ lies in $A$, so $\phi(t)\le
m=\phi(R)$; dividing $\phi(R)-\phi(t)\ge0$ by $R-t>0$ and letting
$t\uparrow R$ gives $\phi'(R)\ge0$. On the other hand, by the
chain rule, the vanishing of the first partials of $h$ at $q$, and
Equation~(\ref{eqnradial}),
$$\phi'(R)=\sum_j\partial_jh(q)\,\nu_j+\eps\,
\frac{d}{dt}\Big|_{t=R}w_\alpha(x_0+t\nu)
=0-2\eps\alpha Re^{-\alpha R^2}<0,$$
a contradiction. Therefore no point $x_0\in D(\bar z,\rho)$ with
$h(x_0)<m$ exists, i.e., $h\equiv m$ on $D(\bar z,\rho)$, and $E$
is open. As explained at the outset, $E=U$ by connectedness, which
is the assertion of the lemma.
\end{proof}

\begin{theorem}[Maximum principle]
\label{thm:maximum-principle}
\uses{rem:laplacian-weak-sense}
\source{morgan-tian}
Let $f$ be a real-valued continuous function on a connected
Riemannian manifold with $\Delta f\ge 0$ in the weak sense. Then $f$
is locally constant near any local maximum. In particular, if $f$
achieves its maximum then it is a constant. (Here, as in
Remark~\ref{rem:laplacian-weak-sense}, the hypothesis $\Delta f\ge0$
in the weak sense on the manifold $M$ means that for every
non-negative compactly supported $C^\infty$-function $\varphi$ on
$M$ we have $\int_Mf\,\Delta\varphi\,d\vol\ge0$; completeness of
the manifold is not assumed.)
\end{theorem}

\begin{proof}
\uses{claim:dirichlet-problem-small-balls,claim:weak-subsolution-comparison,lem:green-identity-compact-support,lem:hopf-strong-maximum,def:exponential-map}
Let $(M,g)$ be the connected Riemannian manifold in question.

{\bf Step 0: choice of a ball around a local maximum.} Let $p$ be a
local maximum of $f$: there is an open set $W\ni p$ with $f\le
f(p)=:m$ on $W$. Let $r_D=r_D(p)>0$ be as in
Claim~\ref{claim:dirichlet-problem-small-balls}. Since $\exp_p$ is
continuous and $\exp_p(0)=p$
(Definition~\ref{def:exponential-map}), the preimage
$\exp_p^{-1}(W)$ is an open neighborhood of $0$ in $T_pM$ and hence
contains a ball $B(0,s_0)$ with $0<s_0\le r_D$. Fix $t$ with
$0<t<s_0$. Then by Item (1) of
Claim~\ref{claim:dirichlet-problem-small-balls},
$$\overline{B}(p,t)=\exp_p\left(\overline{B}(0,t)\right)\subset
\exp_p\left(B(0,s_0)\right)\subset W,$$
so that
\begin{equation}
f\le m\ \text{ on }\overline{B}(p,t),\qquad f(p)=m.
\label{eqnlocmax}
\end{equation}

{\bf Step 1: harmonic replacement and comparison.} Fix $s\in(0,t]$.
By Item (2) of Claim~\ref{claim:dirichlet-problem-small-balls}
applied with boundary values $\varphi:=f|_{S(p,s)}$, there is a
function $h_s$, continuous on $\overline{B}(p,s)$ and smooth on
$B(p,s)$, with $\Delta h_s=0$ in $B(p,s)$ and $h_s=f$ on $S(p,s)$.

We verify that the open set $\Omega:=B(p,s)$ satisfies the
hypotheses of Claim~\ref{claim:weak-subsolution-comparison}: by
Item (1) of Claim~\ref{claim:dirichlet-problem-small-balls} it is
non-empty and {\it connected}, its closure $\overline{B}(p,s)$ is
compact, and its topological boundary $S(p,s)$ is non-empty.

First we apply Claim~\ref{claim:weak-subsolution-comparison} to
$w:=f-h_s$ on $\overline{B}(p,s)$, which is continuous and
vanishes on $S(p,s)$. To see that $\Delta w\ge0$ in the weak sense
on $B(p,s)$, let $\psi\ge0$ be a compactly supported
$C^\infty$-function on $B(p,s)$. Extending $\psi$ by zero on
$M\setminus B(p,s)$ produces a non-negative compactly supported
$C^\infty$-function on $M$ (it is smooth at each point: points of
$B(p,s)$ have the neighborhood $B(p,s)$ on which it agrees with
$\psi$, and points outside the compact set $\supp(\psi)$ have
the neighborhood $M\setminus\supp(\psi)$ on which it
vanishes). The hypothesis on $f$ then gives $\int_Mf\Delta\psi\,
d\vol\ge0$; since $\Delta\psi$ vanishes outside $\supp
(\psi)\subset B(p,s)$, this integral equals $\int_{B(p,s)}f\Delta
\psi\,d\vol$. On the other hand,
Lemma~\ref{lem:green-identity-compact-support} applied to the
smooth function $u=h_s$ on $\Omega=B(p,s)$ and to $\psi$ gives
$$\int_{B(p,s)}h_s\,\Delta\psi\,d\vol=\int_{B(p,s)}\psi\,\Delta
h_s\,d\vol=0.$$
Subtracting,
$$\int_{B(p,s)}w\,\Delta\psi\,d\vol=\int_{B(p,s)}f\,\Delta\psi\,
d\vol\ \ge\ 0,$$
as required. Claim~\ref{claim:weak-subsolution-comparison} now
yields
$$\max_{\overline{B}(p,s)}\left(f-h_s\right)
=\max_{S(p,s)}\left(f-h_s\right)=0,\qquad\text{i.e.,}\qquad f\le
h_s\ \text{ on }\overline{B}(p,s).$$

Next we apply Claim~\ref{claim:weak-subsolution-comparison} to
$h_s$ itself, on the same set $\Omega=B(p,s)$ (whose hypotheses
were verified above). By
Lemma~\ref{lem:green-identity-compact-support} (with $u=h_s$), for
every compactly supported $C^\infty$-function $\psi$ on $B(p,s)$ we
have $\int h_s\Delta\psi\,d\vol=\int\psi\Delta h_s\,d\vol=0$; in
particular $\Delta h_s\ge0$ in the weak sense on $B(p,s)$. Hence
$$\max_{\overline{B}(p,s)}h_s=\max_{S(p,s)}h_s=\max_{S(p,s)}f\le
m,$$
the last inequality by (\ref{eqnlocmax}), since $S(p,s)\subset
\overline{B}(p,t)$. On the other hand $h_s(p)\ge f(p)=m$ by the
comparison just proved. Combining the last two displays,
$$h_s(p)=m=\max_{\overline{B}(p,s)}h_s,$$
so the function $h_s$ attains its maximum over
$\overline{B}(p,s)$ at the interior point $p\in B(p,s)$.

{\bf Step 2: the replacement is constant.} The open set $B(p,s)$ is
connected by Item (1) of
Claim~\ref{claim:dirichlet-problem-small-balls}, the function
$h_s$ is smooth on it with $\Delta h_s=0\ge0$ pointwise, and
$$\sup_{B(p,s)}h_s\le\max_{\overline{B}(p,s)}h_s=m=h_s(p)$$
with $p\in B(p,s)$; so $h_s$ attains its supremum over $B(p,s)$ at
$p$. By Lemma~\ref{lem:hopf-strong-maximum}, $h_s\equiv m$ on
$B(p,s)$, and hence, by continuity, on $\overline{B}(p,s)$. In
particular, on the boundary sphere,
$$f=h_s\equiv m\qquad\text{on } S(p,s).$$

{\bf Step 3: local constancy.} The conclusion of Step 2 holds for
every $s\in(0,t]$. Now let $x\in B(p,t)$ be arbitrary. If $x=p$
then $f(x)=m$; otherwise $s:=d(p,x)$ satisfies $0<s<t$, and $x\in
S(p,s)$ by the definition of the metric sphere, so $f(x)=m$ by
Step 2. Hence $f\equiv m$ on the open neighborhood $B(p,t)$ of
$p$: the function $f$ is constant near any local maximum.

{\bf Step 4: global constancy.} Suppose that $f$ achieves its
maximum: there is $x_0\in M$ with $f(x_0)=\sup_Mf=:m_1$. Let
$E:=f^{-1}(m_1)$. Then $E$ is non-empty ($x_0\in E$) and closed in
$M$ (as $f$ is continuous). Moreover every $x\in E$ is a local
maximum of $f$, since $f\le m_1=f(x)$ on all of $M$; so by Steps
0--3 (applied at $p=x$, with $W=M$) the function $f$ is identically
$m_1$ on some metric ball $B(x,t_x)$ with $t_x>0$, whence
$B(x,t_x)\subset E$. Thus $E$ is also open. Since $M$ is connected
and $E$ is a non-empty, open and closed subset, $E=M$, i.e., $f$ is
constant. This completes the proof.
\end{proof}

\begin{theorem}[Regularity of weakly harmonic functions]
\label{thm:weak-harmonic-regularity}
\uses{rem:laplacian-weak-sense}
Let $(M,g)$ be a Riemannian manifold and let $f\colon M\to\Ar$ be a
continuous function that is harmonic in the weak sense, meaning
(cf.\ Remark~\ref{rem:laplacian-weak-sense}) that for every
compactly supported $C^\infty$-function $\varphi$ on $M$ we have
$$\int_Mf\,\Delta\varphi\,d\vol=0.$$
Then $f$ is smooth, and $\Delta f=0$ in the classical sense.
\end{theorem}

\begin{proof}
\uses{claim:dirichlet-problem-small-balls,claim:weak-subsolution-comparison,lem:green-identity-compact-support}
Smoothness is a local property, so it suffices to show that $f$ is
smooth and classically harmonic on a neighborhood of each point
$p\in M$. Let $r_D=r_D(p)>0$ be as in
Claim~\ref{claim:dirichlet-problem-small-balls} and fix
$s\in(0,r_D]$. By Item (1) of that claim the ball $B(p,s)$ is
non-empty and connected, with compact closure $\overline{B}(p,s)$
and non-empty topological boundary $S(p,s)$; so $B(p,s)$ satisfies
the hypotheses of Claim~\ref{claim:weak-subsolution-comparison}.
By Item (2) of Claim~\ref{claim:dirichlet-problem-small-balls}
applied with boundary values $f|_{S(p,s)}$ there is a function
$u$, continuous on $\overline{B}(p,s)$ and smooth on $B(p,s)$,
with $\Delta u=0$ in $B(p,s)$ and $u=f$ on $S(p,s)$.

We claim that $f=u$ on $\overline{B}(p,s)$. Let $\psi\ge0$ be a
compactly supported $C^\infty$-function on $B(p,s)$. Extending
$\psi$ by zero on $M\setminus B(p,s)$ produces a non-negative
compactly supported $C^\infty$-function on $M$: it is smooth at
each point, since points of $B(p,s)$ have the neighborhood
$B(p,s)$ on which it agrees with $\psi$, while points outside the
compact set $\supp(\psi)$ have the neighborhood
$M\setminus\supp(\psi)$ on which it vanishes. The hypothesis on
$f$ therefore gives
$$\int_{B(p,s)}f\,\Delta\psi\,d\vol=\int_Mf\,\Delta\psi\,d\vol=0,$$
the first equality because $\Delta\psi$ vanishes outside
$\supp(\psi)\subset B(p,s)$; and
Lemma~\ref{lem:green-identity-compact-support} applied to the
smooth function $u$ on $B(p,s)$ gives
$$\int_{B(p,s)}u\,\Delta\psi\,d\vol
=\int_{B(p,s)}\psi\,\Delta u\,d\vol=0.$$
Hence the two continuous functions $w=f-u$ and $-w=u-f$ on
$\overline{B}(p,s)$ each vanish on $S(p,s)$ and each satisfy
$\int_{B(p,s)}(\pm w)\,\Delta\psi\,d\vol=0\ge0$ for every
non-negative compactly supported $C^\infty$-function $\psi$ on
$B(p,s)$, i.e., $\Delta(\pm w)\ge0$ in the weak sense there.
Claim~\ref{claim:weak-subsolution-comparison} applied to $w$ and
to $-w$ yields
$$\max_{\overline{B}(p,s)}(\pm w)=\max_{S(p,s)}(\pm w)=0,$$
so $w\equiv0$, i.e., $f=u$ on $\overline{B}(p,s)$. In particular
$f$ is smooth on $B(p,s)$ and $\Delta f=\Delta u=0$ there in the
classical sense. Since $p\in M$ was arbitrary, $f$ is smooth and
harmonic on all of $M$.
\end{proof}

\begin{remark}[Provenance of the regularity theorem]
\label{rem:weak-harmonic-regularity-citation}
\uses{thm:weak-harmonic-regularity,claim:dirichlet-problem-small-balls,claim:weak-subsolution-comparison,rem:laplacian-weak-sense,lem:laplacian-local-formula}
Theorem~\ref{thm:weak-harmonic-regularity} is Weyl's lemma for the
Laplace--Beltrami operator: in local coordinates $\Delta$ is a
second-order elliptic operator with smooth coefficients
(Lemma~\ref{lem:laplacian-local-formula}), and continuous
distributional solutions of such equations are smooth. The proof
given above derives it from the two citation-backed claims of this
section, which are therefore the only analytic black boxes that
this chapter consumes. For Riemannian manifolds the theorem is
stated as the regularity theorem for harmonic functions in
\cite{Petersen}, p.~279ff.\ (in the third edition: Theorem 7.1.8,
pp.~283--284), and proved there by harmonic replacement; note,
however, that Petersen's development of the maximum principle is
in the barrier sense, so his argument does not transpose verbatim
to the distributional setting of
Remark~\ref{rem:laplacian-weak-sense}. For the underlying PDE
facts see Gilbarg--Trudinger, {\it Elliptic Partial Differential
Equations of Second Order}, Chapters 6 and 8, and H.~Weyl, {\it
The method of orthogonal projection in potential theory}, Duke
Math.~J.\ {\bf 7} (1940), 411--444.
\end{remark}

We shall also need the Bochner formula for functions and the fact
that a complete manifold splits isometrically along a smooth
function with vanishing Hessian and unit gradient, together with
the elementary apparatus of covariant differentiation along smooth
maps that their proofs use. We establish these next, beginning with
the first-order calculus of the function $|\nabla f|^2$.

\begin{lemma}[First derivatives of the square norm of a gradient]
\label{lem:gradient-norm-sq-gradient}
\uses{thm:levi-civita-connection,lem:hessian-symmetric}
\lean{MorganTianLib.metricNormSq,MorganTianLib.contMDiff_metricNormSq,MorganTianLib.dir_metricNormSq,MorganTianLib.dir_metricNormSq_gradientField,MorganTianLib.gradientField_metricNormSq_gradientField,MorganTianLib.hessian_metricNormSq_gradientField}
\leanok
Let $(M,g)$ be a Riemannian manifold with Levi-Civita connection
$\nabla$, let $f$ be a smooth function on $M$, and write
$G=(\nabla f)^*$ for its gradient vector field. Then
$|\nabla f|^2=\langle G,G\rangle$ is a smooth function on $M$ and,
for all smooth vector fields $X$ and $Y$ on $M$:
\begin{enumerate}
\item $Y\left(|\nabla f|^2\right)=2\left\langle\nabla_YG,G\right\rangle$;
\item $\left(\nabla|\nabla f|^2\right)^*=2\,\nabla_GG$;
\item $\Hess\left(|\nabla f|^2\right)(X,Y)
=2\left\langle\nabla_X\left(\nabla_GG\right),Y\right\rangle$.
\end{enumerate}
\end{lemma}

\begin{proof}
\uses{thm:levi-civita-connection,lem:hessian-symmetric}
\leanok
Smoothness of $|\nabla f|^2$ holds because the metric pairing of two
smooth vector fields is a smooth function (in a chart it is a
composition of the smooth metric coefficients with the smooth
component functions of the fields), and $G$ is smooth.

(1) Compatibility of $\nabla$ with the metric
(Theorem~\ref{thm:levi-civita-connection}) gives
$$Y\langle G,G\rangle
=\langle\nabla_YG,G\rangle+\langle G,\nabla_YG\rangle
=2\langle\nabla_YG,G\rangle.$$

(2) By the gradient formula for the Hessian
(Lemma~\ref{lem:hessian-symmetric}),
$\langle\nabla_YG,G\rangle=\Hess(f)(Y,G)$; by the symmetry of the
Hessian this equals $\Hess(f)(G,Y)$, which reads back through the
gradient formula as $\langle\nabla_GG,Y\rangle$. Combining with (1),
$$Y\left(|\nabla f|^2\right)
=\left\langle 2\,\nabla_GG,\,Y\right\rangle
\qquad\text{for every smooth vector field }Y.$$
Since the gradient of $|\nabla f|^2$ is by definition the vector
field pairing to $Y\left(|\nabla f|^2\right)$ against every $Y$, and
tangent vectors with equal pairings against all tangent vectors at a
point coincide (non-degeneracy of $g$; every tangent vector extends
to a smooth vector field), this forces
$\left(\nabla|\nabla f|^2\right)^*=2\,\nabla_GG$.

(3) Apply the gradient formula for the Hessian
(Lemma~\ref{lem:hessian-symmetric}) to the smooth function
$|\nabla f|^2$ and use (2):
$$\Hess\left(|\nabla f|^2\right)(X,Y)
=\left\langle\nabla_X\left(\left(\nabla|\nabla f|^2\right)^*\right),
Y\right\rangle
=\left\langle\nabla_X\left(2\,\nabla_GG\right),Y\right\rangle
=2\left\langle\nabla_X\left(\nabla_GG\right),Y\right\rangle,$$
the constant $2$ passing through $\nabla_X$ by the Leibniz rule
(the derivative of a constant function vanishes).
\end{proof}

\begin{definition}[Square norm of the Hessian]
\label{def:hessian-norm-sq}
\uses{lem:hessian-symmetric}
\lean{MorganTianLib.hessianNormSqAt}
\leanok
Let $f$ be a smooth function on a Riemannian manifold $(M,g)$ and
let $p\in M$. The \textit{square norm of the Hessian} of $f$ at $p$
is
$$|\Hess(f)|^2(p)=\sum_{i,j=1}^n\Hess(f)_p(e_i,e_j)^2,$$
where $\{e_1,\ldots,e_n\}$ is an orthonormal basis of $(T_pM,g_p)$;
by Lemma~\ref{lem:hessian-norm-sq-basic} the value does not depend
on the choice of orthonormal basis. In local coordinates this is
$g^{ik}g^{jl}\Hess(f)_{ij}\Hess(f)_{kl}$, the square norm appearing
in the Bochner formula.
\end{definition}

\begin{lemma}[Basic properties of the Hessian square norm]
\label{lem:hessian-norm-sq-basic}
\uses{def:hessian-norm-sq,lem:hessian-symmetric}
\lean{MorganTianLib.hessianNormSqAt_eq_sum,MorganTianLib.hessianNormSqAt_nonneg,MorganTianLib.hessianNormSqAt_eq_zero_iff,MorganTianLib.sq_laplacianAt_le_finrank_mul_hessianNormSqAt}
\leanok
Let $f$ be smooth on $(M,g)$, $p\in M$, and $n=\dim M$. Then:
\begin{enumerate}
\item every orthonormal basis of $(T_pM,g_p)$ computes the same
value $\sum_{i,j}\Hess(f)_p(e_i,e_j)^2$;
\item $|\Hess(f)|^2(p)\ge0$, with equality if and only if
$\Hess(f)_p=0$ as a bilinear form on $T_pM$;
\item $\left(\triangle f(p)\right)^2\le n\,|\Hess(f)|^2(p)$.
\end{enumerate}
\end{lemma}

\begin{proof}
\uses{lem:hessian-symmetric}
\leanok
(1) Let $\{e_i\}$ and $\{f_a\}$ be orthonormal bases of
$(T_pM,g_p)$ and write $B=\Hess(f)_p$. Expanding
$e_i=\sum_a\langle e_i,f_a\rangle f_a$ in both slots by bilinearity,
$$\sum_{i,j}B(e_i,e_j)^2
=\sum_{i,j}\sum_{a,b,a',b'}\langle e_i,f_a\rangle\langle
e_i,f_{a'}\rangle\langle e_j,f_b\rangle\langle e_j,f_{b'}\rangle
B(f_a,f_b)B(f_{a'},f_{b'}).$$
Summing first over $i$ and $j$ and using Parseval's identity
$\sum_i\langle e_i,f_a\rangle\langle e_i,f_{a'}\rangle
=\langle f_a,f_{a'}\rangle=\delta_{aa'}$ (and likewise for $j$)
collapses the right-hand side to $\sum_{a,b}B(f_a,f_b)^2$.

(2) The sum of squares is non-negative, and vanishes exactly when
$B(e_i,e_j)=0$ for all $i,j$, which by bilinearity and expansion of
arbitrary vectors in the basis $\{e_i\}$ is equivalent to $B=0$.

(3) By definition $\triangle f(p)=\sum_iB(e_i,e_i)$, so the
Cauchy--Schwarz inequality gives
$\left(\triangle f(p)\right)^2\le n\sum_iB(e_i,e_i)^2$, and the
diagonal sum of squares is at most the full sum of squares
$|\Hess(f)|^2(p)$.
\end{proof}

\begin{lemma}[Hessian square norm as a trace along the gradient]
\label{lem:hessian-riesz-trace}
\uses{def:hessian-norm-sq,lem:hessian-symmetric}
\lean{MorganTianLib.metricInner_cov_gradientField_eq_hessianAt,MorganTianLib.sum_metricInner_cov_cov_gradientField_eq_hessianNormSqAt}
\leanok
Let $(M,g)$ be Riemannian with Levi-Civita connection $\nabla$, let
$f$ be smooth with gradient field $G=(\nabla f)^*$, let $p\in M$,
and let $\{e_1,\ldots,e_n\}$ be an orthonormal basis of $T_pM$ with
smooth extensions $E_i$. Then, for every smooth vector field $X$ and
every $w\in T_pM$,
$$\left\langle\left(\nabla_XG\right)(p),w\right\rangle
=\Hess(f)_p\left(X(p),w\right),$$
and consequently
$$\sum_{i=1}^n\left\langle\nabla_{\nabla_{E_i}G}\,G,\,E_i
\right\rangle(p)=|\Hess(f)|^2(p).$$
\end{lemma}

\begin{proof}
\uses{lem:hessian-symmetric}
\leanok
The first identity is the pointwise form of the gradient formula for
the Hessian (Lemma~\ref{lem:hessian-symmetric}): pairing
$\nabla_XG$ against any extension of $w$ computes
$\Hess(f)(X,\cdot)$, and both sides depend on the second argument
only through its value $w$ at $p$. For the second, set
$w_i=\left(\nabla_{E_i}G\right)(p)$; by the first identity (applied
twice) $w_i$ is the Riesz vector of $\Hess(f)_p(e_i,\cdot)$, i.e.
$\langle w_i,e_j\rangle=\Hess(f)_p(e_i,e_j)$ for all $j$, and each
summand equals $\Hess(f)_p(w_i,e_i)$. Expanding
$w_i=\sum_j\langle w_i,e_j\rangle e_j$ and using linearity of the
Hessian in its first slot,
$$\Hess(f)_p(w_i,e_i)
=\sum_j\Hess(f)_p(e_i,e_j)\,\Hess(f)_p(e_j,e_i)
=\sum_j\Hess(f)_p(e_i,e_j)^2,$$
the last step by symmetry of the Hessian
(Lemma~\ref{lem:hessian-symmetric}). Summing over $i$ gives
$|\Hess(f)|^2(p)$ by Definition~\ref{def:hessian-norm-sq}.
\end{proof}

\begin{lemma}[Smooth local orthonormal frames]
\label{lem:local-orthonormal-frame}
\uses{def:riemannian-metric}
\lean{MorganTianLib.chartFrameNorm_section_contMDiffOn,MorganTianLib.exists_orthonormalFrame,MorganTianLib.orthonormalBasisOfMetricInner}
\leanok
Let $(M,g)$ be a Riemannian manifold and let $p\in M$. Then there are
smooth vector fields $F_1,\ldots,F_n$ on $M$ ($n=\dim M$) and a
neighborhood $U$ of $p$ such that for every $q\in U$ the vectors
$F_1(q),\ldots,F_n(q)$ form an orthonormal basis of $(T_qM,g_q)$.
\end{lemma}

\begin{proof}
\uses{def:riemannian-metric}
\leanok
Fix a coordinate chart around $p$ and let
$\partial_1,\ldots,\partial_n$ be its coordinate frame, a smooth frame
of $TM$ over the chart domain. Apply the Gram--Schmidt process
fiberwise with respect to $g$: set recursively
$$w_i=\partial_i-\sum_{j<i}\langle\partial_i,F_j\rangle F_j,
\qquad F_i=\frac{w_i}{\sqrt{\langle w_i,w_i\rangle}}.$$
At each point of the chart domain the vectors
$\partial_1,\ldots,\partial_n$ are linearly independent, so by
induction each $w_i$ is nonzero there: were $w_i=0$, then $\partial_i$
would lie in the span of $F_1,\ldots,F_{i-1}$, which by construction
equals the span of $\partial_1,\ldots,\partial_{i-1}$, contradicting
linear independence. Hence $\langle w_i,w_i\rangle>0$ on the chart
domain, and by induction each $F_i$ is smooth there: the coefficients
$\langle\partial_i,F_j\rangle$ are metric pairings of smooth fields,
hence smooth; the sum defining $w_i$ is a linear combination of smooth
fields with smooth coefficients; and the normalizing factor
$\left(\sqrt{\langle w_i,w_i\rangle}\right)^{-1}$ is a smooth function
of the positive smooth function $\langle w_i,w_i\rangle$. The
resulting frame is $g$-orthonormal at every point of the chart domain
by the usual Gram--Schmidt computation. Finally, multiply each $F_i$
by a smooth bump function equal to $1$ near $p$ and supported in the
chart domain, extending it by zero to a globally defined smooth vector
field; this does not change the frame on the neighborhood of $p$ where
the bump function equals $1$.
\end{proof}

\begin{lemma}[Derivative of the Laplacian as a trace of the second covariant derivative]
\label{lem:laplacian-trace-commutation}
\uses{def:second-covariant-derivative-vector-field,lem:hessian-symmetric}
\lean{MorganTianLib.laplacianAt_eventuallyEq_sum_frame,MorganTianLib.contMDiff_laplacianAt,MorganTianLib.metricInner_secondCov_middle_congr,MorganTianLib.sum_metricInner_secondCov_frame_eq_dir_laplacianAt,MorganTianLib.sum_metricInner_secondCov_gradientField_eq_dir_laplacianAt}
\leanok
Let $f$ be a smooth function on a Riemannian manifold $(M,g)$ with
Levi-Civita connection $\nabla$, let $G=(\nabla f)^*$ be its gradient
field, let $p\in M$, and let $\{e_1,\ldots,e_n\}$ be an orthonormal
basis of $T_pM$, each $e_i$ extended to a smooth vector field $E_i$.
Then $\triangle f$ is a smooth function on $M$ and
$$\sum_{i=1}^n\left\langle\nabla^2G(G,E_i),E_i\right\rangle(p)
=G(\triangle f)(p),$$
where $\nabla^2$ is the second covariant derivative of
Definition~\ref{def:second-covariant-derivative-vector-field}. That
is: the derivative along $G$ of the metric trace $\triangle f$ of
$\nabla G$ is the corresponding metric trace of the second covariant
derivative of $G$ with outer derivative along $G$.
\end{lemma}

\begin{proof}
\uses{lem:local-orthonormal-frame,lem:second-covariant-derivative-tensorial,lem:hessian-symmetric,lem:hessian-riesz-trace,thm:levi-civita-connection}
\leanok
By Lemma~\ref{lem:second-covariant-derivative-tensorial} each summand
$\left\langle\nabla^2G(G,E_i),E_i\right\rangle(p)$ depends on $E_i$
only through $e_i=E_i(p)$, and as a function of the pair of slots
occupied by $e_i$ it is a bilinear form on $T_pM$; the diagonal sum of
a bilinear form over an orthonormal basis is independent of the chosen
orthonormal basis. We may therefore compute the sum with any
orthonormal basis of $T_pM$ and any smooth extensions. Choose, by
Lemma~\ref{lem:local-orthonormal-frame}, a smooth local orthonormal
frame $F_1,\ldots,F_n$ at $p$ and take $e_i=F_i(p)$, $E_i=F_i$.

Since $\{F_i(q)\}$ is an orthonormal basis of $T_qM$ for every $q$
near $p$, the trace formula for the Laplacian
(Equation~(\ref{laplacformula})) together with the gradient formula
for the Hessian (Lemma~\ref{lem:hessian-symmetric}) gives, on a
neighborhood of $p$,
\begin{equation}\label{eq:laplacian-frame-sum}
\triangle f=\sum_i\Hess(f)(F_i,F_i)
=\sum_i\left\langle\nabla_{F_i}G,F_i\right\rangle .
\end{equation}
The right-hand side is a finite sum of metric pairings of smooth
vector fields, hence smooth; as $p$ was arbitrary, $\triangle f$ is a
smooth function on $M$. Differentiating
Equation~(\ref{eq:laplacian-frame-sum}) along $G$ at $p$ and using the
compatibility of $\nabla$ with the metric
(Theorem~\ref{thm:levi-civita-connection}),
$$G(\triangle f)(p)
=\sum_i\left\langle\nabla_G\nabla_{F_i}G,F_i\right\rangle(p)
+\sum_i\left\langle\nabla_{F_i}G,\nabla_GF_i\right\rangle(p).$$
On the other hand, by the definition of the second covariant
derivative,
$$\sum_i\left\langle\nabla^2G(G,F_i),F_i\right\rangle(p)
=\sum_i\left\langle\nabla_G\nabla_{F_i}G,F_i\right\rangle(p)
-\sum_i\left\langle\nabla_{\nabla_GF_i}\,G,F_i\right\rangle(p),$$
so it suffices to show that
$$\sum_i\left[\left\langle\nabla_{F_i}G,\nabla_GF_i\right\rangle
+\left\langle\nabla_{\nabla_GF_i}\,G,F_i\right\rangle\right](p)=0.$$
By the gradient formula for the Hessian in the pointwise Riesz form of
Lemma~\ref{lem:hessian-riesz-trace}, the two summands are the Hessian
values $\Hess(f)_p\left(F_i(p),(\nabla_GF_i)(p)\right)$ and
$\Hess(f)_p\left((\nabla_GF_i)(p),F_i(p)\right)$, which agree by
symmetry of the Hessian (Lemma~\ref{lem:hessian-symmetric}). Expanding
$(\nabla_GF_i)(p)=\sum_ja_{ij}F_j(p)$ over the orthonormal basis, with
$a_{ij}=\left\langle\nabla_GF_i,F_j\right\rangle(p)$, and using
bilinearity of the Hessian, the sum in question equals
$$2\sum_{i,j}a_{ij}\,\Hess(f)_p\left(F_i(p),F_j(p)\right).$$
Since each function $\langle F_i,F_j\rangle$ is constant (equal to
$\delta_{ij}$) near $p$, differentiating it along $G$ and using metric
compatibility gives $a_{ij}+a_{ji}=0$: the matrix $(a_{ij})$ is
antisymmetric. As the matrix
$\Hess(f)_p\left(F_i(p),F_j(p)\right)$ is symmetric, the full
contraction above vanishes, completing the proof of the displayed
identity.
\end{proof}

\begin{lemma}[Laplacian of the square norm of a one-form]
\label{lem:laplacian-square-norm-one-form}
\uses{thm:levi-civita-connection}
Let $\omega$ be a smooth one-form on a Riemannian manifold $(M,g)$ and
let
$$\triangle\omega=g^{ij}\left(\nabla_{\partial_i}\nabla_{\partial_j}\omega
-\nabla_{\nabla_{\partial_i}\partial_j}\omega\right)$$
denote its connection Laplacian, i.e., the operator introduced in the
discussion preceding Lemma~\ref{lem:laplacian-one-form}. Then
$$\frac{1}{2}\triangle|\omega|^2
=\langle\triangle\omega,\omega\rangle+|\nabla\omega|^2,$$
where $\langle\alpha,\beta\rangle=g^{jk}\alpha_j\beta_k$ denotes the inner
product induced by $g$ on one-forms,
$|\omega|^2=\langle\omega,\omega\rangle$, and
$|\nabla\omega|^2=g^{ij}g^{kl}(\nabla\omega)_{ik}(\nabla\omega)_{jl}$ is the
square norm of the two-tensor $(\nabla\omega)(Y,Z)=(\nabla_Y\omega)(Z)$.
\end{lemma}

\begin{proof}
\uses{thm:levi-civita-connection,lem:hessian-symmetric,def:exponential-map,lem:normal-coordinates-expansion}
Both sides of the asserted identity are smooth functions on $M$ defined
independently of any choice of coordinates, so it suffices to verify the
equality at an arbitrary point $p\in M$. Fix Gaussian normal coordinates
$(x^1,\ldots,x^n)$ centered at $p$ (see the discussion following
Definition~\ref{def:exponential-map}). By
Lemma~\ref{lem:normal-coordinates-expansion} we have
$g_{ij}(x)=\delta_{ij}+O(|x|^2)$; hence $g_{ij}(p)=\delta_{ij}$ and
$\partial_kg_{ij}(p)=0$, so by Equation~(\ref{Gamma}) all Christoffel
symbols vanish at $p$: $\Gamma^k_{ij}(p)=0$, and therefore
$\left(\nabla_{\partial_i}\partial_j\right)(p)=0$. Differentiating
$g^{ij}g_{jl}=\delta^i_l$ also gives
$\partial_kg^{ij}(p)=-g^{ia}\left(\partial_kg_{ab}\right)g^{bj}(p)=0$.

\emph{Step 1: compatibility of $\nabla$ with the dual metric.} We claim
that for all smooth one-forms $\alpha,\beta$ and every vector field $W$,
everywhere on $M$,
\begin{equation}\label{eq:dual-metric-compat}
W\langle\alpha,\beta\rangle
=\langle\nabla_W\alpha,\beta\rangle+\langle\alpha,\nabla_W\beta\rangle .
\end{equation}
Recall that the covariant derivative of a one-form (the extension of the
Levi-Civita connection to one-tensors, as in Chapter~\ref{chap:prelim}) is
characterized by the Leibniz rule
$W\left(\alpha(Z)\right)=(\nabla_W\alpha)(Z)+\alpha(\nabla_WZ)$; on
coordinate fields this yields
$(\nabla_{\partial_i}\alpha)_j=\partial_i\alpha_j-\Gamma^k_{ij}\alpha_k$.
In particular $\nabla_W\alpha$ is $C^\infty$-linear in $W$ (replace $W$ by
$fW$ in the Leibniz rule and use $\nabla_{fW}Z=f\nabla_WZ$), and the
left-hand side of Equation~(\ref{eq:dual-metric-compat}) is
$C^\infty$-linear in $W$ as well, being the derivative of a scalar
function in the direction $W$. Hence it suffices to verify
Equation~(\ref{eq:dual-metric-compat}) for $W=\partial_i$ at the
arbitrary point $p$, in the normal coordinates fixed above. There we have
$(\nabla_{\partial_i}\alpha)_j(p)=\partial_i\alpha_j(p)$, since
$\Gamma(p)=0$, and consequently
$$\partial_i\left(g^{jk}\alpha_j\beta_k\right)(p)
=\left(\partial_ig^{jk}\right)(p)\,\alpha_j(p)\beta_k(p)
+\delta^{jk}\left(\partial_i\alpha_j\,\beta_k
+\alpha_j\,\partial_i\beta_k\right)(p)
=\langle\nabla_{\partial_i}\alpha,\beta\rangle(p)
+\langle\alpha,\nabla_{\partial_i}\beta\rangle(p),$$
using $\partial_ig^{jk}(p)=0$ and $g^{jk}(p)=\delta^{jk}$. Since normal
coordinates exist centered at every point of $M$ and both sides of
Equation~(\ref{eq:dual-metric-compat}) are defined independently of
coordinates, the identity holds everywhere on $M$.

\emph{Step 2: two differentiations of $|\omega|^2$.} Set
$u=|\omega|^2=\langle\omega,\omega\rangle$. By
Equation~(\ref{eq:dual-metric-compat}) with $\alpha=\beta=\omega$ we
have, at every point of the coordinate chart,
$$\partial_ju=2\langle\nabla_{\partial_j}\omega,\omega\rangle .$$
For fixed $j$ the field $\nabla_{\partial_j}\omega$ is again a smooth
one-form on the chart, so a second application of
Equation~(\ref{eq:dual-metric-compat}) gives, at every point of the
chart,
$$\partial_i\partial_ju
=2\langle\nabla_{\partial_i}\nabla_{\partial_j}\omega,\omega\rangle
+2\langle\nabla_{\partial_j}\omega,\nabla_{\partial_i}\omega\rangle .$$

\emph{Step 3: evaluation at $p$.} From the definition of the Laplacian of
a function as the trace of its Hessian together with the coordinate
formula for the Hessian in Lemma~\ref{lem:hessian-symmetric},
$$\triangle u
=g^{ij}\left(\partial_i\partial_ju-\Gamma^k_{ij}\partial_ku\right),$$
which at $p$ (where $\Gamma(p)=0$ and $g^{ij}(p)=\delta^{ij}$) reduces,
by Step~2, to
$$\triangle u(p)=\sum_i\partial_i\partial_iu(p)
=2\sum_i\left[\langle\nabla_{\partial_i}\nabla_{\partial_i}\omega,
\omega\rangle+|\nabla_{\partial_i}\omega|^2\right](p).$$
On the other hand, at $p$ the connection Laplacian reduces to
$\triangle\omega(p)
=\sum_i\left(\nabla_{\partial_i}\nabla_{\partial_i}\omega\right)(p)$:
indeed $g^{ij}(p)=\delta^{ij}$, and the term
$\nabla_{\nabla_{\partial_i}\partial_j}\omega$ is tensorial in the
subscript $\nabla_{\partial_i}\partial_j$, which vanishes at $p$. Hence
$$\sum_i\langle\nabla_{\partial_i}\nabla_{\partial_i}\omega,\omega\rangle(p)
=\langle\triangle\omega,\omega\rangle(p),$$
while
$$|\nabla\omega|^2(p)
=g^{ij}g^{kl}(\nabla\omega)_{ik}(\nabla\omega)_{jl}(p)
=\sum_{i,k}\left((\nabla_{\partial_i}\omega)_k(p)\right)^2
=\sum_i|\nabla_{\partial_i}\omega|^2(p).$$
Substituting the last two displays into the formula for $\triangle u(p)$
gives
$$\frac12\triangle|\omega|^2(p)
=\langle\triangle\omega,\omega\rangle(p)+|\nabla\omega|^2(p),$$
and since $p$ was arbitrary the identity holds on all of $M$.
\end{proof}

\begin{lemma}[Bochner formula for functions]
\label{lem:function-bochner-formula}
\uses{def:ricci-curvature}
\lean{MorganTianLib.laplacianAt_metricNormSq_gradientField,MorganTianLib.function_bochner_formula}
\leanok
Let $f$ be a smooth function on a Riemannian manifold $(M,g)$. Then
$$\frac{1}{2}\triangle|\nabla f|^2
=|\Hess(f)|^2
+\left\langle(\nabla f)^*,(\nabla\triangle f)^*\right\rangle
+\Ric\left((\nabla f)^*,(\nabla f)^*\right).$$
Here $(\nabla f)^*$ denotes the gradient of $f$, i.e., the vector field
dual to the one-form $\nabla f=df$; $|\nabla f|$ is the norm of the
one-form $df$, which agrees with the norm of the gradient $(\nabla f)^*$;
$\left\langle(\nabla f)^*,(\nabla\triangle f)^*\right\rangle$ agrees with
the inner product $\langle df,d\triangle f\rangle$ of one-forms; and
$|\Hess(f)|^2=g^{ik}g^{jl}\Hess(f)_{ij}\Hess(f)_{kl}$ is the square norm
of the Hessian.
\end{lemma}

\begin{proof}
\uses{lem:hessian-symmetric,lem:gradient-norm-sq-gradient,def:second-covariant-derivative-vector-field,lem:ricci-commutation-vector-field,lem:ricci-trace-curvature-operator,lem:hessian-riesz-trace,lem:laplacian-trace-commutation,def:hessian-norm-sq,def:ricci-curvature}
\leanok
We first record the compatibility of the two readings of the various
terms. In local coordinates the gradient of a smooth function $u$ is
$(\nabla u)^*=g^{sj}u_{,s}\partial_j$, so for smooth functions $u,v$,
$$\left\langle(\nabla u)^*,(\nabla v)^*\right\rangle
=g_{jk}\left(g^{sj}u_{,s}\right)\left(g^{tk}v_{,t}\right)
=g^{st}u_{,s}v_{,t}
=\langle du,dv\rangle,$$
the inner product of the one-forms $du$ and $dv$; in particular
$\left|(\nabla u)^*\right|=|du|=|\nabla u|$. Note also that
$\triangle f$ is a smooth function by
Lemma~\ref{lem:laplacian-trace-commutation}, so $(\nabla\triangle f)^*$
is defined, and that by the defining property of the gradient
$\left\langle(\nabla f)^*,(\nabla\triangle f)^*\right\rangle
=(\nabla f)^*(\triangle f)$, the derivative of $\triangle f$ in the
direction of the gradient of $f$.

Write $G=(\nabla f)^*$, fix $p\in M$, let $\{e_1,\ldots,e_n\}$ be an
orthonormal basis of $T_pM$, and extend each $e_i$ to a smooth vector
field $E_i$. By Lemma~\ref{lem:gradient-norm-sq-gradient}, the gradient
of the smooth function $|\nabla f|^2$ is $2\nabla_GG$ and its Hessian
is
$$\Hess\left(|\nabla f|^2\right)(X,Y)
=2\left\langle\nabla_X\left(\nabla_GG\right),Y\right\rangle$$
for all vector fields $X,Y$. Taking the metric trace
(Equation~(\ref{laplacformula})),
$$\triangle|\nabla f|^2(p)
=2\sum_i\left\langle\nabla_{E_i}\left(\nabla_GG\right),E_i\right\rangle(p).$$
By the definition of the second covariant derivative
(Definition~\ref{def:second-covariant-derivative-vector-field}),
$$\nabla_{E_i}\left(\nabla_GG\right)
=\nabla^2G(E_i,G)+\nabla_{\nabla_{E_i}G}\,G,$$
and by the Ricci commutation identity
(Lemma~\ref{lem:ricci-commutation-vector-field})
$$\nabla^2G(E_i,G)=\nabla^2G(G,E_i)+{\mathcal R}(E_i,G)G.$$
Substituting and summing the three resulting groups of terms: the
curvature trace is
$$\sum_i\left\langle{\mathcal R}(E_i,G)G,E_i\right\rangle(p)
=\Ric\left(G(p),G(p)\right)$$
by Lemma~\ref{lem:ricci-trace-curvature-operator}; the trace of the
correction terms is
$$\sum_i\left\langle\nabla_{\nabla_{E_i}G}\,G,E_i\right\rangle(p)
=|\Hess(f)|^2(p)$$
by Lemma~\ref{lem:hessian-riesz-trace}; and the trace of the second
covariant derivative is
$$\sum_i\left\langle\nabla^2G(G,E_i),E_i\right\rangle(p)
=G(\triangle f)(p)
=\left\langle(\nabla f)^*,(\nabla\triangle f)^*\right\rangle(p)$$
by Lemma~\ref{lem:laplacian-trace-commutation} and the first
paragraph. Altogether
$$\triangle|\nabla f|^2(p)
=2\left(|\Hess(f)|^2
+\left\langle(\nabla f)^*,(\nabla\triangle f)^*\right\rangle
+\Ric\left((\nabla f)^*,(\nabla f)^*\right)\right)(p),$$
and since $p$ was arbitrary this is the Bochner formula.
\end{proof}

\begin{corollary}[Bochner vanishing]
\label{cor:bochner-vanishing}
\uses{lem:function-bochner-formula,def:hessian-norm-sq,def:ricci-curvature}
\lean{MorganTianLib.hessianNormSqAt_eq_zero_of_bochner,MorganTianLib.hessianAt_eq_zero_of_bochner,MorganTianLib.cov_gradientField_apply_eq_zero_of_bochner}
\leanok
Let $f$ be a smooth function on a Riemannian manifold $(M,g)$ such that
$\triangle f$ and $|\nabla f|^2$ are constant, and let $p\in M$ be a
point with $\Ric\left((\nabla f)^*,(\nabla f)^*\right)(p)\ge0$. Then
$|\Hess(f)|^2(p)=0$, i.e. $\Hess(f)_p=0$. In particular, if $f$ is
harmonic with $|\nabla f|$ constant on a manifold of non-negative Ricci
curvature, then $\Hess(f)\equiv0$ and the gradient $(\nabla f)^*$ is
parallel.
\end{corollary}

\begin{proof}
\uses{lem:function-bochner-formula,def:hessian-norm-sq}
\leanok
Both derivative terms in the Bochner formula
(Lemma~\ref{lem:function-bochner-formula}) vanish at $p$: the
left-hand side $\frac12\triangle|\nabla f|^2$ is the Laplacian of a
constant function, and
$\left\langle(\nabla f)^*,(\nabla\triangle f)^*\right\rangle
=(\nabla f)^*(\triangle f)$ is the derivative of a constant function.
Hence
$$0=|\Hess(f)|^2(p)
+\Ric\left((\nabla f)^*,(\nabla f)^*\right)(p),$$
and since both summands are non-negative — the first as a sum of
squares (Definition~\ref{def:hessian-norm-sq}), the second by
hypothesis — each vanishes. That $\Hess(f)_p=0$ follows since
$|\Hess(f)|^2(p)$ is the sum of the squares of the entries of
$\Hess(f)_p$ in an orthonormal basis. Finally, $\Hess(f)(X,Y)
=\langle\nabla_X(\nabla f)^*,Y\rangle$ for all $X,Y$
(Lemma~\ref{lem:hessian-symmetric}), so $\Hess(f)\equiv0$ forces
$\nabla_X(\nabla f)^*=0$ for every $X$ by non-degeneracy of the
metric: the gradient is parallel.
\end{proof}

\begin{lemma}[Covariant differentiation along smooth maps]
\label{lem:covariant-derivative-along-maps}
\lean{MorganTianLib.chartLocalMap,
  MorganTianLib.chartFieldCoord_curveVelocity_snd_eventually,
  MorganTianLib.chartFieldCoord_curveVelocity_fst_eventually,
  MorganTianLib.curveVelocity_comp_fst_eq,
  MorganTianLib.curveVelocity_comp_snd_eq,
  MorganTianLib.hasCovDerivAlongAt_fst_curveVelocity_snd,
  MorganTianLib.hasCovDerivAlongAt_snd_curveVelocity_fst,
  MorganTianLib.hasCovDerivAlongAt_fst_snd_symm,
  MorganTianLib.HasCovDerivAlongAt,
  MorganTianLib.HasCovDerivAlongAt.unique,
  MorganTianLib.HasCovDerivAlongAt.add,
  MorganTianLib.HasCovDerivAlongAt.const_smul,
  MorganTianLib.HasCovDerivAlongAt.smul_fun,
  MorganTianLib.hasCovDerivAlongAt_comp,
  MorganTianLib.HasCovDerivAlongAt.hasDerivAt_metricInner,
  MorganTianLib.hasGeodesicEquationAt_iff_hasCovDerivAlongAt_velocity_zero,
  MorganTianLib.IsParallelAlong,
  MorganTianLib.IsParallelAlong.metricInner_eq,
  MorganTianLib.IsParallelAlong.apply_eq}
\leanok
\uses{thm:levi-civita-connection,def:geodesic}
Let $(M,g)$ be a Riemannian manifold with Levi-Civita connection
$\nabla$, let $U\subset\Ar^m$ be open with coordinates
$(u^1,\ldots,u^m)$, and let $\alpha\colon U\to M$ be a smooth map. By a
{\em smooth vector field along $\alpha$} we mean a smooth assignment
$u\mapsto W(u)\in T_{\alpha(u)}M$, i.e., in every local coordinate chart
$(x^1,\ldots,x^n)$ on $M$ we can write $W=W^k\,(\partial_k\circ\alpha)$
on the preimage of the chart, with smooth component functions $W^k$.
Then:
\begin{enumerate}
\item For each index $a$, the expressions
$$\frac{DW}{\partial u^a}
=\left(\frac{\partial W^k}{\partial u^a}
+\Gamma^k_{ij}(\alpha)\,\frac{\partial\alpha^i}{\partial u^a}\,W^j\right)
\partial_k ,$$
computed in local coordinate charts on $M$ (where
$\alpha^i=x^i\circ\alpha$), agree on overlaps and hence define a smooth
vector field $\frac{DW}{\partial u^a}$ along $\alpha$. The operation
$W\mapsto\frac{DW}{\partial u^a}$ is $\Ar$-linear, satisfies the Leibniz
rule
$$\frac{D(fW)}{\partial u^a}
=\frac{\partial f}{\partial u^a}\,W+f\,\frac{DW}{\partial u^a}
\qquad\text{for smooth } f\colon U\to\Ar,$$
and satisfies
$\frac{D(Z\circ\alpha)}{\partial u^a}=\nabla_{\partial\alpha/\partial
u^a}Z$ for every vector field $Z$ defined on an open subset of $M$, the
identity holding on the preimage of that subset.
\item For smooth vector fields $W_1,W_2$ along $\alpha$,
$$\frac{\partial}{\partial u^a}\langle W_1,W_2\rangle
=\left\langle\frac{DW_1}{\partial u^a},W_2\right\rangle
+\left\langle W_1,\frac{DW_2}{\partial u^a}\right\rangle .$$
\item For all indices $a,b$,
$$\frac{D}{\partial u^a}\frac{\partial\alpha}{\partial u^b}
=\frac{D}{\partial u^b}\frac{\partial\alpha}{\partial u^a}.$$
\item A smooth curve $\gamma\colon I\to M$ ($I$ an open interval; the
case $m=1$, $t=u^1$) is a geodesic if and only if
$\frac{D\gamma'}{dt}=0$, where $\gamma'=d\gamma/dt$.
\item If $W_1,W_2$ are smooth vector fields along a smooth curve
$\gamma\colon I\to M$ with $\frac{DW_1}{dt}=\frac{DW_2}{dt}=0$ ({\em
parallel fields along $\gamma$}), then $\langle W_1,W_2\rangle$ is
constant on $I$; and if in addition $W_1(t_0)=W_2(t_0)$ for some
$t_0\in I$, then $W_1=W_2$ on all of $I$.
\end{enumerate}
\end{lemma}

\begin{proof}
\leanok
\uses{thm:levi-civita-connection,def:geodesic}
(1) Fix a chart $(O;x^1,\ldots,x^n)$ on $M$ and denote by $D_a$ the
displayed expression computed in this chart, an operation on smooth
vector fields along the restriction of $\alpha$ to $\alpha^{-1}(O)$. It
is visibly $\Ar$-linear and smooth, and replacing $W^k$ by $fW^k$
produces exactly the extra term $\frac{\partial f}{\partial
u^a}W^k\partial_k$, which is the Leibniz rule. For a vector field
$Z=Z^k\partial_k$ defined on an open subset $O'\subset O$ and
$W=Z\circ\alpha$ we have $W^k=Z^k\circ\alpha$, hence $\frac{\partial
W^k}{\partial u^a}=\left(\partial_iZ^k\circ\alpha\right)\frac{\partial
\alpha^i}{\partial u^a}$ on $\alpha^{-1}(O')$, and the displayed
expression becomes
$$D_a(Z\circ\alpha)
=\left(\left(\partial_iZ^k+\Gamma^k_{ij}Z^j\right)\circ\alpha\right)
\frac{\partial\alpha^i}{\partial u^a}\,\partial_k
=\nabla_{\partial\alpha/\partial u^a}Z,$$
by the coordinate formula
$\nabla_{\partial_i}Z=\left(\partial_iZ^k+\Gamma^k_{ij}Z^j\right)\partial_k$
for the Levi-Civita connection (Theorem~\ref{thm:levi-civita-connection}
and the discussion following it) and the $C^\infty$-linearity of
$\nabla_XZ$ in $X$. Thus each chart expression has the three asserted
properties on the preimage of its chart.

For the independence of the chart, let $(\widetilde
O;\tilde x^1,\ldots,\tilde x^n)$ be a second chart with corresponding
expression $\widetilde D_a$, and let $W$ be a smooth field along
$\alpha$. On $\alpha^{-1}(O\cap\widetilde O)$ write
$W=\sum_kW^k(\partial_k\circ\alpha)$ using the first chart. Since
$\widetilde D_a$ is additive, satisfies the Leibniz rule, and satisfies
the composed-field property for vector fields on open subsets of
$O\cap\widetilde O$ --- in particular for the coordinate fields
$\partial_k$ of the first chart, which are defined on $O\cap\widetilde
O$ --- we get
$$\widetilde D_aW
=\sum_k\left(\frac{\partial W^k}{\partial u^a}(\partial_k\circ\alpha)
+W^k\,\widetilde D_a(\partial_k\circ\alpha)\right)
=\sum_k\left(\frac{\partial W^k}{\partial u^a}(\partial_k\circ\alpha)
+W^k\,\nabla_{\partial\alpha/\partial u^a}\partial_k\right),$$
and by the identical computation the right-hand side also equals
$D_aW$. Hence $D_aW=\widetilde D_aW$ on $\alpha^{-1}(O\cap\widetilde
O)$, so the chart expressions patch together to a globally defined
operation $\frac{D}{\partial u^a}$ with the stated properties.

(2) Compute in a chart. On the one hand,
$$\frac{\partial}{\partial u^a}\left(g_{kl}(\alpha)W_1^kW_2^l\right)
=\left(\partial_ig_{kl}\right)(\alpha)\,
\frac{\partial\alpha^i}{\partial u^a}W_1^kW_2^l
+g_{kl}(\alpha)\left(\frac{\partial W_1^k}{\partial u^a}W_2^l
+W_1^k\frac{\partial W_2^l}{\partial u^a}\right).$$
On the other hand, expanding
$\left\langle\frac{DW_1}{\partial u^a},W_2\right\rangle+\left\langle
W_1,\frac{DW_2}{\partial u^a}\right\rangle$ by the formula in (1)
produces the same second group of terms together with
$$\left(g_{ml}\Gamma^m_{ik}+g_{km}\Gamma^m_{il}\right)
\frac{\partial\alpha^i}{\partial u^a}W_1^kW_2^l .$$
The two agree because
$\partial_ig_{kl}=g_{ml}\Gamma^m_{ik}+g_{km}\Gamma^m_{il}$, which is the
coordinate form of the metric-compatibility of the Levi-Civita
connection displayed after Equation~(\ref{Gamma}) (see also the
existence part of the proof of
Theorem~\ref{thm:levi-civita-connection}).

(3) In a chart, by the formula in (1) applied to
$W=\partial\alpha/\partial u^b$ (with components $\partial
\alpha^k/\partial u^b$),
$$\frac{D}{\partial u^a}\frac{\partial\alpha}{\partial u^b}
=\left(\frac{\partial^2\alpha^k}{\partial u^a\partial u^b}
+\Gamma^k_{ij}(\alpha)\frac{\partial\alpha^i}{\partial u^a}
\frac{\partial\alpha^j}{\partial u^b}\right)\partial_k ,$$
which is symmetric in $a$ and $b$ by the equality of mixed partial
derivatives and the symmetry $\Gamma^k_{ij}=\Gamma^k_{ji}$ of the
Christoffel symbols (the torsion-free condition of
Theorem~\ref{thm:levi-civita-connection}).

(4) For $W=\gamma'$ the formula in (1) reads
$$\frac{D\gamma'}{dt}
=\left(\ddot\gamma^k+\Gamma^k_{ij}(\gamma)\,
\dot\gamma^i\dot\gamma^j\right)\partial_k ,$$
whose vanishing is precisely the geodesic equation in local coordinates
displayed after Definition~\ref{def:geodesic}.

(5) The first assertion is immediate from (2):
$\frac{d}{dt}\langle W_1,W_2\rangle=0$. For the second, let $A=\{t\in
I\,|\,W_1(t)=W_2(t)\}$; it is non-empty ($t_0\in A$) and closed in $I$
by continuity. It is also open: around any $t_1\in A$ choose a chart of
$M$ containing $\gamma(t_1)$; on the subinterval of $I$ around $t_1$
that $\gamma$ maps into the chart, the equation $\frac{DW}{dt}=0$ reads
$$\frac{dW^k}{dt}=-\Gamma^k_{ij}(\gamma)\,\dot\gamma^i\,W^j,$$
a linear first-order system of ordinary differential equations with
smooth coefficients, so by the classical uniqueness theorem for such
systems two solutions agreeing at $t_1$ agree on that subinterval.
Since $I$ is an interval, hence connected, $A=I$.
\end{proof}

\begin{lemma}[Covariant differentiation along a curve]
\label{lem:cov-deriv-along-curve}
\lean{MorganTianLib.HasCovDerivAlongAt,
  MorganTianLib.chartFieldCoord,
  MorganTianLib.HasCovDerivAlongAt.unique,
  MorganTianLib.HasCovDerivAlongAt.add,
  MorganTianLib.HasCovDerivAlongAt.const_smul,
  MorganTianLib.HasCovDerivAlongAt.sub,
  MorganTianLib.HasCovDerivAlongAt.smul_fun,
  MorganTianLib.metricInner_eq_chartMetricInner,
  MorganTianLib.HasCovDerivAlongAt.hasDerivAt_metricInner,
  MorganTianLib.IsParallelAlong,
  MorganTianLib.IsParallelAlong.metricInner_eq,
  MorganTianLib.IsParallelAlong.apply_eq,
  MorganTianLib.curveVelocity,
  MorganTianLib.hasGeodesicEquationAt_iff_hasCovDerivAlongAt_velocity_zero,
  MorganTianLib.fieldChartRep,
  MorganTianLib.hasCovDerivAlongAt_comp,
  MorganTianLib.cov_apply_eq_fderiv_add_chartChristoffelContraction,
  MorganTianLib.hasCovDerivAlongAt_comp_cov,
  MorganTianLib.hasCovDerivAlongAt_comp_zero,
  MorganTianLib.isParallelAlong_comp_of_cov_eq_zero,
  MorganTianLib.isParallelAlong_curveVelocity_of_isGeodesic}
\leanok
\uses{thm:levi-civita-connection,def:geodesic}
Let $(M,g)$ be a Riemannian manifold, $\gamma\colon\Ar\to M$ a curve, and
let $W$ be a vector field along $\gamma$, i.e., an assignment $t\mapsto
W(t)\in T_{\gamma(t)}M$. Fix $t_0\in\Ar$, suppose that $\gamma(s)$ lies in
the domain of the canonical chart centred at $\gamma(t_0)$ for all $s$ near
$t_0$, and write $u=\varphi\circ\gamma$ for the chart curve and $V$ for the
coordinate representation of $W$ in this chart (through the induced
trivialization of $TM$ over the chart domain). If $u$ and $V$ are
differentiable at $t_0$, say that $W$ has \emph{covariant derivative}
$$\frac{DW}{dt}(t_0)
=\dot V(t_0)+\Gamma\bigl(\dot u(t_0),W(t_0)\bigr)(u(t_0))
\in T_{\gamma(t_0)}M$$
at $t_0$, where $\Gamma(v,w)(y)=\sum_k\bigl(\sum_{i,j}
\Gamma^k_{ij}(y)v^iw^j\bigr)e_k$ is the contraction against the Christoffel
symbols of $g$ in the chart (Theorem~\ref{thm:levi-civita-connection}), and
$T_{\gamma(t_0)}M$ is identified with the model space through the chart at
$\gamma(t_0)$ itself. Then:
\begin{enumerate}
\item The value $\frac{DW}{dt}(t_0)$ is uniquely determined; the operation
is additive and $\Ar$-homogeneous in $W$, and satisfies the Leibniz rule
$$\frac{D(fW)}{dt}(t_0)=f'(t_0)\,W(t_0)+f(t_0)\,\frac{DW}{dt}(t_0)$$
for scalar functions $f$ differentiable at $t_0$. Moreover, for a smooth
vector field $Z$ on $M$ the restriction $Z\circ\gamma$ has covariant
derivative
$$\frac{D(Z\circ\gamma)}{dt}(t_0)
=\partial_{v}\widehat Z(u(t_0))
+\Gamma\bigl(v,Z(\gamma(t_0))\bigr)(u(t_0)),
\qquad v=\dot u(t_0),$$
where $\widehat Z$ is the coordinate representation of $Z$ in the chart at
$\gamma(t_0)$. This value is the Levi-Civita covariant derivative of $Z$
in the direction of the velocity: for every smooth vector field $X$ on $M$
with $X(\gamma(t_0))=v$,
$$\frac{D(Z\circ\gamma)}{dt}(t_0)=\bigl(\nabla_XZ\bigr)(\gamma(t_0));$$
equivalently, at the centre $x$ of its own chart,
$$(\nabla_XZ)(x)=\partial_{X(x)}\widehat Z(\varphi(x))
+\Gamma\bigl(X(x),Z(x)\bigr)(\varphi(x)),$$
so that $(\nabla_XZ)(x)$ depends on $X$ only through $X(x)$. In
particular, a smooth field $Z$ with $\nabla Z\equiv0$ restricts to a
parallel field (in the sense of item (4)) along every curve that satisfies
the standing chart-regularity hypotheses at each time.
\item (Metric compatibility.) If $W_1,W_2$ have covariant derivatives
$D_1,D_2$ at $t_0$, then $t\mapsto\langle W_1(t),W_2(t)\rangle$ is
differentiable at $t_0$ with derivative
$\langle D_1,W_2(t_0)\rangle+\langle W_1(t_0),D_2\rangle$.
\item A curve $\gamma$ whose chart curve at $\gamma(t_0)$ is differentiable
near $t_0$ and which stays in the chart at $\gamma(t_0)$ near $t_0$
satisfies the geodesic equation at $t_0$ if and only if its velocity field
$\gamma'$ has covariant derivative $0$ at $t_0$.
\item (Parallel fields.) Call $W$ \emph{parallel along $\gamma$} when
$\frac{DW}{dt}(t)=0$ for every $t$. The inner product of two parallel
fields along $\gamma$ is constant, and two parallel fields that agree at
one time agree at every time.
\end{enumerate}
\end{lemma}

\begin{proof}
\leanok
\uses{thm:levi-civita-connection,def:geodesic,lem:covariant-derivative-along-maps}
(1) Uniqueness holds because derivatives are unique. The trivialization of
$TM$ over the chart domain is linear on each fibre, so the coordinate
representation of $W_1+W_2$ (resp.\ $fW$) is $V_1+V_2$ (resp.\ $fV$);
additivity and homogeneity follow since differentiation and
$\Gamma(v,\cdot)(y)$ are linear, and for $fW$ the product rule gives
$\frac{d}{dt}(fV)=f'V+f\dot V$, whence the Leibniz rule using
$V(t_0)=W(t_0)$ (the trivialization is the identity on the fibre over the
centre of the chart) and the homogeneity of $\Gamma$ in its second slot.
For the composed field: the coordinate representation of $Z\circ\gamma$
equals $\widehat Z\circ u$ near $t_0$, where $\widehat Z$, the coordinate
representation of the section $Z$, is smooth near $u(t_0)$ because $Z$ is a
smooth section of $TM$ and the trivialization and the inverse chart are
smooth; the chain rule gives $\dot V(t_0)=\partial_v\widehat Z(u(t_0))$ and
the displayed formula follows from the definition.

For the identification with $\nabla_XZ$ (the \emph{frame bridge}), work at
$x=\gamma(t_0)$ in the chart centred at $x$, with chart frame
$\partial_1,\dots,\partial_n$. Since sections of $TM$ read through the
trivialization are smooth over the chart domain, each chart coordinate
$Z^k$ of $Z$ (and $X^i$ of $X$) is smooth there, and multiplication by a
bump function equal to $1$ near $x$ and supported in the chart domain
extends it to a globally smooth function agreeing with $Z^k$ (resp.\
$X^i$) near $x$. Likewise extend the chart frame to global smooth fields
$X_1,\dots,X_n$ agreeing with $\partial_1,\dots,\partial_n$ near $x$; by
the local coordinate form of the Levi-Civita connection
(Theorem~\ref{thm:levi-civita-connection}),
$(\nabla_{X_i}X_k)(x)=\sum_m\Gamma^m_{ik}(\varphi(x))\,X_m(x)$. The value
$(\nabla_XZ)(x)$ depends only on the germ of $Z$ at $x$ — a field
vanishing near $x$ can be written as $fZ'$ with $f$ smooth, $f(x)=0$, and
the Leibniz rule kills both terms at $x$ — and only on the value $X(x)$,
by $\mathcal D(M)$-linearity of $\nabla$ in the direction slot. Hence we
may replace $Z$ by $\sum_kZ^kX_k$ and $X$ by $\sum_iX^iX_i$. Expanding
with additivity, $\mathcal D(M)$-homogeneity and the Leibniz rule, and
evaluating at $x$ (where $X_k(x)=e_k$),
$$(\nabla_XZ)(x)
=\sum_k\Bigl(X(Z^k)(x)\,e_k+Z^k(x)\sum_iX^i(x)(\nabla_{X_i}X_k)(x)\Bigr)
=\partial_{X(x)}\widehat Z(\varphi(x))
+\Gamma\bigl(X(x),Z(x)\bigr)(\varphi(x)):$$
the first summand is the chain rule through the chart applied
coordinatewise ($X(Z^k)(x)=\partial_{X(x)}\widehat Z^k(\varphi(x))$,
summed against the model basis), and the second collects the Christoffel
terms into the contraction $\Gamma$. With $X(x)=v$ this is the
composed-field display, proving
$\frac{D(Z\circ\gamma)}{dt}(t_0)=(\nabla_XZ)(\gamma(t_0))$. Finally, if
$\nabla Z\equiv0$ then for each time $t$ choose (extension of a tangent
vector to a global smooth field) some $X$ with $X(\gamma(t))=\dot u(t)$;
then $\frac{D(Z\circ\gamma)}{dt}(t)=(\nabla_XZ)(\gamma(t))=0$, so
$Z\circ\gamma$ is parallel along $\gamma$.

(2) In the chart at $\gamma(t_0)$, for $s$ near $t_0$,
$$\langle W_1(s),W_2(s)\rangle
=\sum_{i,j}G_{ij}(u(s))\,V_1^i(s)\,V_2^j(s),$$
where $G_{ij}$ is the Gram matrix of $g$ in the chart: expand both vectors
in the chart frame, with coefficients the coordinates of their fibre
representations. Differentiating at $t_0$ by the product and chain rules
and inserting the coordinate identity
$$\partial_kG_{ij}
=\sum_m\left(G_{mj}\Gamma^m_{ki}+G_{im}\Gamma^m_{kj}\right),$$
the coordinate form of the metric compatibility of the Levi-Civita
connection (as displayed in the proof of
Lemma~\ref{lem:covariant-derivative-along-maps}(2)), regroups the terms
into
$$\left\langle\dot V_1+\Gamma(\dot u,V_1)(u),V_2\right\rangle
+\left\langle V_1,\dot V_2+\Gamma(\dot u,V_2)(u)\right\rangle$$
evaluated at $t_0$, which is the asserted formula since
$V_i(t_0)=W_i(t_0)$.

(3) Both statements are read in the chart at $\gamma(t_0)$. The geodesic
equation at $t_0$ is $\ddot u(t_0)+\Gamma(\dot u(t_0),\dot
u(t_0))(u(t_0))=0$ (Definition~\ref{def:geodesic}). The coordinate
representation of the velocity field $\gamma'$ in the chart at
$\gamma(t_0)$ agrees with $\dot u$ near $t_0$: for $s$ near $t_0$ the
velocity $\gamma'(s)$ is read in the chart at $\gamma(s)$, the transition
between the charts at $\gamma(s)$ and at $\gamma(t_0)$ is smooth with
fibre derivative the coordinate change of tangent vectors, and the chain
rule together with the fact that the coordinate change from a chart to
itself is the identity converts the derivative of the chart curve at
$\gamma(s)$ into $\dot u(s)$. Hence the covariant derivative of $\gamma'$
at $t_0$ is $\ddot u(t_0)+\Gamma(\dot u(t_0),\dot u(t_0))(u(t_0))$, whose
vanishing is precisely the geodesic equation.

(4) By part (2) the function $t\mapsto\langle W_1(t),W_2(t)\rangle$ has
vanishing derivative at every time when both fields are parallel, hence is
constant by the mean value theorem. If moreover $W_1(t_1)=W_2(t_1)$, then
$W=W_1-W_2$ is parallel by part (1), so $|W|^2$ is constant and vanishes at
$t_1$; hence $\langle W(t),W(t)\rangle=0$ for all $t$, and
positive-definiteness of $g$ forces $W(t)=0$ for every $t$, i.e.,
$W_1=W_2$.
\end{proof}


\begin{lemma}[Flow of a parallel gradient field]
\label{lem:parallel-gradient-flow}
\lean{MorganTianLib.isParallelAlong_gradientField_comp_of_bochner,
  MorganTianLib.curveVelocity_eq_gradientField_of_bochner,
  MorganTianLib.metricInner_curveVelocity_self_of_bochner,
  MorganTianLib.hasDerivAt_comp_of_bochner,
  MorganTianLib.comp_eq_add_mul_of_bochner,
  MorganTianLib.isMIntegralCurve_gradientField_of_bochner,
  MorganTianLib.isMIntegralCurve_smoothVectorField_eq,
  MorganTianLib.isMIntegralCurve_smoothVectorField_comp_add,
  MorganTianLib.IsContGeodesicallyComplete,
  MorganTianLib.isContGeodesicallyComplete_of_complete,
  MorganTianLib.exists_isMIntegralCurve_gradientField_of_complete,
  MorganTianLib.exists_gradientFlowLine_of_bochner,
  MorganTianLib.smoothVectorFieldFlow,
  MorganTianLib.smoothVectorFieldFlow_zero,
  MorganTianLib.smoothVectorFieldFlow_add,
  MorganTianLib.smoothVectorFieldFlow_bijective,
  MorganTianLib.continuous_smoothVectorFieldFlow,
  MorganTianLib.smoothVectorFieldFlowHomeomorph,
  MorganTianLib.comp_smoothVectorFieldFlow_gradientField_of_bochner,
  MorganTianLib.fderiv_fieldChartRep_gradientField_of_bochner,
  MorganTianLib.chartMetricInner_variational_eq_left,
  MorganTianLib.exists_localFlow_hasStrictFDerivAt,
  MorganTianLib.metricPreserving_of_localFlowRep,
  MorganTianLib.exists_flowIsometryBoxAt,
  MorganTianLib.metricPreserving_smoothVectorFieldFlow_of_bochner,
  HasStrictFDerivAt.continuousAt_derivMap,
  MorganTianLib.contDiffOn_one_of_hasStrictFDerivAt,
  MorganTianLib.contMDiffAt_smoothVectorFieldFlow_of_bochner,
  MorganTianLib.contMDiff_smoothVectorFieldFlow_of_bochner,
  MorganTianLib.exists_localFlow_hasStrictFDerivAt_uncurry,
  MorganTianLib.exists_flowJointC1BoxAt,
  MorganTianLib.contMDiffAt_smoothVectorFieldFlow_uncurry_of_bochner,
  MorganTianLib.contMDiff_smoothVectorFieldFlow_uncurry_of_bochner,
  MorganTianLib.mfderiv_smoothVectorFieldFlow_gradientField_of_bochner,
  MorganTianLib.pathELength_comp_eq_of_metricPreserving,
  MorganTianLib.edist_smoothVectorFieldFlow_le_of_bochner,
  MorganTianLib.isometry_smoothVectorFieldFlow_of_bochner}
% Formalization status: parts (1), (3), (4) fully covered; part (2) covered at
% the level of existence, uniqueness, group law, bijectivity, joint continuity
% (homeomorphism), C^1 regularity of each map theta_t in the space
% variable (contMDiff_smoothVectorFieldFlow_of_bochner), and JOINT C^1
% regularity of theta as a map R x X -> X
% (contMDiff_smoothVectorFieldFlow_uncurry_of_bochner: the chart-local flow
% is jointly C^1 on an open space-time box,
% exists_localFlow_hasStrictFDerivAt_uncurry, by continuity of the Neumann
% resolvent of the variational Volterra equation in the base point plus
% continuous-partials => joint strict differentiability; the manifold
% statement glues one-sided joint flow boxes, exists_flowJointC1BoxAt, along
% the compact orbit arc with a backward-anchor group-law reparametrization).
% The flow differential also carries the gradient field to itself,
% d(theta_t)(grad f) = grad f o theta_t
% (mfderiv_smoothVectorFieldFlow_gradientField_of_bochner). Part (4) is
% formalized in BOTH senses: the differential preserves the Riemannian inner
% product (metricPreserving_smoothVectorFieldFlow_of_bochner) AND theta_t is
% a genuine isometry of the metric space (M, d) when the ambient distance is
% the Riemannian distance of g (isometry_smoothVectorFieldFlow_of_bochner,
% by pushing C^1 paths through theta_t with preserved pathELength). Only the
% joint C-infinity smoothness of theta in (t,x) asserted in (2) remains
% unformalized (it needs C^k flow dependence for all k); nothing downstream
% in this chapter consumes more than the joint C^1 + isometry data.
\uses{thm:levi-civita-connection,def:geodesic,thm:hopf-rinow,def:exponential-map}
Let $(X,g)$ be a complete Riemannian manifold and let $B\colon X\to\Ar$
be a smooth function with $\Hess(B)\equiv 0$ and $|\nabla B|\equiv 1$.
Denote by $V=(\nabla B)^*$ the gradient vector field of $B$, so that
$\langle V,W\rangle=dB(W)=W(B)$ for every vector field $W$ and
$|V|=|\nabla B|\equiv 1$. Then:
\begin{enumerate}
\item $V$ is parallel: $\nabla_WV=0$ for every vector field $W$.
\item For every $x\in X$ the curve $t\mapsto\exp_x(tV(x))$ is defined
for all $t\in\Ar$ and is the maximal integral curve of $V$ through $x$;
it is a unit-speed geodesic. Consequently $V$ is a complete vector
field, whose flow
$$\theta\colon\Ar\times X\to X,\qquad
\theta_t(x)=\exp_x\left(tV(x)\right),$$
is smooth and satisfies $\theta_0=\mathrm{id}$ and
$\theta_s\circ\theta_t=\theta_{s+t}$; in particular each
$\theta_t$ is a diffeomorphism of $X$ with inverse $\theta_{-t}$.
\item $B(\theta_t(x))=B(x)+t$ for all $x\in X$ and $t\in\Ar$.
\item Each $\theta_t$ is an isometry of $(X,g)$, i.e.,
$\theta_t^*g=g$.
\end{enumerate}
\end{lemma}

\begin{proof}
\uses{thm:levi-civita-connection,lem:hessian-symmetric,def:geodesic,thm:hopf-rinow,def:exponential-map,lem:covariant-derivative-along-maps}
First we record the norm identities. In local coordinates
$V=g^{sj}B_{,s}\partial_j$, so
$$\langle V,V\rangle
=g_{jk}\left(g^{sj}B_{,s}\right)\left(g^{tk}B_{,t}\right)
=g^{st}B_{,s}B_{,t}=|\nabla B|^2=1,
\qquad
dB(V)=B_{,j}g^{sj}B_{,s}=|\nabla B|^2=1.$$

(1) For vector fields $W,Y$, metric compatibility of the Levi-Civita
connection (Theorem~\ref{thm:levi-civita-connection}) gives
$$\langle\nabla_WV,Y\rangle
=W\langle V,Y\rangle-\langle V,\nabla_WY\rangle
=W(Y(B))-(\nabla_WY)(B)
=\Hess(B)(W,Y)=0,$$
using $\langle V,Y\rangle=Y(B)$ and the definition of the Hessian,
Equation~(\ref{Hessian}) (cf.\ Lemma~\ref{lem:hessian-symmetric}). Since
$Y$ is arbitrary and $g$ is non-degenerate, $\nabla_WV=0$; as
$\nabla_WV$ is tensorial in $W$, this holds for every tangent vector.

(2) Fix $x\in X$ and let $\gamma_x$ be the geodesic with
$\gamma_x(0)=x$ and $\gamma_x'(0)=V(x)$; since $X$ is complete it is
defined for all $t\in\Ar$ (Theorem~\ref{thm:hopf-rinow}). Moreover
$\gamma_x(t)=\exp_x(tV(x))$ for every $t$: the curve
$s\mapsto\gamma_x(ts)$ is again a geodesic, with initial velocity
$tV(x)$ (this rescaling property is immediate from the coordinate form
of the geodesic equation), so by the uniqueness of geodesics with
prescribed initial conditions it is the geodesic $\gamma_{tV(x)}$ of
Definition~\ref{def:exponential-map}, and evaluating at $s=1$ gives
$\exp_x(tV(x))=\gamma_{tV(x)}(1)=\gamma_x(t)$. Along $\gamma_x$ consider the
two vector fields $\gamma_x'$ and $V\circ\gamma_x$. The first is
parallel by Lemma~\ref{lem:covariant-derivative-along-maps}(4), since
$\gamma_x$ is a geodesic; the second is parallel since, by
Lemma~\ref{lem:covariant-derivative-along-maps}(1) and part~(1) of the
present lemma,
$$\frac{D(V\circ\gamma_x)}{dt}=\nabla_{\gamma_x'}V=0 .$$
The two fields agree at $t=0$, so by
Lemma~\ref{lem:covariant-derivative-along-maps}(5) they agree for all
$t$:
$$\gamma_x'(t)=V\left(\gamma_x(t)\right)\qquad\text{for all }t\in\Ar,$$
i.e., $\gamma_x$ is an integral curve of $V$ through $x$, defined on all
of $\Ar$; it has unit speed since $|V|\equiv1$. Integral curves of $V$
are unique: in a chart the equation $\dot c=V(c)$ is a first-order
system with smooth right-hand side, so two integral curves agreeing at
one time agree near it by the classical uniqueness theorem, and the set
of times where two integral curves with a common value agree is open
and closed in their common (connected) interval of definition. Hence
$\gamma_x$ is the unique maximal integral curve of $V$ through $x$ and
$V$ is complete, with flow $\theta_t(x)=\gamma_x(t)=\exp_x(tV(x))$. The
map $\theta$ is smooth in $(t,x)$: solutions of the geodesic equation
depend smoothly on their initial conditions (classical ODE theory, as
recalled after Definition~\ref{def:geodesic}), so
$(x,v)\mapsto\exp_x(v)$ is smooth on $TX$ (which, by completeness, is
its whole domain), and $\theta$ is the composition of this map with the
smooth map $(t,x)\mapsto tV(x)$. For the group law, fix $t$ and $x$;
both $s\mapsto\theta_{s+t}(x)$ and $s\mapsto\theta_s(\theta_t(x))$ are
integral curves of $V$ taking the value $\theta_t(x)$ at $s=0$, hence
they coincide by uniqueness. In particular
$\theta_t\circ\theta_{-t}=\theta_0=\mathrm{id}=\theta_{-t}\circ\theta_t$,
so each $\theta_t$ is a diffeomorphism.

(3) For fixed $x$ the function $t\mapsto B(\theta_t(x))$ has derivative
$$dB\left(\frac{\partial}{\partial t}\theta_t(x)\right)
=dB\left(V(\theta_t(x))\right)=1$$
by the first paragraph, and value $B(x)$ at $t=0$; hence
$B(\theta_t(x))=B(x)+t$.

(4) Fix $x\in X$, $w\in T_xX$ and $t\in\Ar$. Choose a smooth curve
$\sigma\colon(-\epsilon,\epsilon)\to X$ with $\sigma(0)=x$ and
$\sigma'(0)=w$, and consider the smooth map
$$\alpha\colon(-\epsilon,\epsilon)\times\Ar\to X,\qquad
\alpha(s,t)=\theta_t(\sigma(s)).$$
For each fixed $s$ the curve $t\mapsto\alpha(s,t)$ is an integral curve
of $V$, so $\partial\alpha/\partial t=V\circ\alpha$. By
Lemma~\ref{lem:covariant-derivative-along-maps}(3), then
Lemma~\ref{lem:covariant-derivative-along-maps}(1) applied to the
composed field $V\circ\alpha$, and finally part~(1),
$$\frac{D}{\partial t}\frac{\partial\alpha}{\partial s}
=\frac{D}{\partial s}\frac{\partial\alpha}{\partial t}
=\frac{D(V\circ\alpha)}{\partial s}
=\nabla_{\partial\alpha/\partial s}V=0 .$$
Hence for each fixed $s$ the field $t\mapsto\partial\alpha/\partial
s(s,t)$ is parallel along the curve $t\mapsto\alpha(s,t)$ --- by the
chart formula of Lemma~\ref{lem:covariant-derivative-along-maps}(1),
covariant differentiation in $t$ along $\alpha$ agrees with covariant
differentiation along the curve $t\mapsto\alpha(s,t)$ for each fixed
$s$ --- and so by Lemma~\ref{lem:covariant-derivative-along-maps}(5)
its norm is independent of $t$. Now
$$\frac{\partial\alpha}{\partial s}(0,t)
=\frac{d}{ds}\Bigl|_{s=0}\Bigr.\theta_t(\sigma(s))
=d(\theta_t)_x(w),$$
and at $t=0$ this is $w$. Therefore
$\left|d(\theta_t)_x(w)\right|=|w|$ for every $x$, $w$ and $t$. Since
$d(\theta_t)_x$ is linear and preserves norms, the polarization identity
$\langle v,w\rangle=\frac14\left(|v+w|^2-|v-w|^2\right)$ shows that it
preserves inner products; thus $\theta_t^*g=g$, and by part~(2) each
$\theta_t$ is a diffeomorphism, hence an isometry of $(X,g)$.
\end{proof}

\begin{lemma}[Level sets of a function with parallel unit gradient]
\label{lem:parallel-gradient-level-sets}
\lean{MorganTianLib.nonempty_preimage_of_bochner,
  MorganTianLib.isClosed_preimage_of_continuous,
  MorganTianLib.metricInner_curveVelocity_gradientField_of_bochner,
  MorganTianLib.hasDerivAt_comp_zero_of_bochner,
  MorganTianLib.comp_eq_of_bochner_of_metricInner_eq_zero,
  MorganTianLib.mapsTo_preimage_of_bochner_of_metricInner_eq_zero,
  MorganTianLib.metricInner_gradientField_eq_zero_of_mfderiv_eq_zero,
  MorganTianLib.isConnected_preimage_of_bochner,
  MorganTianLib.levelSetFlowIsometry,
  MorganTianLib.exists_openPartialHomeomorph_comp_symm_eq_affine,
  MorganTianLib.hasFDerivAt_comp_extChartAt_symm,
  MorganTianLib.mfderiv_ne_zero_of_metricNormSq_gradientField_eq_one,
  MorganTianLib.exists_extChartAt_openPartialHomeomorph_comp_symm_eq_affine,
  MorganTianLib.levelDifferential,
  MorganTianLib.levelRieszVector,
  MorganTianLib.mem_orthogonal_levelRieszVector_iff,
  MorganTianLib.levelHyperplaneBasis,
  MorganTianLib.sliceProj,
  MorganTianLib.sliceEmb,
  MorganTianLib.adaptedStraightening,
  MorganTianLib.levelSliceChart,
  MorganTianLib.levelSetChartedSpace,
  MorganTianLib.contDiffOn_levelSliceChart_trans,
  MorganTianLib.isManifold_levelSet,
  MorganTianLib.contMDiff_levelSet_val,
  MorganTianLib.mfderiv_levelSet_val_injective,
  MorganTianLib.range_mfderiv_levelSet_val,
  MorganTianLib.exists_hasCovDerivAlongAt_curveVelocity,
  MorganTianLib.metricInner_gradientField_curveVelocity_eq_zero_of_level,
  MorganTianLib.metricInner_gradientField_covDerivVelocity_eq_zero_of_level,
  MorganTianLib.levelSetInducedMetric,
  MorganTianLib.levelSetInducedMetric_metricInner,
  MorganTianLib.metricInner_mfderiv_levelSet_val,
  MorganTianLib.exists_hasCovDerivAlongAt_curveVelocity_levelSet}
\uses{lem:parallel-gradient-flow,def:geodesic,def:riemannian-metric}
Let $(X,g)$, $B$ and $V=(\nabla B)^*$ be as in
Lemma~\ref{lem:parallel-gradient-flow}, i.e., $X$ is complete,
$\Hess(B)\equiv0$ and $|\nabla B|\equiv1$. For $c\in\Ar$ let
$N_c=B^{-1}(c)$, equipped with the induced Riemannian metric $g_{N_c}$.
Then:
\begin{enumerate}
\item $N_c$ is non-empty, closed in $X$, and a smooth embedded
hypersurface: around each $y\in N_c$ there is a chart of $X$ in which
$B$ is the last coordinate function, so that $N_c$ is locally the
slice where the last coordinate equals $c$. For each $y\in N_c$ we have
$T_yN_c=\ker(dB_y)=V(y)^\perp$, so $V$ restricts to a unit normal field
along $N_c$. Moreover every smooth map from a smooth manifold into $X$
with image contained in $N_c$ is smooth as a map into $N_c$.
\item Every geodesic of $X$ whose initial point lies on $N_c$ and whose
initial velocity is tangent to $N_c$ remains in $N_c$.
\item For every smooth curve $\gamma$ in $N_c$ the covariant derivatives
of $\gamma'$ along $\gamma$ in $(X,g)$ and in $(N_c,g_{N_c})$ coincide;
in particular $\gamma$ is a geodesic of $(N_c,g_{N_c})$ if and only if
it is a geodesic of $(X,g)$. Hence $N_c$ is totally geodesic: every
geodesic of the induced metric on $N_c$ is a geodesic of $X$, and, by
item~2, every $X$-geodesic starting on $N_c$ tangent to $N_c$ stays in
$N_c$ and is a geodesic of the induced metric.
\end{enumerate}
\end{lemma}

\begin{proof}
\uses{lem:parallel-gradient-flow,lem:covariant-derivative-along-maps,def:geodesic,def:exponential-map,thm:levi-civita-connection}
Throughout, $\nabla$ and $\frac{D}{dt}$ denote the Levi-Civita
connection of $(X,g)$ and covariant differentiation along curves and
parametrized surfaces in $X$
(Lemma~\ref{lem:covariant-derivative-along-maps}); a superscript $N$
indicates the corresponding objects of the Riemannian manifold
$(N_c,g_{N_c})$.

(1) $N_c$ is non-empty: for any $x\in X$,
Lemma~\ref{lem:parallel-gradient-flow}(3) gives
$B\left(\theta_{c-B(x)}(x)\right)=c$. It is closed as the preimage of
the closed set $\{c\}$ under the continuous map $B$. Next,
$dB_y(V(y))=|V(y)|^2=1$, so $dB$ vanishes nowhere. Fix $y\in N_c$ and
any chart $(x^1,\ldots,x^n)$ around $y$; since $dB_y\neq0$ we may, after
permuting the coordinates, assume $\partial B/\partial x^n(y)\neq0$.
Then the map $x\mapsto\left(x^1,\ldots,x^{n-1},B(x)\right)$ has
invertible differential at $y$ and hence restricts to a chart on some
neighborhood of $y$ by the inverse function theorem. In this chart $B$
is the $n^{th}$ coordinate function, and the intersection of $N_c$ with
the chart domain is the slice where the last coordinate equals $c$.
Consequently $N_c$ is a smooth embedded hypersurface, and its tangent
space at $y$ is spanned by the first $n-1$ coordinate fields of the
slice chart, on which $dB$ (the differential of the last coordinate)
vanishes; by equality of dimensions, $T_yN_c=\ker(dB_y)$. Since
$\langle V(y),w\rangle=dB_y(w)$ for $w\in T_yX$, we have
$\ker(dB_y)=V(y)^\perp$, and $|V(y)|=1$, so $V$ is a unit normal field
along $N_c$. Finally, if $P$ is a smooth manifold and $h\colon P\to X$
is smooth with $h(P)\subset N_c$, then, composing with a slice chart,
the last component of $h$ is the constant $c$ while the remaining $n-1$
components are smooth and are precisely local coordinates on $N_c$;
hence $h$ is smooth as a map into $N_c$. The same argument shows: a
smooth vector field $W$ along a smooth curve $\gamma$ in $N_c$ (in the
sense of fields along curves in $X$) which satisfies $W(t)\in
T_{\gamma(t)}N_c$ for all $t$ is also smooth as a vector field along
$\gamma$ in $N_c$ --- in a slice chart its last component vanishes
identically and its remaining components are smooth --- and conversely.

(2) Let $\gamma\colon I\to X$ be a geodesic with $0\in I$,
$\gamma(0)\in N_c$ and $\gamma'(0)\in T_{\gamma(0)}N_c$. Consider
$h(t)=\left\langle V(\gamma(t)),\gamma'(t)\right\rangle
=dB(\gamma'(t))=(B\circ\gamma)'(t)$. By
Lemma~\ref{lem:covariant-derivative-along-maps}(2), then
Lemma~\ref{lem:covariant-derivative-along-maps}(1) together with
Lemma~\ref{lem:parallel-gradient-flow}(1), and
Lemma~\ref{lem:covariant-derivative-along-maps}(4),
$$h'(t)=\left\langle\frac{D(V\circ\gamma)}{dt},\gamma'\right\rangle
+\left\langle V\circ\gamma,\frac{D\gamma'}{dt}\right\rangle
=\left\langle\nabla_{\gamma'}V,\gamma'\right\rangle+0=0 .$$
Since $h(0)=dB(\gamma'(0))=0$, the function $B\circ\gamma$ has
identically vanishing derivative, so $B\circ\gamma\equiv c$ and
$\gamma$ remains in $N_c$.

(3) Let $\gamma\colon I\to X$ be a smooth curve with image in $N_c$; by
item~1 it is also a smooth curve in $N_c$.

\emph{Tangency of the acceleration.} Since $B\circ\gamma\equiv c$ we
have $\left\langle V\circ\gamma,\gamma'\right\rangle
=dB(\gamma')\equiv0$; differentiating as in item~2,
$$0=\left\langle\frac{D(V\circ\gamma)}{dt},\gamma'\right\rangle
+\left\langle V\circ\gamma,\frac{D\gamma'}{dt}\right\rangle
=\left\langle V\circ\gamma,\frac{D\gamma'}{dt}\right\rangle,$$
so that, by item~1,
\begin{equation}\label{eq:tangential-acceleration}
\frac{D\gamma'}{dt}(t)\in T_{\gamma(t)}N_c\qquad\text{for all }t\in I .
\end{equation}
% Formalization status (item 3): the tangential-acceleration identity
% \eqref{eq:tangential-acceleration} is fully formalized. The acceleration
% D\gamma'/dt exists for every smooth curve
% (exists_hasCovDerivAlongAt_curveVelocity, via C^2 regularity of the chart
% curve), the pairing <V o gamma, gamma'> vanishes pointwise on N_c
% (metricInner_gradientField_curveVelocity_eq_zero_of_level), and
% differentiating it with the metric product rule together with the
% parallelism of V o gamma yields <V o gamma, D\gamma'/dt> = 0, i.e. the
% acceleration is orthogonal to the unit normal
% (metricInner_gradientField_covDerivVelocity_eq_zero_of_level). The induced
% Riemannian metric g_{N_c} on the level set is now formalized
% (levelSetInducedMetric, the pullback of g through the smooth immersion
% i : N_c -> X, DoCarmo DCInducedMetric), with the pullback identity
% <v,w>_{g_{N_c}} = <di v, di w>_g (levelSetInducedMetric_metricInner,
% metricInner_mfderiv_levelSet_val) and the existence of the intrinsic
% acceleration D^{N_c}gamma'/dt (exists_hasCovDerivAlongAt_curveVelocity_levelSet);
% D^{N_c} is available both as the along-curve operator on (N_c, g_{N_c}) and via
% DoCarmo's RiemannianMetric.leviCivitaConnection. What remains unformalized for
% the full totally-geodesic conclusion is the Gauss-formula comparison
% di(D^{N_c}gamma'/dt) = D gamma'/dt (the two covariant derivatives agree because
% the ambient acceleration is already tangent to N_c by
% eq:tangential-acceleration); equivalently, the first-variation-of-energy
% comparison below. This step is a chart-Christoffel identity relating the
% Christoffel symbols of g_{N_c} in slice charts to those of g.

\emph{First variation of the energy.} We record a computation in an
arbitrary Riemannian manifold $(M,h)$, to be applied both to $(X,g)$
and to $(N_c,g_{N_c})$. Let $\gamma\colon I\to M$ be a smooth curve and
let $[a,b]\subset I$ be a compact subinterval. By a {\em variation of
$\gamma$ supported in $(a,b)$} we mean a smooth map
$\tilde\gamma\colon J\times(-\epsilon,\epsilon)\to M$, where
$\epsilon>0$ and $J$ is an open interval with $[a,b]\subset J\subset
I$, such that $\tilde\gamma(t,0)=\gamma(t)$ for all $t\in J$ and such
that $\tilde\gamma(t,u)=\gamma(t)$ for all $u$ whenever $t$ lies
outside some compact subset of $(a,b)$; its {\em variation field} is
the smooth vector field $Y(t)=\partial\tilde\gamma/\partial u\,(t,0)$
along $\gamma|_J$, which vanishes outside a compact subset of $(a,b)$.
Set
$$E(u)=\frac12\int_a^b
h\left(\frac{\partial\tilde\gamma}{\partial t},
\frac{\partial\tilde\gamma}{\partial t}\right)dt .$$
Differentiating under the integral sign ($\tilde\gamma$ is smooth on a
neighborhood of $[a,b]\times\{0\}$) and using
Lemma~\ref{lem:covariant-derivative-along-maps}(2), then (3),
\begin{eqnarray*}
E'(u) & = & \int_a^bh\left(\frac{D}{\partial u}
\frac{\partial\tilde\gamma}{\partial t},
\frac{\partial\tilde\gamma}{\partial t}\right)dt
\ =\ \int_a^bh\left(\frac{D}{\partial t}
\frac{\partial\tilde\gamma}{\partial u},
\frac{\partial\tilde\gamma}{\partial t}\right)dt\\
& = & \Bigl[h\left(\frac{\partial\tilde\gamma}{\partial u},
\frac{\partial\tilde\gamma}{\partial t}\right)\Bigr]_{t=a}^{t=b}
-\int_a^bh\left(\frac{\partial\tilde\gamma}{\partial u},
\frac{D}{\partial t}\frac{\partial\tilde\gamma}{\partial t}\right)dt ,
\end{eqnarray*}
where the last equality uses
Lemma~\ref{lem:covariant-derivative-along-maps}(2) in the form
$h\left(\frac{D}{\partial t}\frac{\partial\tilde\gamma}{\partial
u},\frac{\partial\tilde\gamma}{\partial t}\right)
=\frac{\partial}{\partial t}
h\left(\frac{\partial\tilde\gamma}{\partial u},
\frac{\partial\tilde\gamma}{\partial t}\right)
-h\left(\frac{\partial\tilde\gamma}{\partial u},
\frac{D}{\partial t}\frac{\partial\tilde\gamma}{\partial t}\right)$
together with the fundamental theorem of calculus. The boundary term
vanishes because $\tilde\gamma(t,u)=\gamma(t)$ for $t$ near $a$ and
$b$, so $\partial\tilde\gamma/\partial u$ vanishes there. Evaluating at
$u=0$,
\begin{equation}\label{eq:first-variation-energy}
E'(0)=-\int_a^bh\left(Y,\frac{D\gamma'}{dt}\right)dt .
\end{equation}

\emph{Realizing a prescribed variation field.} Given a smooth field
$Y$ along $\gamma$ vanishing outside a compact subset of $(a,b)$,
there is a variation of $\gamma$ supported in $(a,b)$ whose variation
field is $Y$. Indeed, choose $\delta>0$ with
$[a-\delta,b+\delta]\subset I$ and set $J=(a-\delta,b+\delta)$. The
geodesic equation is a smooth second-order ODE, so the set of
$(x,v)\in TM$ for which the geodesic with initial condition $(x,v)$
exists up to time $1$ is open, contains the zero section, and on it
$(x,v)\mapsto\exp_x(v)$ is smooth (classical ODE theory, as in the
discussion following Definition~\ref{def:geodesic} and
Definition~\ref{def:exponential-map}). The preimage of this open set
under the continuous map $\Ar^2\supset J\times\Ar\to TM$,
$(t,u)\mapsto\left(\gamma(t),uY(t)\right)$, is open and contains
$[a-\delta/2,b+\delta/2]\times\{0\}$; by compactness of
$[a-\delta/2,b+\delta/2]$ it contains
$[a-\delta/2,b+\delta/2]\times(-\epsilon,\epsilon)$ for some
$\epsilon>0$. Hence
$$\tilde\gamma(t,u)=\exp_{\gamma(t)}\left(uY(t)\right)$$
is defined and smooth on
$(a-\delta/2,b+\delta/2)\times(-\epsilon,\epsilon)$. Since
$\exp_x(0)=x$ we have $\tilde\gamma(t,0)=\gamma(t)$, and
$\tilde\gamma(t,u)=\gamma(t)$ for all $u$ whenever $Y(t)=0$, in
particular outside a compact subset of $(a,b)$. Finally
$\partial\tilde\gamma/\partial u\,(t,0)=Y(t)$, since
$u\mapsto\exp_{\gamma(t)}(uY(t))$ is the geodesic with initial
velocity $Y(t)$ (immediate from the rescaling property of the geodesic
equation, as in the proof of
Lemma~\ref{lem:parallel-gradient-flow}(2)).

\emph{Comparison of the two covariant derivatives.} Now let $\gamma$ be
our curve in $N_c$, let $[a,b]\subset I$, and let $Y$ be a smooth field
along $\gamma$, tangent to $N_c$ at every point and vanishing outside a
compact subset of $(a,b)$; by the last remark in item~1, $Y$ is a
smooth field along $\gamma$ in $N_c$ as well. Applying the previous
paragraph in the Riemannian manifold $(N_c,g_{N_c})$, construct a
variation $\tilde\gamma$ of $\gamma$ supported in $(a,b)$ with
variation field $Y$ and with values in $N_c$. Composing with the
inclusion, $\tilde\gamma$ is also a smooth map into $X$, and as such it
is again a variation of $\gamma$ supported in $(a,b)$ with variation
field $Y$. Since $g_{N_c}$ is the induced metric and each curve
$t\mapsto\tilde\gamma(t,u)$ lies in $N_c$, the integrand
$\left|\partial\tilde\gamma/\partial t\right|^2$ is the same whether
computed with $g_{N_c}$ or with $g$; hence the function $E(u)$ is the
same for the two Riemannian manifolds, and
Equation~(\ref{eq:first-variation-energy}) applied to this one
variation in $(X,g)$ and in $(N_c,g_{N_c})$ yields
$$\int_a^b\left\langle Y,\frac{D\gamma'}{dt}\right\rangle dt
=-E'(0)
=\int_a^bg_{N_c}\left(Y,\frac{D^N\gamma'}{dt}\right)dt
=\int_a^b\left\langle Y,\frac{D^N\gamma'}{dt}\right\rangle dt,$$
the last equality because both entries are tangent to $N_c$ and
$g_{N_c}$ is induced by $g$. Therefore the field
$$F=\frac{D\gamma'}{dt}-\frac{D^N\gamma'}{dt}$$
along $\gamma$ satisfies $\int_a^b\langle Y,F\rangle\,dt=0$ for every
{\em admissible} $Y$, i.e., for every smooth field $Y$ along $\gamma$
which is tangent to $N_c$ at every point and vanishes outside a compact
subset of $(a,b)$. The field $F$ is tangent to $N_c$
(Equation~(\ref{eq:tangential-acceleration}) for the first summand; the
second lies in $TN_c$ by its definition) and smooth. Fix $t_0\in I$,
choose $a<t_0<b$ with $[a,b]\subset I$ and a smooth function
$\varphi\ge0$ with $\varphi(t_0)=1$ vanishing outside a compact subset
of $(a,b)$; then $Y=\varphi F$ is admissible (smooth as a field in
$N_c$ by the slice-chart remark of item~1), and
$$0=\int_a^b\varphi\,|F|^2\,dt$$
with a continuous non-negative integrand forces $\varphi|F|^2\equiv0$;
at $t_0$ this gives $F(t_0)=0$. Since $t_0\in I$ was arbitrary,
$$\frac{D\gamma'}{dt}=\frac{D^N\gamma'}{dt}\qquad\text{along
}\gamma .$$
By Lemma~\ref{lem:covariant-derivative-along-maps}(4), applied in
$(X,g)$ and in $(N_c,g_{N_c})$, the curve $\gamma$ is a geodesic of
$(N_c,g_{N_c})$ if and only if it is a geodesic of $(X,g)$. The
concluding statements of item~3 follow by combining this equivalence
with item~2.
\end{proof}

\begin{proposition}[Splitting along a parallel unit gradient]
\label{prop:parallel-gradient-splitting}
\lean{MorganTianLib.bochnerSplittingHomeomorph,
  MorganTianLib.bochnerSplittingHomeomorphOfComplete,
  MorganTianLib.mfderivReal,
  MorganTianLib.abs_metricInner_le_sqrt_mul_sqrt,
  MorganTianLib.ofReal_abs_sub_le_pathELength,
  MorganTianLib.ofReal_abs_sub_le_edist_of_bochner,
  MorganTianLib.abs_sub_le_dist_of_bochner,
  MorganTianLib.edist_smoothVectorFieldFlow_self_of_bochner,
  MorganTianLib.edist_smoothVectorFieldFlow_flow_of_bochner,
  MorganTianLib.isMinGeodesicOn_smoothVectorFieldFlow_of_bochner,
  MorganTianLib.edist_smoothVectorFieldFlow_pair_le_of_bochner,
  MorganTianLib.ofReal_abs_le_edist_smoothVectorFieldFlow_pair_of_bochner,
  MorganTianLib.edist_le_edist_smoothVectorFieldFlow_pair_add_of_bochner,
  MorganTianLib.isComplete_preimage_of_continuous,
  MorganTianLib.isGeodesicRay_smoothVectorFieldFlow_of_bochner,
  MorganTianLib.busemann_smoothVectorFieldFlow_of_bochner,
  MorganTianLib.sub_le_busemann_smoothVectorFieldFlow_of_bochner,
  MorganTianLib.mfderiv_smoothVectorFieldFlow_comp_apply_one_of_bochner,
  MorganTianLib.enorm_mfderiv_smoothVectorFieldFlow_comp_of_bochner,
  MorganTianLib.pathELength_smoothVectorFieldFlow_comp_of_bochner,
  MorganTianLib.pathELength_smoothVectorFieldFlow_comp_levelSet_of_bochner,
  MorganTianLib.hasDerivAt_comp_gradientField,
  MorganTianLib.enorm_mfderiv_eq_sqrt_metricInner_of_hasMFDerivAt,
  MorganTianLib.enorm_mfderiv_levelProjection_of_bochner,
  MorganTianLib.edist_levelProjection_le_of_bochner,
  MorganTianLib.ofReal_mul_add_mul_le_edist_smoothVectorFieldFlow_pair_of_bochner,
  MorganTianLib.ofReal_sqrt_le_edist_smoothVectorFieldFlow_pair_of_bochner,
  MorganTianLib.continuous_metricInner_mfderiv_of_contMDiff,
  MorganTianLib.pathELength_levelProjection_le_of_bochner,
  MorganTianLib.pathELength_eq_ofReal_integral_of_contMDiff,
  MorganTianLib.edist_smoothVectorFieldFlow_levelPath_le_of_bochner,
  MorganTianLib.edist_smoothVectorFieldFlow_pair_le_ofReal_sqrt_of_bochner,
  MorganTianLib.edist_smoothVectorFieldFlow_pair_eq_of_bochner,
  MorganTianLib.dist_smoothVectorFieldFlow_pair_of_bochner,
  MorganTianLib.busemann_smoothVectorFieldFlow_eq_of_bochner,
  MorganTianLib.prod_dist_eq_sqrt_of_L2,
  MorganTianLib.dist_eq_sqrt_dist_levelProjection_of_bochner,
  MorganTianLib.bochnerSplittingIsometry,
  MorganTianLib.metricInner_flowVariation_of_bochner,
  MorganTianLib.levelSetInducedMetric,
  MorganTianLib.levelSetProductMetric,
  MorganTianLib.DCProductMetric_metricInner_mk,
  MorganTianLib.levelSetProductMetric_metricInner,
  MorganTianLib.levelSetProductMetric_metricInner_ambient,
  MorganTianLib.mfderiv_prod_eq_add,
  MorganTianLib.mfderiv_bochnerSplittingMap_apply,
  MorganTianLib.metricInner_mfderiv_bochnerSplittingMap,
  MorganTianLib.mfderiv_bochnerSplittingMap_injective}
% Formalization status: Step 1 is covered at the topological level
% (bochnerSplittingHomeomorph: M ≃ₜ B⁻¹(0) × ℝ). For Step 4, the metric
% layer is now formalized: B is 1-Lipschitz for the Riemannian distance
% (abs_sub_le_dist_of_bochner, via Cauchy–Schwarz + FTC along C¹ paths,
% ofReal_abs_sub_le_pathELength); each flow line is a unit-speed
% MINIMIZING line, d(θ_s x, θ_t x) = |s − t|
% (edist_smoothVectorFieldFlow_flow_of_bochner,
% isMinGeodesicOn_smoothVectorFieldFlow_of_bochner — so every point of X
% lies on a minimizing line, and the Busemann line lemmas apply to flow
% lines); the level sets are complete (isComplete_preimage_of_continuous);
% and the splitting map satisfies the two-sided product-metric bounds
% d(θ_s x, θ_t y) ≤ d(x,y) + |s−t|, |B x − B y + (s−t)| ≤ d(θ_s x, θ_t y),
% d(x,y) ≤ d(θ_s x, θ_t y) + |s−t|. The flow lines are moreover wired into
% the metric Busemann layer: each is an IsGeodesicRay, its Busemann
% function is −u along the line (busemann_smoothVectorFieldFlow_of_bochner)
% and its Busemann function IS the affine function B_{λ_x} = f(x) − f
% (busemann_smoothVectorFieldFlow_eq_of_bochner: ≥ from the 1-Lipschitz
% bound, ≤ by evaluating the approximants √(D² + (t−a)²) − t ↓ −a with
% the sharp ℓ² product formula).
% The length layer of the sharp ℓ² formula is formalized
% (TiltedPathLength.lean): the tilted-path length formula
% L(θ_{r(u)}(c(u))) = ∫√(r'²c₁ + 2r'⟨∇f,c'⟩ + |c'|²)
% (pathELength_smoothVectorFieldFlow_comp_of_bochner, with the
% Pythagorean level-set form), the level projection x ↦ θ_{v₀−f(x)}(x)
% is 1-Lipschitz (edist_levelProjection_le_of_bochner, via the pointwise
% speed drop enorm_mfderiv_levelProjection_of_bochner |ĉ'|² = |σ'|² −
% ⟨∇f,σ'⟩²), and the SHARP ℓ² LOWER BOUND holds: for x, y in a common
% level set, √(d(x,y)² + (s−t)²) ≤ d(θ_s x, θ_t y)
% (ofReal_sqrt_le_edist_smoothVectorFieldFlow_pair_of_bochner), proven
% by decomposing any competitor σ = θ_ρ(ĉ) along the flow and the level
% projection and integrating the pointwise Cauchy–Schwarz bound
% α|ĉ'| + βρ ≤ |σ'| over each direction (α,β) of the unit disc — no
% integral Minkowski inequality and no measurability of the speed
% needed. The matching ℓ² UPPER bound is also formalized
% (edist_smoothVectorFieldFlow_pair_le_ofReal_sqrt_of_bochner): the
% speed of a C¹ curve is continuous
% (continuous_metricInner_mfderiv_of_contMDiff, via the tangent lift
% into the Riemannian tangent bundle), so a near-minimizing path from x
% to y can be level-projected and tilted through the flow with time
% profile proportional to margin-δ arclength A(u) = ∫(|ĉ'| + δ)
% (edist_smoothVectorFieldFlow_levelPath_le_of_bochner), giving
% d(θ_s x, θ_t y) ≤ √((L+2δ)² + (s−t)²) → √(d(x,y)² + (s−t)²).
% Hence the SHARP ℓ² PRODUCT FORMULA holds as an equality
% (edist_smoothVectorFieldFlow_pair_eq_of_bochner):
% d(θ_s x, θ_t y) = √(d(x,y)² + (s−t)²) for x, y in a common level set —
% the distance of M IS the ℓ² product distance of f⁻¹(v₀) × ℝ under Ψ.
% The METRIC-SPACE SPLITTING is now assembled (SplittingMetric.lean):
% the two-point formula d(x,y) = √(d(π₀x, π₀y)² + (Bx − By)²) holds for
% ALL x, y ∈ X (dist_eq_sqrt_dist_levelProjection_of_bochner, with
% π₀ = θ_{−B(·)}(·) the level projection), and the splitting map Ψ is
% an ISOMETRY M ≃ᵢ B⁻¹(0) ×₂ ℝ onto the ℓ² product WithLp 2 (N × ℝ)
% of the level set (with induced metric) and the line
% (bochnerSplittingIsometry) — the full metric-space Cheeger–Gromoll
% splitting for the Bochner package.
% The RIEMANNIAN FORM Φ*g = g_N ⊕ dt² is formalized in AMBIENT form
% (SplittingRiemannianForm.lean, metricInner_flowVariation_of_bochner):
% for level-tangent vectors v, w ∈ ker dB_x and a, b ∈ ℝ,
% g(dθ_t v + a·∇B, dθ_t w + b·∇B) = g(v, w) + ab at θ_t(x) — the image
% of ((v,a),(w,b)) under dΦ, since the time leg of Φ(y,s) = θ_s(y) is
% the flow velocity ∇B. Combines metric preservation ⟨dθ_t v, dθ_t w⟩ =
% ⟨v,w⟩, flow-invariance of the gradient pairing, and |∇B| = 1.
% The smooth submanifold structure of N (slice charts,
% LevelSetChartedSpace) and its induced Riemannian metric g_N
% (levelSetInducedMetric, the pullback of g through the inclusion, with
% ⟨v,w⟩_{g_N} = ⟨dι v, dι w⟩_g) are now formalized, and the product metric
% g_N ⊕ dt² is DoCarmo's DCProductMetric g_N dt². The product metric now has
% its evaluation formula: for split tangent vectors, ⟨(w,a),(w',b)⟩_{g_N⊕dt²}
% = ⟨w,w'⟩_{g_N} + a·b (levelSetProductMetric_metricInner, via the reusable
% DCProductMetric_metricInner_mk), and its ambient reading ⟨dι w, dι w'⟩_g + a·b
% (levelSetProductMetric_metricInner_ambient) — this is EXACTLY the right-hand
% side produced by metricInner_flowVariation_of_bochner (with v = dι w, which
% satisfies dB(dι w) = 0 since range dι = ker dB). The Riemannian form
% Φ*g = g_N ⊕ dt² is now FULLY FORMALIZED as an actual isometry identity on the
% product manifold N × ℝ (SplittingRiemannianIsometry.lean). The differential of
% the splitting parametrization Φ(y,t) = θ_t(ι y) is computed:
% dΦ_{(y,t)}(w,a) = dθ_t(dι w) + a·V(θ_t ι y)
% (mfderiv_bochnerSplittingMap_apply), assembled by mfderiv_prod_eq_add from the
% joint C¹ regularity of the flow (contMDiff_smoothVectorFieldFlow_uncurry_of_bochner):
% the space leg dθ_t(dι w) is the chain rule for θ_t ∘ ι; the time leg a·V(θ_t ι y)
% is the integral-curve equation isMIntegralCurve_smoothVectorFieldFlow, whose
% differential is ContinuousLinearMap.smulRight 1 (V(θ_t ι y)). Feeding this into
% metricInner_flowVariation_of_bochner (with v = dι w ∈ ker dB) and matching to
% levelSetProductMetric_metricInner_ambient gives the pullback identity
% g(dΦ(w,a), dΦ(w',b)) = ⟨w,w'⟩_{g_N} + a·b, i.e. Φ*g = g_N ⊕ dt²
% (metricInner_mfderiv_bochnerSplittingMap) — the Riemannian Cheeger–Gromoll
% splitting for the Bochner package. As an immediate corollary of the isometry,
% dΦ is INJECTIVE at every point (mfderiv_bochnerSplittingMap_injective): Φ is an
% immersion — the infinitesimal half of Step 1. All three declarations are
% axiom-clean. The remaining unformalized parts of the proposition are the
% global-diffeomorphism smoothness of Ψ = Φ⁻¹ and the connectedness of N.
\uses{def:geodesic,thm:hopf-rinow,def:exponential-map,thm:levi-civita-connection,def:riemannian-metric}
Let $(X,g)$ be a complete, connected Riemannian manifold and let
$B\colon X\to\Ar$ be a smooth function with $\Hess(B)\equiv0$ and
$|\nabla B|\equiv1$. Then $N=B^{-1}(0)$, equipped with the induced
Riemannian metric $g_N$, is a non-empty, connected, smooth, complete,
totally geodesic embedded hypersurface of $X$ (totally geodesic meaning
that every geodesic of $(N,g_N)$ is a geodesic of $(X,g)$), and the map
$$\Phi\colon N\times\Ar\To X,\qquad
\Phi(y,t)=\exp_y\left(t\,(\nabla B)^*(y)\right),$$
where $(\nabla B)^*$ denotes the gradient of $B$, is an isometry from
the Riemannian product $\left(N\times\Ar,\,g_N\oplus dt^2\right)$ onto
$(X,g)$. Here $g_N\oplus dt^2$ is the product metric, defined under the
identification $T_{(y,t)}(N\times\Ar)=T_yN\oplus T_t\Ar$ by
$$\left(g_N\oplus dt^2\right)\left((w,a),(w',a')\right)
=g_N(w,w')+aa' .$$
\end{proposition}

\begin{proof}
\uses{lem:parallel-gradient-flow,lem:parallel-gradient-level-sets,thm:levi-civita-connection,lem:hessian-symmetric,def:riemannian-metric}
Write $V=(\nabla B)^*$ and let $\theta_t$ be the flow of $V$, so that
$\theta_t(x)=\exp_x(tV(x))$, $B\circ\theta_t=B+t$, and each $\theta_t$
is an isometry of $(X,g)$
(Lemma~\ref{lem:parallel-gradient-flow}(2),(3),(4)). By
Lemma~\ref{lem:parallel-gradient-level-sets} with $c=0$, the level set
$N=B^{-1}(0)$ is a non-empty, closed, smooth embedded hypersurface with
$T_yN=\ker(dB_y)=V(y)^\perp$ for every $y\in N$, and $(N,g_N)$ is
totally geodesic in the asserted sense. It remains to prove that $\Phi$
is a diffeomorphism onto $X$ satisfying
$\Phi^*g=g_N\oplus dt^2$, that $N$ is connected, and that $(N,g_N)$ is
complete.

\emph{Step 1: $\Phi$ is a diffeomorphism onto $X$.} By
Lemma~\ref{lem:parallel-gradient-flow}(2),
$\Phi(y,t)=\exp_y(tV(y))=\theta_t(y)$, so $\Phi$ is the restriction to
$N\times\Ar\subset X\times\Ar$ of the smooth map
$(x,t)\mapsto\theta_t(x)$, hence smooth. Define
$$\Psi\colon X\To N\times\Ar,\qquad
\Psi(x)=\left(\theta_{-B(x)}(x),\,B(x)\right).$$
The first component indeed lies in $N$, since
$B\left(\theta_{-B(x)}(x)\right)=B(x)-B(x)=0$ by
Lemma~\ref{lem:parallel-gradient-flow}(3); the map
$x\mapsto\theta_{-B(x)}(x)$ is smooth as a map into $X$ (it is the
composition of $x\mapsto(-B(x),x)$ with the smooth flow), with image in
$N$, hence smooth as a map into $N$ by
Lemma~\ref{lem:parallel-gradient-level-sets}(1). Thus $\Psi$ is smooth.
For $y\in N$ and $t\in\Ar$ we have
$B\left(\Phi(y,t)\right)=B(\theta_t(y))=B(y)+t=t$, so, using the group
law of the flow,
$$\Psi\left(\Phi(y,t)\right)
=\left(\theta_{-t}\left(\theta_t(y)\right),\,t\right)=(y,t),$$
while conversely
$$\Phi\left(\Psi(x)\right)
=\theta_{B(x)}\left(\theta_{-B(x)}(x)\right)=\theta_0(x)=x .$$
Hence $\Phi$ is a diffeomorphism of $N\times\Ar$ onto $X$, with inverse
$\Psi$; note also that $B\circ\Phi$ is the projection onto the $\Ar$
factor, so $\Phi$ carries $N\times\{t\}$ onto $B^{-1}(t)$.

\emph{Step 2: $N$ is connected.} If $N=A\sqcup A'$ with $A,A'$
non-empty, open and disjoint in $N$, then
$N\times\Ar=(A\times\Ar)\sqcup(A'\times\Ar)$ would be a separation of
$N\times\Ar$ into non-empty open sets, and its image under the
diffeomorphism $\Phi$ a separation of $X$, contradicting the
connectedness of $X$. Hence $N$ is connected (and non-empty, by
Lemma~\ref{lem:parallel-gradient-level-sets}(1)).

\emph{Step 3: $\Phi^*g$ is the product metric.} Fix $(y,t)\in
N\times\Ar$ and use the identification
$T_{(y,t)}(N\times\Ar)=T_yN\oplus T_t\Ar$, writing $\partial_t$ for the
standard unit vector of the second factor. First,
$$d\Phi_{(y,t)}(\partial_t)
=\frac{d}{ds}\Bigl|_{s=0}\Bigr.\theta_{t+s}(y)
=V\left(\theta_t(y)\right),$$
since $s\mapsto\theta_{s+t}(y)$ is an integral curve of $V$
(Lemma~\ref{lem:parallel-gradient-flow}(2)). Second, for $w\in T_yN$,
since $\Phi(\cdot,t)=\theta_t|_N$,
$$d\Phi_{(y,t)}(w)=d(\theta_t)_y(w).$$
We compute the three types of inner products. Since $|V|\equiv1$,
$$\left\langle d\Phi(\partial_t),d\Phi(\partial_t)\right\rangle
=\left|V(\theta_t(y))\right|^2=1 .$$
Next, for $w\in T_yN$, using that $V$ is the gradient of $B$ and that
$B\circ\theta_t=B+t$ (Lemma~\ref{lem:parallel-gradient-flow}(3)),
$$\left\langle d\Phi(w),d\Phi(\partial_t)\right\rangle
=\left\langle d(\theta_t)_y(w),V(\theta_t(y))\right\rangle
=dB\left(d(\theta_t)_y(w)\right)
=d(B\circ\theta_t)_y(w)
=dB_y(w)=0,$$
because $w\in T_yN=\ker(dB_y)$. Finally, for $w,w'\in T_yN$, since
$\theta_t$ is an isometry of $(X,g)$
(Lemma~\ref{lem:parallel-gradient-flow}(4)) and $g_N$ is the induced
metric,
$$\left\langle d\Phi(w),d\Phi(w')\right\rangle
=\left\langle d(\theta_t)_y(w),d(\theta_t)_y(w')\right\rangle
=\langle w,w'\rangle=g_N(w,w') .$$
These are exactly the inner products assigned by $g_N\oplus dt^2$ on
$T_yN\oplus T_t\Ar$; hence $\Phi^*g=g_N\oplus dt^2$, and by Step~1,
$\Phi$ is an isometry from the Riemannian product onto $(X,g)$.

\emph{Step 4: $(N,g_N)$ is complete.} Since $\Phi$ is a diffeomorphism
with $\Phi^*g=g_N\oplus dt^2$, it preserves the lengths of smooth
paths:
$$L_g(\Phi\circ\beta)
=\int\left|d\Phi(\dot\beta)\right|_g
=\int|\dot\beta|_{g_N\oplus dt^2}
=L_{g_N\oplus dt^2}(\beta),$$
and since smooth paths between corresponding endpoints correspond
bijectively under $\Phi$, the associated distance functions satisfy
$d_X\left(\Phi(\xi),\Phi(\xi')\right)=d_{N\times\Ar}(\xi,\xi')$ for all
$\xi,\xi'\in N\times\Ar$ ($N\times\Ar$ is connected by Step~2, so these
distances are finite). Now let $y,y'\in N$. Any smooth path $\beta$ in
$N$ from $y$ to $y'$ gives the smooth path $t\mapsto(\beta(t),0)$ in
$N\times\Ar$ from $(y,0)$ to $(y',0)$ of the same length, so
$d_{N\times\Ar}\left((y,0),(y',0)\right)\le d_N(y,y')$. Conversely, if
$(\beta_1,\beta_2)$ is a smooth path in $N\times\Ar$ from $(y,0)$ to
$(y',0)$, then $\beta_1$ is a smooth path in $N$ from $y$ to $y'$ with
$$|\dot\beta_1|_{g_N}
\le\left(|\dot\beta_1|_{g_N}^2+|\dot\beta_2|^2\right)^{1/2},$$
so $d_N(y,y')\le d_{N\times\Ar}\left((y,0),(y',0)\right)$. Hence, using
$\Phi(y,0)=\theta_0(y)=y$,
$$d_N(y,y')=d_{N\times\Ar}\left((y,0),(y',0)\right)=d_X(y,y')
\qquad\text{for all }y,y'\in N:$$
the distance function of the induced metric on $N$ is the restriction
of $d_X$. Let now $(y_k)$ be a Cauchy sequence in $(N,d_N)$. Then
$(y_k)$ is Cauchy in $(X,d_X)$, which is complete by hypothesis, so
$y_k\to x$ for some $x\in X$; since $N$ is closed in $X$
(Lemma~\ref{lem:parallel-gradient-level-sets}(1)), $x\in N$; and
$d_N(y_k,x)=d_X(y_k,x)\to0$. Hence $(N,g_N)$ is complete. This
establishes all the assertions of the proposition.
\end{proof}


\begin{lemma}[Busemann functions along a minimizing line]
\label{lem:busemann-along-line}
\lean{MorganTianLib.busemann_apply_line}
\leanok
\uses{def:busemann-function,def:minimizing-geodesic-window}
\source{morgan-tian}
Let $\lambda\colon\Ar\to M$ be a minimizing geodesic line in a
metric space and let $B_+=B_{\lambda_+}$ be the Busemann function of
its forward ray $\lambda_+(t)=\lambda(t)$, $t\ge0$. Then
$$B_+(\lambda(u))=-u\qquad\text{for all } u\in\Ar,$$
extending Lemma~\ref{lem:busemann-along-ray} from $u\ge0$ to all
real $u$. Applied to the reversed line $u\mapsto\lambda(-u)$, which
is again minimizing, this also gives $B_-(\lambda(u))=u$ for the
Busemann function $B_-$ of the backward ray
$\lambda_-(t)=\lambda(-t)$.
\end{lemma}

\begin{proof}
\leanok
\uses{lem:busemann-le-approx}
For $t\ge\max(u,0)$ minimality gives
$d(\lambda(t),\lambda(u))=t-u$, so
$B_{\lambda_+,t}(\lambda(u))=(t-u)-t=-u$; in particular, taking
$t=\max(u,0)$ and using Lemma~\ref{lem:busemann-le-approx},
$B_+(\lambda(u))\le-u$. Conversely, for every $t\ge0$,
$$B_{\lambda_+,t}(\lambda(u))=|t-u|-t\ge(t-u)-t=-u,$$
so the infimum over $t\ge0$ is at least $-u$. Hence
$B_+(\lambda(u))=-u$. The statement for $B_-$ follows by applying
the result to the reversed line at the parameter $-u$.
\end{proof}

\begin{lemma}[The Busemann pair of a line is non-negative]
\label{lem:busemann-pair-nonneg}
\lean{MorganTianLib.busemann_add_busemann_neg_nonneg}
\leanok
\uses{def:busemann-function,def:minimizing-geodesic-window}
\source{morgan-tian}
Let $\lambda\colon\Ar\to M$ be a minimizing geodesic line in a
metric space, with Busemann functions $B_\pm$ of its forward and
backward rays $\lambda_\pm(t)=\lambda(\pm t)$. Then
$$B_+(x)+B_-(x)\ge0\qquad\text{for all } x\in M.$$
\end{lemma}

\begin{proof}
\leanok
\uses{def:busemann-approx}
Fix $x\in M$ and $s,t\ge0$. Since $\lambda$ is minimizing,
$d(\lambda(t),\lambda(-s))=s+t$, and the triangle inequality gives
$$d(x,\lambda(t))+d(x,\lambda(-s))\ge d(\lambda(t),\lambda(-s))
=s+t,$$
that is,
$$B_{\lambda_+,t}(x)+B_{\lambda_-,s}(x)
=\bigl(d(x,\lambda(t))-t\bigr)+\bigl(d(x,\lambda(-s))-s\bigr)
\ge0.$$
Taking the infimum over $t\ge0$ and then over $s\ge0$ yields
$B_+(x)+B_-(x)\ge0$.
\end{proof}

\begin{corollary}[The Busemann pair vanishes on its line]
\label{cor:busemann-pair-vanishes-on-line}
\lean{MorganTianLib.busemann_add_busemann_neg_apply_line}
\leanok
\uses{lem:busemann-along-line}
In the situation of Lemma~\ref{lem:busemann-pair-nonneg}, for every
$u\in\Ar$,
$$B_+(\lambda(u))+B_-(\lambda(u))=0.$$
\end{corollary}

\begin{proof}
\leanok
\uses{lem:busemann-along-line}
Immediate from $B_+(\lambda(u))=-u$ and $B_-(\lambda(u))=u$
(Lemma~\ref{lem:busemann-along-line}).
\end{proof}

\begin{lemma}[Minimizing line implies isometric product]
\label{lem:minimizing-line-implies-product}
\label{line}
\uses{def:ricci-curvature}
\source{morgan-tian}
(Cheeger--Gromoll; this result was formulated as the main theorem in
\cite{CheegerGromoll2}.)
Any connected, complete Riemannian manifold $X$ of non-negative Ricci
curvature containing a minimizing line is isometric to a product
$N\times \Ar$ for some Riemannian manifold $N$.
\end{lemma}

\begin{proof}
\uses{def:busemann-function,rem:laplacian-weak-sense,prop:busemann-laplacian,thm:maximum-principle,thm:weak-harmonic-regularity,def:asymptotic-ray,lem:asymptotic-ray-exists,lem:asymptotic-ray,lem:busemann-gradient-norm-one,lem:function-bochner-formula,prop:parallel-gradient-splitting,lem:busemann-along-line,lem:busemann-pair-nonneg,cor:busemann-pair-vanishes-on-line}
Given a minimizing line $\lambda\colon \Ar\to X$, define
$\lambda_\pm\colon [0,\infty)\to X$ by $\lambda_+(t)=\lambda(t)$ and
$\lambda_-(t)=\lambda(-t)$. These are minimizing geodesic rays, and we have
the Busemann functions $B_+=B_{\lambda_+}$ and $B_-=B_{\lambda_-}$ of
Definition~\ref{def:busemann-function}; each is Lipschitz with Lipschitz
constant $1$, and in particular continuous.

Proposition~\ref{prop:busemann-laplacian} applies to both $B_+$ and $B_-$
and shows that $\triangle B_+\le 0$ and $\triangle B_-\le 0$ in the weak sense;
that is (Remark~\ref{rem:laplacian-weak-sense}), for every non-negative
test function $\varphi$ (i.e., every non-negative, compactly supported
$C^\infty$-function $\varphi$ on $X$) we have
$$\int_X B_+\triangle\varphi\,d\vol\le 0
\qquad\text{and}\qquad
\int_X B_-\triangle\varphi\,d\vol\le 0.$$
Adding these two inequalities for one and the same non-negative test
function $\varphi$ gives
$$\int_X (B_++B_-)\triangle\varphi\,d\vol\le 0,$$
that is, $\triangle(B_++B_-)\le 0$ in the weak sense.

On the other hand, since $\lambda$ is distance minimizing,
Lemma~\ref{lem:busemann-pair-nonneg} gives $B_++B_-\ge 0$
everywhere, and
Corollary~\ref{cor:busemann-pair-vanishes-on-line} gives
$B_+(\lambda(u))+B_-(\lambda(u))=0$ for every point $\lambda(u)$ in
the image of $\lambda$; recall also from
Lemma~\ref{lem:busemann-along-line} the individual values
$B_+(\lambda(u))=-u$ and $B_-(\lambda(u))=u$.

Thus the continuous function $f=-(B_++B_-)$ satisfies $f\le 0$ everywhere
and $f(\lambda(0))=0$, so that $f$ achieves its maximum value $0$; and
$\triangle f\ge 0$ in the weak sense, since for every non-negative test
function $\varphi$
$$\int_X f\,\triangle\varphi\,d\vol
=-\int_X (B_++B_-)\triangle\varphi\,d\vol\ge 0.$$
Since $X$ is connected, the maximum principle
(Theorem~\ref{thm:maximum-principle}) applies to $f$ and shows that $f$ is
constant, equal to its maximum value $0$. That is, $B_-=-B_+$ on all of
$X$.

It now follows that $\triangle B_+=0$ in the weak sense. Indeed, for every
non-negative test function $\varphi$ we have the two inequalities
$$\int_X B_+\triangle\varphi\,d\vol\le 0
\qquad\text{and}\qquad
\int_X B_+\triangle\varphi\,d\vol
=-\int_X B_-\triangle\varphi\,d\vol\ge 0,$$
the first being Proposition~\ref{prop:busemann-laplacian} for $B_+$ and the
second being Proposition~\ref{prop:busemann-laplacian} for $B_-$ combined
with $B_+=-B_-$. Hence $\int_X B_+\triangle\varphi\,d\vol=0$ for every
non-negative test function $\varphi$, and therefore for every test function
$\varphi$: choosing a compactly supported $C^\infty$-function
$\chi\colon X\to[0,1]$ with $\chi\equiv 1$ on $\supp\varphi$ and setting
$\psi=\bigl(\sup_X|\varphi|\bigr)\chi$, both $\psi$ and $\varphi+\psi$ are
non-negative test functions and $\varphi=(\varphi+\psi)-\psi$. By
Theorem~\ref{thm:weak-harmonic-regularity} the continuous function $B_+$
is smooth and harmonic: $\triangle B_+=0$ in the classical sense.

Next, we show that $|\nabla B_+(x)|=1$ for all $x\in X$. Fix $x\in X$.
Since $X$ is complete, Lemma~\ref{lem:asymptotic-ray} provides a minimizing
geodesic ray $\mu$ emanating from $x$ asymptotic to the minimizing ray
$\lambda_+$. The function $B_+$ is smooth, hence in particular
differentiable at $x$, so Lemma~\ref{lem:busemann-gradient-norm-one} gives
$\nabla B_+(x)=-\mu'(0)$; as $\mu$ is a unit-speed geodesic, it follows
that $|\nabla B_+(x)|=1$.

Now apply the Bochner formula for functions
(Lemma~\ref{lem:function-bochner-formula}) to the smooth function $B_+$:
$$\frac{1}{2}\triangle|\nabla B_+|^2
=|\Hess B_+|^2+\langle\nabla B_+,\nabla(\triangle B_+)\rangle
+\Ric(\nabla B_+,\nabla B_+).$$
The left-hand side vanishes because $|\nabla B_+|^2\equiv 1$, and the
middle term on the right vanishes because $\triangle B_+\equiv 0$. Hence
$$0=|\Hess B_+|^2+\Ric(\nabla B_+,\nabla B_+)\ge|\Hess B_+|^2,$$
because $\Ric\ge 0$. Therefore $\Hess B_+\equiv 0$.

Thus $B_+$ is a smooth function on the connected, complete manifold $X$
with $\Hess B_+\equiv 0$ and $|\nabla B_+|\equiv 1$.
Proposition~\ref{prop:parallel-gradient-splitting} now shows that
$N=B_+^{-1}(0)$, with the induced metric, is a smooth, complete, totally
geodesic hypersurface of $X$ --- non-empty, since $B_+(\lambda(0))=0$ ---
and that $(y,t)\mapsto \exp_y\bigl(t\,\nabla B_+(y)\bigr)$ is an isometry
from the Riemannian product $N\times\Ar$ onto $X$. Hence $X$ is isometric
to a product $N\times\Ar$, as claimed.
\end{proof}

\begin{theorem}[Splitting Theorem]
\label{thm:splitting-theorem}
\label{splitting}
\uses{def:ricci-curvature,def:space-of-ends}
\source{morgan-tian}
Suppose that $M$ is complete, of non-negative Ricci curvature and suppose
that $M$ has at least two ends. Then $M$ is isometric to a product
$N\times \Ar$ where $N$ is a compact manifold.
\end{theorem}

\begin{proof}
\uses{lem:ends-exist,lem:minimizing-line-implies-product,def:space-of-ends,lem:compact-set-in-codim-zero-submanifold}
The hypothesis that $M$ has at least two ends presupposes, via
Definition~\ref{def:space-of-ends}, that $M$ is connected. Also, $M$ is
non-compact: were $M$ compact, then $M$ itself would be one of the
compact, codimension-$0$ submanifolds indexing the inverse system of
Definition~\ref{def:space-of-ends}, and the corresponding term
$\pi_0(M\setminus M)=\pi_0(\emptyset)$ of that system would be empty; an
end selects an element of every term of the system, so $M$ would have no
ends at all, contrary to hypothesis. Thus $M$ is complete, connected,
non-compact and has more than one end, so by Lemma~\ref{lem:ends-exist}
it contains a minimizing line.

By Lemma~\ref{lem:minimizing-line-implies-product} there are a Riemannian
manifold $N$ and an isometry $\Phi\colon N\times\Ar\to M$. The manifold
$N$ is connected, being the image of the connected space $M$ under the
continuous map $\mathrm{pr}_N\circ\Phi^{-1}$. It remains to show that $N$
is compact. Suppose to the contrary that $N$ is non-compact, and set
$P=N\times\Ar$.

We first note that $P$ has at least two ends. The isometry $\Phi$ is in
particular a diffeomorphism. Hence $K\mapsto\Phi(K)$ is a bijection from
the set of compact, codimension-$0$ submanifolds of $P$ onto the set of
those of $M$, and for each such $K$ the homeomorphism
$P\setminus K\to M\setminus\Phi(K)$ induced by $\Phi$ gives a bijection
$\pi_0(P\setminus K)\to\pi_0(M\setminus\Phi(K))$. These bijections commute
with the maps of the two inverse systems of
Definition~\ref{def:space-of-ends}, all of which are induced by inclusions
of complements, and hence they induce a bijection from the space of ends
of $P$ onto the space of ends of $M$. Since $M$ has at least two ends, so
does $P$.

We now derive a contradiction by showing that, because $N$ is connected
and non-compact, $P=N\times\Ar$ has at most one end. We proceed in two
steps; we shall also use the fact that every compact subset of a smooth
manifold is contained in a compact, codimension-$0$ submanifold
(Lemma~\ref{lem:compact-set-in-codim-zero-submanifold}).

\textit{Step 1: for every compact subset $K\subset P=N\times\Ar$ there is
a compact set $W$ with $K\subseteq W\subset P$ such that $P\setminus W$ is
path-connected.} Let $p_N\colon P\to N$ and $p_{\Ar}\colon P\to\Ar$ be the
projections. The sets $C=p_N(K)\subset N$ and $p_{\Ar}(K)\subset\Ar$ are
compact, being continuous images of the compact set $K$; choose $a>0$ with
$p_{\Ar}(K)\subseteq[-a,a]$, which is possible since a compact subset of
$\Ar$ is bounded, and set $W=C\times[-a,a]$. Then $W$ is compact, and
$K\subseteq W$ since every $(y,t)\in K$ has $y\in p_N(K)=C$ and
$t\in p_{\Ar}(K)\subseteq[-a,a]$. Since $N$ is non-compact and $C$ is
compact, there is a point $z\in N\setminus C$; set $q=(z,a+1)$, a point of
$P\setminus W$. Being a connected manifold, $N$ is locally path-connected
and connected, hence path-connected. We claim that every point
$(y,t)\in P\setminus W$ can be joined to $q$ by a path in $P\setminus W$;
there are three cases.
\begin{enumerate}
\item If $t>a$: the vertical segment $s\mapsto\bigl(y,(1-s)t+s(a+1)\bigr)$,
$s\in[0,1]$, has second coordinate lying between $t$ and $a+1$, hence
$>a$ throughout, so it avoids $W$ (every point of $W$ has second
coordinate $\le a$) and joins $(y,t)$ to $(y,a+1)$. Next choose a path
$\alpha$ in $N$ from $y$ to $z$; then $s\mapsto(\alpha(s),a+1)$ lies in
$N\times\{a+1\}$, which is disjoint from $W$, and joins $(y,a+1)$ to $q$.
\item If $|t|\le a$: then $y\notin C$, since $(y,t)\notin W=C\times[-a,a]$.
Hence the vertical segment from $(y,t)$ to $(y,a+1)$ lies in
$\{y\}\times\Ar$, which is disjoint from $W$; from $(y,a+1)$ proceed to
$q$ as in the first case.
\item If $t<-a$: the vertical segment from $(y,t)$ to $(y,-(a+1))$ has
second coordinate $<-a$ throughout and so avoids $W$. Choose a path
$\alpha$ in $N$ from $y$ to $z$; then $s\mapsto(\alpha(s),-(a+1))$ lies in
$N\times\{-(a+1)\}$, disjoint from $W$, and joins $(y,-(a+1))$ to
$(z,-(a+1))$. Finally, the vertical segment from $(z,-(a+1))$ to
$(z,a+1)=q$ lies in $\{z\}\times\Ar$, which is disjoint from $W$ since
$z\notin C$.
\end{enumerate}
Hence any two points of $P\setminus W$ can be joined by a path through
$q$, and $P\setminus W$ is path-connected, in particular connected. This
proves Step 2.

\textit{Step 2: $P$ has at most one end.} Let ${\mathcal E}$ and
${\mathcal E}'$ be ends of $P$. By Definition~\ref{def:space-of-ends}
they are elements of the inverse limit of the system
$K\mapsto\pi_0(P\setminus K)$: for each compact, codimension-$0$
submanifold $K\subset P$, the end ${\mathcal E}$ selects a component
${\mathcal E}_K\in\pi_0(P\setminus K)$, and whenever $K\subseteq L$ the
map of the system $\pi_0(P\setminus L)\to\pi_0(P\setminus K)$ --- which
sends a component of $P\setminus L$ to the unique component of
$P\setminus K$ containing it --- sends ${\mathcal E}_L$ to
${\mathcal E}_K$; equivalently, ${\mathcal E}_L\subseteq{\mathcal E}_K$.
Now fix a compact, codimension-$0$ submanifold $K\subset P$. By Step 1
there is a compact set $W\supseteq K$ such that $P\setminus W$ is
connected, and by
Lemma~\ref{lem:compact-set-in-codim-zero-submanifold} there is a
compact, codimension-$0$ submanifold
$L\subset P$ with $W\subseteq L$. The connected set $P\setminus W$ is
contained in $P\setminus K$, hence lies in a single component
$V\in\pi_0(P\setminus K)$. The component ${\mathcal E}_L$ of
$P\setminus L$ is non-empty and contained in
$P\setminus L\subseteq P\setminus W$, so ${\mathcal E}_L\subseteq V$; on
the other hand $K\subseteq L$, so ${\mathcal E}_L\subseteq{\mathcal E}_K$.
Thus ${\mathcal E}_K$ and $V$ are components of $P\setminus K$ with
${\mathcal E}_K\cap V\supseteq{\mathcal E}_L\ne\emptyset$, whence
${\mathcal E}_K=V$. The identical argument gives ${\mathcal E}'_K=V$, so
${\mathcal E}_K={\mathcal E}'_K$. Since $K$ was arbitrary and two elements
of an inverse limit are equal exactly when all of their coordinates agree,
${\mathcal E}={\mathcal E}'$. This proves Step 2.

Step 2 contradicts the fact, established above, that $P$ has at least two
ends. Hence $N$ is compact, and the isometry $\Phi\colon N\times\Ar\to M$
exhibits $M$ as a product $N\times\Ar$ with $N$ compact. This completes
the proof of the theorem.
\end{proof}

\section{$\epsilon$-necks}

Certain types of (incomplete) Riemannian manifolds play an
especially important role in our analysis. The purpose of this
section is to introduce these manifolds and use them to prove one
essential result in Riemannian geometry.



For all of the following definitions we fix $0<\epsilon<1/2$. Set
$k$ equal to the greatest integer less than or equal to
$\epsilon^{-1}$. In particular, $k\ge 2$.

\begin{definition}[Epsilon-closeness in $C^{[1/\epsilon]}$-topology]
\label{def:epsilon-close-topology}
\label{epsclose}
\uses{thm:levi-civita-connection}
Suppose that we have a fixed metric $g_0$ on a manifold $M$ and an
open submanifold $X\subset M$. We say that another metric $g$ on $X$
is {\em within $\epsilon$ of $g_0|_X$ in the
$C^{[1/\epsilon]}$-topology} if, setting $k=[1/\epsilon]$ we have
\begin{equation}\label{eqnepsclose}\sup_{x\in X}\left(|g(x)-g_0(x)|_{g_0}^2+
\sum_{\ell=1}^k|\nabla_{g_0}^\ell g(x)|_{g_0}^2\right)<
\epsilon^2,\end{equation} where the covariant derivative
$\nabla^\ell_{g_0}$ is the Levi-Civita connection of $g_0$ and
norms are the pointwise $g_0$-norms on
$$\Sym^2T^*M\otimes \underbrace{T^*M\otimes\cdots\otimes
T^*M}_{\ell-\mathrm{times}}.$$

More generally, given two smooth families of metrics $g(t)$ and
$g_0(t)$ on $M$ defined for $t$ in some interval $I$ we say that
the family $g(t)|_X$ is within $\epsilon$ of the family $g_0(t)|_X$
in the $C^{[1/\epsilon]}$-topology if we have
$$\sup_{(x,t)\in X\times I}\left(\left| g(x,t)
-g_0(x,t)\right|_{g_0(t)}^2 +\sum_{\ell=1}^k\left|\nabla_{g_0}^\ell
g(x,t)\right|_{g_0}^2\right)<\epsilon^2.$$
\end{definition}

\begin{remark}[Ricci flow and metric topology]
\label{rem:ricci-flow-metric-topology}
\uses{def:epsilon-close-topology}
Notice that if we view a one-parameter family of metrics $g(t)$ as a
curve in the space of metrics on $X$ with the
$C^{[1/\epsilon]}$-topology then this is the statement that the two
paths are pointwise within $\epsilon$ of each other. It says nothing
about the derivatives of the paths, or equivalently about the time
derivatives of the metrics and of their covariant derivatives. We
will always be considering paths of metrics satisfying the Ricci
flow equation. In this context two one-parameter families of metrics
that are close in the $C^{2k}$-topology exactly when the $r^{th}$
time derivatives of the $s^{th}$-covariant derivatives are close for
all $r,s$ with $s+2r\le 2k$.
\end{remark}

The first object of interest is one that, up to scale, is close to a
long, round cylinder.


\begin{definition}[Epsilon-neck structure]
\label{def:epsilon-neck}
\label{epsneck}
\uses{def:sectional-curvature,def:epsilon-close-topology}
Let $(N,g)$ be a Riemannian manifold and $x\in N$ a point. Then {\em
an $\epsilon$-neck structure on $(N,g)$
centered at $x$} consists of a diffeomorphism
$$\varphi\colon S^2\times (-\epsilon^{-1},\epsilon^{-1})\to N,$$
with $x\in \varphi(S^2\times\{0\})$, such that the metric
$R(x)\varphi^*g$ is within $\epsilon$ in the
$C^{[1/\epsilon]}$-topology of the product of the usual Euclidean
metric on the open interval with the metric of constant Gaussian
curvature $1/2$ on $S^2$. We also use the terminology {\em $N$ is an
$\epsilon$-neck centered at $x$}. The image under $\varphi$ of the
family of submanifolds $S^2\times \{t\}$ is called the {\em family
of $2$-spheres of the $\epsilon$-neck}. The submanifold
$\varphi(S^2\times \{0\})$ is called {\em the central $2$-sphere} of
the $\epsilon$-neck structure. We denote by $s_N\colon N\to \Ar$
the composition $p_2\circ \varphi^{-1}$, where $p_2$ is the
projection of $S^2\times (-\epsilon^{-1},\epsilon^{-1})$ to the
second factor. There is also the vector field $\partial/\partial
s_N$ on $N$ which is $\varphi_*$ of the standard vector field in the
interval-direction of the product. We also use the terminology of
the {\em plus} and {\em minus} end of the $\epsilon$-neck in the
obvious sense. The opposite (or reversed) $\epsilon$-neck structure
is the one obtained by composing the structure map with $\mathrm{Id}_{S^2}\times -1$. We define the
{\em positive half of the neck} to be the region
$s_N^{-1}(0,\epsilon^{-1})$ and the {\em negative half} to be the
region $s_N^{-1} (-\epsilon^{-1},0)$. For any other fraction, e.g.,
the left-hand three-quarters, the right-hand one-quarter, there are
analogous notions, all measured with respect to $s_N\colon N\to
(-\epsilon^{-1},\epsilon^{-1})$. We also use the terminology the
middle one-half, or middle one-third of the $\epsilon$-neck; again
these regions have their obvious meaning when measured via $s_N$.


{\em An $\epsilon$-neck} in a Riemannian manifold $X$ is a
codimension-zero submanifold $N$ and an $\epsilon$-structure on $N$
centered at some point $x\in N$.

The {\em scale} of an $\epsilon$-neck $N$ centered
at $x$ is $R(x)^{-1/2}$. The scale of $N$ is denoted $r_N$. Intuitively, this
is a measure of the radius of the cross-sectional $S^2$ in the neck. In fact,
the extrinsic diameter of any $S^2$ factor in the neck is close to $\sqrt{2}\pi
r_N$. See the figure in the introduction.
\end{definition}



Here is the result that will be so important in our later arguments.

\begin{proposition}[Positively curved manifolds don't have small epsilon-necks]
\label{prop:epsilon-neck-narrow}
\label{narrows}
\uses{def:epsilon-neck,def:sectional-curvature}
The following
holds for any $\epsilon>0$ sufficiently small.
Let $(M,g)$ be a complete, positively curved Riemannian
$3$-manifold. Then $(M,g)$ does not contain $\epsilon$-necks of
arbitrarily small scale.
\end{proposition}



\begin{proof}
\uses{thm:soul-theorem,prop:busemann-laplacian,lem:neck-geodesic-directions}
The result is obvious if $M$ is compact, so we assume that $M$ is
non-compact. Let $p\in M$ be a soul for $M$ (Theorem~\ref{thm:soul-theorem}),
and let $f$ be the distance function from $p$. Then $f^{-1}(s)$ is
connected for all $s>0$.


\begin{lemma}[Epsilon-neck separates soul from end]
\label{lem:epsilon-neck-separates}
\label{neckseparate}
\uses{def:epsilon-neck,lem:metric-sphere-meets-neck-topologically}
Suppose that $\epsilon>0$ is sufficiently small that Lemma~\ref{topiso} from
the appendix holds. Let $(M,g)$ be a non-compact $3$-manifold of positive
curvature and let $p\in M$ be a soul for it. Then for any $\epsilon$-neck $N$
disjoint from $p$
the central $2$-sphere of $N$ separates the soul from the end of the manifold. In particular, if
two $\epsilon$-necks $N_1$ and $N_2$ in $M$ are disjoint from each other and from $p$, then
the central $2$-spheres of $N_1$ and $N_2$ are the boundary
components of a region in $M$ diffeomorphic to $S^2\times I$.
\end{lemma}

\begin{proof}
\uses{lem:metric-sphere-meets-neck-topologically}
Let $N$ be an $\epsilon$-neck disjoint from $p$. By
Lemma~\ref{topiso} for any point $z$ in the middle third of $N$, the
boundary of the metric ball $B(p,d(z,p))$ is a topological
$2$-sphere in $N$ isotopic in $N$ to the central $2$-sphere of $N$.
Hence, the central $2$-sphere separates the soul from the end of
$M$. The second statement follows immediately by applying this to
$N_1$ and $N_2$.
\end{proof}

Let $N_1$ and $N_2$ be disjoint $\epsilon$-necks, each disjoint from
the soul. By the previous lemma, the central $2$-spheres $S_1$ and
$S_2$ of these necks are smoothly isotopic to each other and they
are the boundary components of a region diffeomorphic to $S^2\times
I$. Reversing the indices if necessary we can assume that $N_2$ is
closer to $\infty$ than $N_1$, i.e., further from the soul.
Reversing the directions of the necks if necessary, we can arrange
that for $i=1,2$ the function $s_{N_i}$ is increasing as we go away
from the soul. We define $C^\infty$- functions $\psi_i$ on $N_i$,
functions depending only on $s_{N_i}$, as follows. The function
$\psi_1$ is zero on the negative side of the middle third of $N_1$
and increases to be identically one on the positive side of the
middle third. The function $\psi_2$ is one on the negative side of
the middle third of $N_2$ and decreases to be zero on the positive
side. We extend $\psi_1,\psi_2$ to a function $\psi$ defined on all
of $M$ by requiring that it be identically one on the region $X$
between $N_1$ and $N_2$ and to be identically zero on $M\setminus
(N_1\cup X\cup N_2)$.

Let $\lambda$ be a geodesic ray from the soul of $M$ to infinity,
and $B_\lambda$ its Busemann function. Let $N$ be any
$\epsilon$-neck disjoint from the soul, with $s_N$ direction chosen
so that it points away from the soul. At any point of the middle
third of $N$ where $B_\lambda$ is smooth, $\nabla B_\lambda$ is a
unit vector in the direction of the unique minimal geodesic ray from
the end of $\lambda$ to this point. Invoking Lemma~\ref{directions}
from the appendix we see that at such points $\nabla B_\lambda$ is
close to $-R(x)^{1/2}\partial/\partial s_N$, where $x\in N$ is the
center of the $\epsilon$-neck. Since $\nabla B_\lambda$ is $L^2$ its
non-smooth points have measure zero and hence, the restriction of
$\nabla B_\lambda$ to the middle third of $N$ is close in the
$L^2$-sense to $-R(x)^{1/2}\partial/\partial s_N$.

Applying this to $N_1$ and $N_2$ we see that
\begin{equation}\label{locallabel}
\int_M\langle \nabla B_\lambda,\nabla\psi\rangle d\vol=\left(\alpha_2R(x_2)^{-1}-\alpha_1R(x_1)^{-1}\right)\Vol_{h_0}(S^2)),\end{equation} where $h(0)$ is the round metric of
scalar curvature $1$ and where each of $\alpha_1$ and $\alpha_2$
limits to $1$ as $\epsilon$ goes to zero. Since $\psi\ge 0$,
Proposition~\ref{prop:busemann-laplacian} tells us that the left-hand side of
Equation~(\ref{locallabel}) must be $\ge 0$. This shows that,
provided that $\epsilon$ is sufficiently small, $R(x_2)$ is bounded
above by $2R(x_1)$. This completes the proof of the proposition.
\end{proof}


\begin{corollary}[Scale bound between separated $\epsilon$-necks]
\label{cor:epsilon-neck-scale-bound}
\label{strnarrows}
\uses{def:epsilon-neck,thm:soul-theorem,lem:epsilon-neck-separates,lem:metric-sphere-meets-neck-topologically}
Fix $\epsilon>0$ sufficiently small so that Lemma~\ref{topiso} holds. Then
there is a constant $C<\infty$ depending on $\epsilon$ such that the following
holds. Suppose that $M$ is a non-compact $3$-manifold of positive sectional
curvature. Suppose that $N$ is an $\epsilon$-neck in $M$ centered at a point
$x$ and disjoint from a soul $p$ of $M$. Then for any $\epsilon$-neck $N'$ that
is separated from $p$ by $N$ with center $x'$ we have $R(x')\le CR(x)$.
\end{corollary}

\section{Forward difference quotients}

Let us review quickly some standard material on forward difference quotients.

\begin{definition}[Forward difference quotient bounds]
\label{def:forward-difference-quotient}
\lean{MorganTianLib.ForwardDiffQuotientLE}
\leanok
\source{morgan-tian}
Let $f\colon [a,b]\to \Ar$ be a continuous function on an interval.
We say that the {\em forward difference quotient of $f$ at a point
$t\in [a,b)$, denoted $\frac{df}{dt}(t)$, is less than or equal to
$c$} provided that
$$\overline{\lim}_{\triangle t\rightarrow
0^+}\frac{f(t+\triangle t)-f(t)}{\triangle t}\le c,$$
that is, provided that for every $r>c$ there is $\delta>0$ such that
$\frac{f(t+\triangle t)-f(t)}{\triangle t}<r$ for all $\triangle t\in
(0,\delta)$. (The second formulation is the meaning we attach to the
upper limit condition; it is robust even when the difference quotients
are unbounded near $t$.)
\end{definition}

Symmetrically, we say that the forward difference quotient of $f$ at
$t$ is greater than or equal to $c'$ if $$c'\le
\underline{\lim}_{\triangle t\rightarrow 0^+}\frac{f(t+\triangle
t)-f(t)}{\triangle t}.$$ Two immediate consequences of the
definition:

\begin{lemma}[Weakening a forward difference quotient bound]
\label{lem:forward-difference-quotient-mono}
\lean{MorganTianLib.ForwardDiffQuotientLE.mono}
\leanok
\uses{def:forward-difference-quotient}
If the forward difference quotient of $f$ at $t$ is at most $c$ and
$c\le c'$, then it is at most $c'$.
\end{lemma}

\begin{proof}
\leanok
Immediate from Definition~\ref{def:forward-difference-quotient}: any
$r>c'$ also satisfies $r>c$.
\end{proof}

\begin{lemma}[Right derivatives bound forward difference quotients]
\label{lem:right-derivative-forward-difference-quotient}
\lean{MorganTianLib.HasDerivWithinAt.forwardDiffQuotientLE}
\leanok
\uses{def:forward-difference-quotient}
If $f$ is differentiable from the right at $t$ with right derivative
$c$, then the forward difference quotient of $f$ at $t$ is at most
$c$.
\end{lemma}

\begin{proof}
\leanok
The difference quotients $\frac{f(t+\triangle t)-f(t)}{\triangle t}$
converge to $c$ as $\triangle t\rightarrow 0^+$; in particular, for
every $r>c$ they are eventually less than $r$.
\end{proof}

Standard comparison arguments show:

\begin{lemma}[Forward difference quotient comparison lemma]
\label{lem:forward-difference-quotient}
\label{fordiffquot}
\lean{MorganTianLib.le_of_forwardDiffQuotientLE}
\leanok
\uses{def:forward-difference-quotient}
\source{morgan-tian}
Suppose that $f\colon [a,b]\to \Ar$ is a continuous function. Suppose that
$\psi$ is a $C^1$-function on $[a,b]\times \Ar$ and suppose that
$\frac{df}{dt}(t)\le \psi(t,f(t))$ for every $t\in [a,b)$ in the sense of
forward difference quotients. Suppose also that there is a function $G(t)$
defined on $[a,b]$ that satisfies the differential equation
$G'(t)=\psi(t,G(t))$ and has $f(a)\le G(a)$. Then $f(t)\le G(t)$ for all $t\in
[a,b]$.
\end{lemma}

\begin{proof}
\uses{def:forward-difference-quotient}
\leanok
Throughout, of $G$ we only use that it is continuous on $[a,b]$ and
satisfies $G'(t)=\psi(t,G(t))$ in the sense of right derivatives at
every $t\in[a,b)$.

Since $f$ and $G$ are continuous on the compact interval $[a,b]$,
there is $C<\infty$ with $|f(t)|\le C$ and $|G(t)|\le C$ for all
$t\in[a,b]$. Set
$$S=[a,b]\times[-(C+1),C+1].$$
Since $\psi$ is $C^1$ on $[a,b]\times\Ar$, its partial derivative in
the second variable is continuous, hence bounded on the compact set
$S$; by the mean value theorem in the second variable (the segment
from $(t,y)$ to $(t,y')$ stays in the convex set $S$) we may
therefore fix $K\ge 1$ such that
$$|\psi(t,y)-\psi(t,y')|\le K|y-y'|\qquad\text{whenever }t\in[a,b],\
y,y'\in[-(C+1),C+1].$$

Fix any $\epsilon>0$ small enough that $\epsilon e^{2K(b-a)}\le 1$,
and define the barrier
$$B_\epsilon(t)=G(t)+\epsilon e^{2K(t-a)},\qquad t\in[a,b].$$
Then $B_\epsilon$ is continuous on $[a,b]$, satisfies
$B_\epsilon>G$ pointwise, and has right derivative
$$B_\epsilon'(t)=\psi(t,G(t))+2K\epsilon e^{2K(t-a)}$$
at every $t\in[a,b)$.

{\em Claim: $f\le B_\epsilon$ on $[a,b]$.} Let
$$T=\{t\in[a,b]\ :\ f\le B_\epsilon\ \text{on}\ [a,t]\},\qquad
t^*=\sup T.$$
Since $f(a)\le G(a)<B_\epsilon(a)$ we have $a\in T$; $T$ is a
subinterval of $[a,b]$ containing $a$, and by continuity of $f$ and
$B_\epsilon$ we have $f\le B_\epsilon$ on $[a,t^*]$, so $t^*\in T$.
Suppose for contradiction that $t^*<b$. We produce $\eta>0$ with
$t^*+\eta\le b$ and $f\le B_\epsilon$ on $[t^*,t^*+\eta]$, whence
$t^*+\eta\in T$, contradicting $t^*=\sup T$.

{\em Case 1: $f(t^*)<B_\epsilon(t^*)$.} By continuity $f<B_\epsilon$
on $[t^*,t^*+\eta]$ for some $\eta>0$.

{\em Case 2: $f(t^*)=B_\epsilon(t^*)$.} Since $0<\epsilon
e^{2K(t^*-a)}\le\epsilon e^{2K(b-a)}\le1$ we have
$f(t^*)=G(t^*)+\epsilon e^{2K(t^*-a)}\in[-(C+1),C+1]$, so both
$(t^*,f(t^*))$ and $(t^*,G(t^*))$ lie in $S$ and the Lipschitz bound
gives
$$\psi(t^*,f(t^*))\le\psi(t^*,G(t^*))+K\epsilon e^{2K(t^*-a)}
<\psi(t^*,G(t^*))+2K\epsilon e^{2K(t^*-a)}=B_\epsilon'(t^*).$$
Choose $r$ with $\psi(t^*,f(t^*))<r<B_\epsilon'(t^*)$. By hypothesis
the forward difference quotient of $f$ at $t^*\in[a,b)$ is at most
$\psi(t^*,f(t^*))$, so by
Definition~\ref{def:forward-difference-quotient} there is $\eta>0$,
with $t^*+\eta\le b$, such that
$$\frac{f(t^*+\triangle t)-f(t^*)}{\triangle t}<r\qquad\text{for all
}\triangle t\in(0,\eta];$$
after shrinking $\eta$, the right differentiability of $B_\epsilon$
at $t^*$ gives also
$$\frac{B_\epsilon(t^*+\triangle t)-B_\epsilon(t^*)}{\triangle
t}>r\qquad\text{for all }\triangle t\in(0,\eta].$$
Subtracting the two displays and using $f(t^*)=B_\epsilon(t^*)$ we
conclude $f<B_\epsilon$ on $(t^*,t^*+\eta]$.

In both cases the maximality of $t^*$ is contradicted, so $t^*=b$
and the claim is proved.

Finally, fix $t\in[a,b]$. For every sufficiently small $\epsilon>0$
the claim gives
$$f(t)\le G(t)+\epsilon e^{2K(t-a)}\le G(t)+\epsilon e^{2K(b-a)},$$
and letting $\epsilon\rightarrow0^+$ yields $f(t)\le G(t)$.
\end{proof}


We first establish the {\em product case} of the maximum property, in
which the manifold is a product $X\times\Ar$ and the vector field is
$\partial/\partial t$; this case contains all of the analytic content
and is the form used for Ricci flow on compact manifolds.

\begin{lemma}[The maximum of a compact family is continuous]
\label{lem:max-function-continuous-product}
\lean{MorganTianLib.continuous_iSup_of_compact}
\leanok
Let $X$ be a compact topological space and let $F\colon X\times\Ar\to
\Ar$ be continuous. Then
$$F_{\mathrm{max}}(t)=\sup_{x\in X}F(x,t)$$
(a maximum when $X\not=\emptyset$) is a continuous function of $t$.
\end{lemma}

\begin{proof}
\leanok
We may assume $X\not=\emptyset$. Fix $t_0$ and $\epsilon>0$. Since a
continuous function on the nonempty compact set $X\times\{t_0\}$
attains its maximum, there is $x_0\in X$ with
$F(x_0,t_0)=F_{\mathrm{max}}(t_0)$; by continuity of $t\mapsto
F(x_0,t)$, for $t$ close to $t_0$ we get
$F_{\mathrm{max}}(t)\ge F(x_0,t)>F_{\mathrm{max}}(t_0)-\epsilon$.
Conversely, for each $x\in X$ continuity of $F$ at $(x,t_0)$ provides
an open neighborhood $U_x\ni x$ and $\delta_x>0$ such that
$F(y,t)<F(x,t_0)+\epsilon\le F_{\mathrm{max}}(t_0)+\epsilon$ for all
$y\in U_x$ and $|t-t_0|<\delta_x$. Finitely many $U_{x_1},\dots,
U_{x_m}$ cover the compact $X$; for $|t-t_0|<
\min_i\delta_{x_i}$ every $y\in X$ lies in some $U_{x_i}$, so
$F_{\mathrm{max}}(t)\le F_{\mathrm{max}}(t_0)+\epsilon$. Hence
$|F_{\mathrm{max}}(t)-F_{\mathrm{max}}(t_0)|\le\epsilon$ for $t$
close to $t_0$.
\end{proof}

\begin{lemma}[Forward difference quotient of the maximum of a compact family]
\label{lem:forward-difference-maximum-product}
\lean{MorganTianLib.forwardDiffQuotientLE_iSup}
\leanok
\uses{def:forward-difference-quotient,lem:max-function-continuous-product}
Let $X$ be a nonempty compact topological space, let $t_0<b$, and let
$F\colon X\times\Ar\to\Ar$ be continuous. Suppose that $F'\colon
X\times\Ar\to\Ar$ is continuous and that for each $x\in X$ the
function $s\mapsto F(x,s)$ is differentiable on $(t_0,b)$ with
derivative $F'(x,s)$. Let $c\in\Ar$ be such that $F'(x,t_0)\le c$ for
every $x$ with $F(x,t_0)=F_{\mathrm{max}}(t_0)$. Then the forward
difference quotient of $F_{\mathrm{max}}$ at $t_0$ is at most $c$.
\end{lemma}

\begin{proof}
\leanok
\uses{lem:max-function-continuous-product}
Fix $r>c$; we must show that
$(F_{\mathrm{max}}(z)-F_{\mathrm{max}}(t_0))/(z-t_0)<r$ for all
$z>t_0$ sufficiently close to $t_0$. Suppose not. Then there are
$z$ arbitrarily close to $t_0$ in $(t_0,b)$ with
$$\frac{F_{\mathrm{max}}(z)-F_{\mathrm{max}}(t_0)}{z-t_0}\ge r.$$
For each such $z$ choose $x_z\in X$ with
$F(x_z,z)=F_{\mathrm{max}}(z)$ (the maximum is attained). Since
$F(x_z,t_0)\le F_{\mathrm{max}}(t_0)$,
$$\frac{F(x_z,z)-F(x_z,t_0)}{z-t_0}\ge
\frac{F_{\mathrm{max}}(z)-F_{\mathrm{max}}(t_0)}{z-t_0}\ge r,$$
so by the mean value theorem applied to $s\mapsto F(x_z,s)$ on
$[t_0,z]$ there is $s_z\in(t_0,z)$ with $F'(x_z,s_z)\ge r$.

Now let $z$ run over a net (or, if $X$ is metrizable, a sequence)
tending to $t_0^+$ along which the above holds. The points
$(x_z,s_z,z)$ lie in the compact set $X\times[t_0,b]\times[t_0,b]$,
so they have a cluster point $(x_0,s_0,z_0)$. Since $z\rightarrow
t_0$ we get $z_0=t_0$, and since $t_0<s_z<z$ also $s_0=t_0$. Since
$F'$ is continuous and $F'(x_z,s_z)\ge r$ along the net,
$F'(x_0,t_0)\ge r$. Since $F$ and $F_{\mathrm{max}}$ are continuous
(Lemma~\ref{lem:max-function-continuous-product}) and
$F(x_z,z)=F_{\mathrm{max}}(z)$ along the net,
$F(x_0,t_0)=F_{\mathrm{max}}(t_0)$; that is, $x_0$ is a maximizer at
time $t_0$. By hypothesis $F'(x_0,t_0)\le c<r$, a contradiction.
\end{proof}

\begin{corollary}[Forward difference maximum property, product case]
\label{cor:forward-difference-maximum-product}
\lean{MorganTianLib.iSup_le_of_forwardDiff_max}
\leanok
\uses{def:forward-difference-quotient}
Let $X$ be a nonempty compact topological space and let $F,F'\colon
X\times\Ar\to\Ar$ be continuous with $s\mapsto F(x,s)$ differentiable
on $(a,b)$ with derivative $F'(x,s)$, for each $x\in X$. Let $\psi$
be a $C^1$-function on $[a,b]\times\Ar$ and suppose that for each
$t\in[a,b)$ and each $x\in X$ with $F(x,t)=F_{\mathrm{max}}(t)$ we
have $F'(x,t)\le\psi(t,F_{\mathrm{max}}(t))$. Suppose that $G(t)$
satisfies the differential equation $G'(t)=\psi(t,G(t))$ on $[a,b]$
and has initial condition $F_{\mathrm{max}}(a)\le G(a)$. Then for all
$t\in [a,b]$ we have $F_{\mathrm{max}}(t)\le G(t)$.
\end{corollary}

\begin{proof}
\leanok
\uses{lem:max-function-continuous-product,lem:forward-difference-maximum-product,lem:forward-difference-quotient}
$F_{\mathrm{max}}$ is continuous by
Lemma~\ref{lem:max-function-continuous-product}, and for each
$t\in[a,b)$ its forward difference quotient at $t$ is at most
$\psi(t,F_{\mathrm{max}}(t))$ by
Lemma~\ref{lem:forward-difference-maximum-product} (applied with
$c=\psi(t,F_{\mathrm{max}}(t))$, whose hypothesis is the displayed
condition at maximizers). Now apply the comparison
Lemma~\ref{lem:forward-difference-quotient} to $f=F_{\mathrm{max}}$.
\end{proof}

The general statement replaces the product structure by the flow of a
vector field. We record the standard flow-box statement it needs;
its proof is classical ODE theory (existence, uniqueness and smooth
dependence for flows of smooth vector fields), which we do not
reproduce here yet.

\begin{lemma}[Flow of a vector field near a compact set]
\label{lem:vector-field-flow-near-compact}
\lean{MorganTianLib.exists_forall_hasDerivWithinAt_continuousOn_of_contDiffAt,
  MorganTianLib.exists_localFlowAt,
  MorganTianLib.exists_localFlow_of_isCompact,
  MorganTianLib.exists_isLocalFlow_of_isCompact}
\leanok
Let $M$ be a smooth manifold (Hausdorff and without boundary), $\chi$ a
smooth vector field on $M$, and $Z\subset M$ a compact subset. Then there
are $\eta>0$, an open neighborhood $U\supset Z$, and a jointly continuous
map $\Phi\colon U\times(-\eta,\eta)\to M$, $\Phi(x,s)=\varphi_s(x)$, such
that $\varphi_0=\mathrm{id}_U$ and, for each $x\in U$, the curve
$s\mapsto\varphi_s(x)$ is an integral curve of $\chi$:
$$\frac{d}{ds}\varphi_s(x)=\chi(\varphi_s(x))\qquad(x\in U,\ |s|<\eta).$$
Consequently, for every smooth function $F\colon M\to\Ar$ the composition
$s\mapsto F(\varphi_s(x))$ is differentiable with
$$\frac{d}{ds}F(\varphi_s(x))=\chi(F)\left(\varphi_s(x)\right),$$
and if ${\bf t}\colon M\to\Ar$ is a smooth function with $\chi({\bf t})=1$,
then ${\bf t}(\varphi_s(x))={\bf t}(x)+s$ for all $x\in U$, $|s|<\eta$.
\end{lemma}

\begin{proof}
\leanok
{\em Step 1: a uniform flow box for a $C^1$ vector field on a normed
space.} Let $f$ be a $C^1$ vector field defined near a point $z_0$ of a
Banach space $E$, and let $W$ be a neighborhood of $z_0$. We claim there
are $r>0$, $\epsilon>0$ and a map $Z\colon\overline B(z_0,r)\times
[-\epsilon,\epsilon]\to E$, jointly continuous, with $Z(z,0)=z$,
$\frac{d}{ds}Z(z,s)=f(Z(z,s))$, and $Z(z,s)\in W$ throughout. Choose
$a>0$ so small that $\overline B(z_0,a)\subset W$, $f$ is bounded by some
$L$ and Lipschitz with some constant $K$ on $\overline B(z_0,a)$ (both
exist since $f$ is $C^1$). For $r<a$ and $\epsilon\le(a-r)/L$, the Picard
integral operator $\alpha\mapsto\left(s\mapsto z+\int_0^s
f(\alpha(\tau))\,d\tau\right)$ maps the complete space of continuous
curves $[-\epsilon,\epsilon]\to\overline B(z_0,a)$ starting at
$z\in\overline B(z_0,r)$ to itself, and a sufficiently high iterate is a
contraction (its Lipschitz constant is bounded by $(K\epsilon)^n/n!$);
by the Banach fixed point theorem there is a unique solution
$Z(z,\cdot)$ through each $z\in\overline B(z_0,r)$, staying in
$\overline B(z_0,a)\subset W$, with
$\frac{d}{ds}Z(z,s)=f(Z(z,s))$ by the fundamental theorem of calculus.
Comparing the Picard iterations for two initial points $x,y$ (or applying
Gr\"onwall's inequality to $|Z(x,s)-Z(y,s)|$) gives
$$\left|Z(x,s)-Z(y,s)\right|\le L'\,|x-y|
\qquad\left(x,y\in\overline B(z_0,r),\ |s|\le\epsilon\right)$$
with a constant $L'$ independent of $s$; since moreover each solution is
$L$-Lipschitz in $s$ (its derivative $f(Z(z,s))$ is bounded by $L$), the
map $Z$ is jointly continuous: continuity in $s$ for fixed $z$, plus a
uniform-in-$s$ Lipschitz bound in $z$, give continuity on the product.

{\em Step 2: a flow box in one chart of the manifold.} Fix $z\in M$ and
a chart $\psi$ at $z$ with image in the model space $E$. The coordinate
representative $f=d\psi\circ\chi\circ\psi^{-1}$ of $\chi$ is $C^1$ near
$\psi(z)$, since $\chi$ is smooth. Apply Step 1 at $z_0=\psi(z)$ with
$W$ the (open) image of the chart: this gives $r,\epsilon>0$ and a
jointly continuous chart flow $Z$ confined to the chart image. Let
$V_z=\psi^{-1}\left(B(\psi(z),r)\right)$, an open neighborhood of $z$,
and define
$$\Phi_z\colon V_z\times(-\epsilon,\epsilon)\to M,\qquad
\Phi_z(x,s)=\psi^{-1}\left(Z(\psi(x),s)\right).$$
Then $\Phi_z(x,0)=x$; each curve $s\mapsto\Phi_z(x,s)$ is an integral
curve of $\chi$, because the confinement keeps $Z(\psi(x),s)$ in the
chart image, where $\psi$ is a diffeomorphism identifying solutions of
$\dot e=f(e)$ with integral curves of $\chi$; and $\Phi_z$ is jointly
continuous as the composition $\psi^{-1}\circ Z\circ(\psi\times
\mathrm{id})$ of continuous maps.

{\em Step 3: gluing over the compact set.} The open sets $V_z$, $z\in Z$,
cover the compact $Z$; choose a finite subcover
$V_{z_1},\dots,V_{z_k}$ with times $\epsilon_1,\dots,\epsilon_k$ (if
$Z=\emptyset$ take $U=\emptyset$, $\eta=1$). Set
$\eta=\min_i\epsilon_i>0$ and $U=\bigcup_iV_{z_i}\supset Z$, and for
$x\in U$ define $\Phi(x,s)=\Phi_{z_i}(x,s)$ for any $i$ with $x\in
V_{z_i}$. This is well defined: for $x\in V_{z_i}\cap V_{z_j}$, both
$s\mapsto\Phi_{z_i}(x,s)$ and $s\mapsto\Phi_{z_j}(x,s)$ are integral
curves of $\chi$ on $(-\eta,\eta)$ passing through $x$ at $s=0$, and
integral curves of a $C^1$ vector field on a Hausdorff manifold through
a common point coincide on their common connected domain (in a chart
this is the classical uniqueness theorem for locally Lipschitz
right-hand sides; the agreement set is open by uniqueness and closed by
continuity and the Hausdorff property, hence all of $(-\eta,\eta)$). The
same uniqueness argument shows $\Phi$ agrees with the jointly continuous
$\Phi_{z_i}$ on all of $(U\cap V_{z_i})\times(-\eta,\eta)$, an open set,
so $\Phi$ is jointly continuous, and $\varphi_0=\mathrm{id}_U$.

{\em Step 4: the derived identities.} For smooth $F$ the chain rule
along the integral curve $s\mapsto\varphi_s(x)$ gives
$\frac{d}{ds}F(\varphi_s(x))=dF\left(\chi(\varphi_s(x))\right)
=\chi(F)(\varphi_s(x))$. If $\chi({\bf t})=1$ then the function
$s\mapsto{\bf t}(\varphi_s(x))-({\bf t}(x)+s)$ has everywhere vanishing
derivative on the interval $(-\eta,\eta)$ and vanishes at $s=0$, so it
vanishes identically by the mean value inequality; that is,
${\bf t}(\varphi_s(x))={\bf t}(x)+s$.
\end{proof}

The flow $\Phi$ is in fact jointly \emph{smooth} — this is the classical
smooth dependence of solutions of ODEs on their initial conditions — but
joint continuity is all the maximum principle below consumes, and all we
formalize here.

Given the local flow, the maximum principle itself is a purely
topological statement, which we isolate; it makes no smoothness
demands beyond differentiability of $F$ along the flow lines.

\begin{lemma}[Forward difference maximum property along a local flow]
\label{lem:forward-difference-maximum-flow}
\lean{MorganTianLib.IsLocalFlow,
  MorganTianLib.continuousOn_levelMax,
  MorganTianLib.forwardDiffQuotientLE_levelMax,
  MorganTianLib.levelMax_le_of_isLocalFlow}
\leanok
\uses{def:forward-difference-quotient}
Let $M$ be a topological space and let
${\bf t},F,\chi F\colon M\to\Ar$ be functions with ${\bf t}$ and
$\chi F$ continuous. Let
$${\mathcal Z}=\left\{x\in M\ :\ F(y)\le F(x)\ \text{for all}\ y\
\text{with}\ {\bf t}(y)={\bf t}(x)\right\}$$
be the set of fibrewise maximizers, and let $a\le b$. Assume:
\begin{enumerate}
\item ${\mathcal Z}$ is compact, and for every $t\in[a,b]$ there is
$x\in{\mathcal Z}$ with ${\bf t}(x)=t$;
\item (\emph{local flow at unit ${\bf t}$-speed}) there are an open set
$U\supset{\mathcal Z}$, a number $\eta>0$, and a jointly continuous map
$\Phi\colon U\times(-\eta,\eta)\to M$ such that for all $x\in U$ and
$|s|<\eta$
$$\Phi(x,0)=x,\qquad {\bf t}(\Phi(x,s))={\bf t}(x)+s,\qquad
\frac{d}{ds}F(\Phi(x,s))=\chi F(\Phi(x,s));$$
\item $\chi F(x)\le\psi({\bf t}(x),F(x))$ for every $x\in{\mathcal Z}$,
where $\psi$ is a $C^1$-function on $[a,b]\times\Ar$.
\end{enumerate}
Set $F_{\mathrm{max}}(t)=\max_{{\bf t}(x)=t}F(x)$ for $t\in[a,b]$; the
maximum is attained at any point of ${\mathcal Z}$ in the fibre over
$t$. Then $F_{\mathrm{max}}$ is continuous on $[a,b]$, and its forward
difference quotient at each $t\in[a,b)$ is at most
$\psi(t,F_{\mathrm{max}}(t))$. Consequently, if $G$ is continuous on
$[a,b]$ and satisfies $G'(t)=\psi(t,G(t))$ on $[a,b)$ with
$F_{\mathrm{max}}(a)\le G(a)$, then $F_{\mathrm{max}}(t)\le G(t)$ for
all $t\in[a,b]$.
\end{lemma}

\begin{proof}
\leanok
\uses{lem:forward-difference-quotient}
Write $\eta'=\eta/2$. The set
$K=\Phi\left({\mathcal Z}\times[-\eta',\eta']\right)$ is compact, being
the image of a compact set under the jointly continuous map $\Phi$; fix
a bound $C$ for $|\chi F|$ on $K$. For $x\in{\mathcal Z}$ the function
$s\mapsto F(\Phi(x,s))$ has derivative $\chi F(\Phi(x,s))$ on
$[-\eta',\eta']$, of absolute value at most $C$ there; the mean value
inequality gives
$$\left|F(\Phi(x,s_2))-F(\Phi(x,s_1))\right|\le C\,|s_2-s_1|
\qquad\left(x\in{\mathcal Z},\ s_1,s_2\in[-\eta',\eta']\right).
\qquad(*)$$
Note also that if $x\in{\mathcal Z}$ then
$F_{\mathrm{max}}({\bf t}(x))=F(x)$, while every $y$ with
${\bf t}(y)={\bf t}(x)$ satisfies $F(y)\le F_{\mathrm{max}}({\bf t}(x))$.

{\em Step 1: $F_{\mathrm{max}}$ is continuous on $[a,b]$.} Let $t_1\le
t_2$ lie in $[a,b]$ with $\triangle t=t_2-t_1\le\eta'$, and choose
$x_1,x_2\in{\mathcal Z}$ with ${\bf t}(x_i)=t_i$ (hypothesis (1)). The
point $\Phi(x_1,\triangle t)$ lies in the fibre over $t_2$, so by $(*)$
$$F_{\mathrm{max}}(t_2)\ge F\left(\Phi(x_1,\triangle t)\right)\ge
F(x_1)-C\triangle t=F_{\mathrm{max}}(t_1)-C\triangle t;$$
symmetrically $\Phi(x_2,-\triangle t)$ lies in the fibre over $t_1$, so
$F_{\mathrm{max}}(t_1)\ge F_{\mathrm{max}}(t_2)-C\triangle t$. Hence
$|F_{\mathrm{max}}(t_2)-F_{\mathrm{max}}(t_1)|\le C\,|t_2-t_1|$ for
nearby times in $[a,b]$, and $F_{\mathrm{max}}$ is locally Lipschitz,
hence continuous, on $[a,b]$.

{\em Step 2: the forward difference quotient bound.} Fix $t_0\in[a,b)$
and set $c=\psi(t_0,F_{\mathrm{max}}(t_0))$. Every $x\in{\mathcal Z}$
with ${\bf t}(x)=t_0$ has $F(x)=F_{\mathrm{max}}(t_0)$ and hence, by
hypothesis (3), $\chi F(x)\le\psi(t_0,F_{\mathrm{max}}(t_0))=c$. Fix
$r>c$ and suppose for contradiction that there are $z>t_0$ arbitrarily
close to $t_0$ with
$$\frac{F_{\mathrm{max}}(z)-F_{\mathrm{max}}(t_0)}{z-t_0}\ge r;$$
we may restrict attention to $z<\min(b,t_0+\eta')$. For each such $z$
choose $y_z\in{\mathcal Z}$ with ${\bf t}(y_z)=z$ (hypothesis (1)), so
that $F(\Phi(y_z,0))=F(y_z)=F_{\mathrm{max}}(z)$, while
$\Phi(y_z,-(z-t_0))$ lies in the fibre over $t_0$ and hence
$F\left(\Phi(y_z,-(z-t_0))\right)\le F_{\mathrm{max}}(t_0)$. The
difference quotient of $s\mapsto F(\Phi(y_z,s))$ over $[-(z-t_0),0]$ is
therefore at least
$\left(F_{\mathrm{max}}(z)-F_{\mathrm{max}}(t_0)\right)/(z-t_0)\ge r$,
and the mean value theorem yields $\tau_z\in(-(z-t_0),0)$ with
$$\chi F\left(\Phi(y_z,\tau_z)\right)\ge r.$$
The pairs $(y_z,\tau_z)$ lie in the compact set
${\mathcal Z}\times[-\eta',0]$, so as $z\to t_0^+$ the triples
$(y_z,\tau_z,z)$ cluster at some $(y_0,\tau_0,t_0)$ with
$y_0\in{\mathcal Z}$. From $-(z-t_0)<\tau_z<0$ and $z\to t_0$ we get
$\tau_0=0$; from ${\bf t}(y_z)=z$ and continuity of ${\bf t}$ we get
${\bf t}(y_0)=t_0$. By joint continuity of $\Phi$ the points
$\Phi(y_z,\tau_z)$ cluster at $\Phi(y_0,0)=y_0$, so continuity of
$\chi F$ gives $\chi F(y_0)\ge r$. But $y_0\in{\mathcal Z}$ with
${\bf t}(y_0)=t_0$, so $\chi F(y_0)\le c<r$, a contradiction. Hence for
every $r>c$ the right difference quotients of $F_{\mathrm{max}}$ at
$t_0$ are eventually $<r$; that is, the forward difference quotient of
$F_{\mathrm{max}}$ at $t_0$ is at most $c=\psi(t_0,F_{\mathrm{max}}(t_0))$.

{\em Step 3: conclusion.} $F_{\mathrm{max}}$ is continuous on $[a,b]$
(Step 1) with forward difference quotient at most
$\psi(t,F_{\mathrm{max}}(t))$ at every $t\in[a,b)$ (Step 2), so the
comparison Lemma~\ref{lem:forward-difference-quotient} yields
$F_{\mathrm{max}}(t)\le G(t)$ on $[a,b]$.
\end{proof}


The application we shall make of these results is the following.

\begin{proposition}[Forward difference maximum property]
\label{prop:forward-difference-maximum}
\label{fordiffmax}
\lean{MorganTianLib.levelMax_le_of_dir_le}
\leanok
\uses{def:forward-difference-quotient}
\source{morgan-tian}
Let $M$ be a smooth manifold with a smooth vector field $\chi$ and a
smooth function ${\bf t}\colon M\to [a,b]$ with $\chi({\bf t})=1$.
Suppose also that $F\colon M\to \Ar$ is a smooth function with the
properties:
\begin{enumerate}
\item for each $t_0\in [a,b]$ the restriction of $F$ to the level set ${\bf
t}^{-1}(t_0)$ achieves its maximum, and \item the subset ${\mathcal Z}$ of $M$
consisting of all $x$ for which $F(x)\ge F(y)$ for all $y\in {\bf t}^{-1}({\bf
t}(x))$ is a compact set.
\end{enumerate}
Suppose also that at each $x\in {\mathcal Z}$ we have $\chi(F(x))\le \psi({\bf
t}(x),F(x))$, where $\psi$ is a $C^1$-function on $[a,b]\times\Ar$. Set $F_{\mathrm{max}}(t)=\max_{x\in {\bf
t}^{-1}(t)}F(x)$. Then $F_{\mathrm{max}}(t)$ is a continuous function and
$$\frac{dF_{\mathrm{max}}}{dt}(t)\le \psi(t,F_{\mathrm{max}}(t))$$
in the sense of forward difference quotients.
Suppose that $G(t)$ satisfies the
differential equation
$$G'(t)=\psi(t,G(t))$$ and has initial condition $F_{\mathrm{max}}(a)\le G(a)$. Then for all $t\in [a,b]$ we have
$$F_{\mathrm{max}}(t)\le G(t).$$
\end{proposition}

\begin{proof}
\leanok
\uses{lem:vector-field-flow-near-compact,lem:forward-difference-maximum-flow}
Apply Lemma~\ref{lem:vector-field-flow-near-compact} to the compact set
${\mathcal Z}$ of fibrewise maximizers: it yields $\eta>0$, an open
$U\supset{\mathcal Z}$, and a jointly continuous local flow
$\Phi\colon U\times(-\eta,\eta)\to M$ of $\chi$ with $\Phi(x,0)=x$ and
${\bf t}(\Phi(x,s))={\bf t}(x)+s$, along which
$$\frac{d}{ds}F(\Phi(x,s))=dF\left(\chi(\Phi(x,s))\right)
=\chi(F)\left(\Phi(x,s)\right),$$
and $\chi(F)$ and ${\bf t}$ are continuous since $F$, ${\bf t}$ and
$\chi$ are smooth. Hypothesis (1) says that every fibre
${\bf t}^{-1}(t)$, $t\in[a,b]$, contains a fibrewise maximizer, i.e.\ a
point of ${\mathcal Z}$, and hypothesis (2) says that ${\mathcal Z}$ is
compact. All hypotheses of
Lemma~\ref{lem:forward-difference-maximum-flow} are therefore satisfied
(with $\chi F=\chi(F)$), and its conclusions are precisely the
continuity of $F_{\mathrm{max}}$, the forward difference quotient bound
$dF_{\mathrm{max}}/dt\le\psi(t,F_{\mathrm{max}}(t))$, and the comparison
$F_{\mathrm{max}}(t)\le G(t)$ on $[a,b]$.
\end{proof}

